\documentclass[12pt,a4paper]{article}

% Shared HSC style packages (see ../HSC-Common/).
\makeatletter
\def\input@path{{../HSC-Common/styles/}{./}}
\makeatother
\newcommand{\HSCBookletHeaderTitle}{Inequality Functional Equations}
\usepackage{amsmath}
\usepackage{amssymb}
\usepackage{amsthm}
\usepackage{geometry}
\usepackage{graphicx}
\usepackage{enumitem}
\usepackage{tikz}
\usepackage[hidelinks]{hyperref}
\hypersetup{hidelinks}

\usepackage{dl101-colors}
\usepackage{dl101-hints}
\usepackage{dl101-boxes}
\usepackage{dl101-theorems}
\usepackage{dl101-hsc-problems}
\usepackage{dl101-layout}

% Compatibility wrapper for the booklet's existing inline hint syntax.
\newcommand{\upsidehint}[1]{%
  \begin{hint}[#1]\end{hint}%
}

\usetikzlibrary{calc}

\title{\textbf{Inequality Functional Equations}\\
\large An Olympiad Problem Collection}
\author{Vu Hung Nguyen}
\date{2026}

\renewcommand{\contentsname}{Table of Contents}

\begin{document}

\thispagestyle{empty}%
\begin{tikzpicture}[remember picture, overlay]
  \node[inner sep=0pt] at (current page.center) {
    \includegraphics[width=\paperwidth,height=\paperheight]{assets/spiral-matte-bg.png}%
  };
  \fill[black, opacity=0.4] (current page.south west) rectangle (current page.north east);
  \node[anchor=center, text=white, align=center, font=\sffamily] at (current page.center) {
    {\Huge \bfseries IMO-Style}\\[1.5em]
    {\LARGE Inequality Functional Equations}\\[3em]
    {\Large Vu Hung Nguyen}%
  };
\end{tikzpicture}

\newpage

\setcounter{page}{1}
% Shared QR code page content (page 2): GitHub repo + Vu’s Maths Hub, side by side.
\vspace*{\fill}
\begin{center}
\begin{minipage}[c]{0.4\textwidth}
  \centering
  \includegraphics[width=0.9\linewidth]{assets/qrcode.png}\\[0.4cm]
  {\small\itshape Download PDF and resources:}\\[0.15cm]
  {\footnotesize\url{https://github.com/vuhung16au/math-olympiad-ml}}
\end{minipage}%
\hfill
\begin{minipage}[c]{0.4\textwidth}
  \centering
  \includegraphics[width=0.9\linewidth]{assets/qrcode-vumaths.png}\\[0.4cm]
  {\small\itshape Read online at Vu's Maths Hub:}\\[0.15cm]
  {\footnotesize\url{https://vumaths.com/}}
\end{minipage}
\end{center}
\vspace*{\fill}

\begingroup\small\noindent This work is licensed under CC BY-NC 4.0 (Attribution-Non-Commercial), see the LICENSE file on the github page for more info.\par\endgroup

\newpage
\thispagestyle{empty}
\vspace*{\fill}
\begin{center}
\itshape
``Exercise is the beste instrument in learnyng.''

\vspace{1em}
\normalfont Robert Recorde
\end{center}
\vspace*{\fill}

\newpage
\tableofcontents
\newpage

\section*{Preface}
\addcontentsline{toc}{section}{Preface}

This is a growing booklet of olympiad-style functional equations and
inequalities. It is written especially for students in Years 10--12 preparing
for mathematical olympiads, as well as for anyone who enjoys practising with
challenging, unusual, and unfamiliar problems. The aim is to develop flexible
problem-solving habits, not just to collect answers.

The collection includes problems with different domains, methods, and levels
of rigidity. Some problems have more than one solution, while others admit a
parameter family or an exceptional value until the final hypotheses are used.

\section*{Introduction}
\addcontentsline{toc}{section}{Introduction}

The problems focus on olympiad-style functional equations built around
inequalities. Their domains include the positive real numbers, the real
numbers, the positive rationals, and the positive integers. These domains
require different tools: order and limits on the positive reals, symmetry and
regularity on the real line, prime exponents on the positive rationals, and
counting or divisibility on the integers. The booklet is arranged in five
themed parts: additive and multiplicative rigidity; reciprocal and harmonic
transforms; iteration, fixed points, and growth; mean and norm squeezes; and
real-line and polynomial inequalities. Problems come first, grouped by theme,
and the full solutions appear together at the end.

For each problem, first read the statement carefully and try to solve it on
your own. Pay attention to the domain and to what the inequalities say for
special choices of inputs. If you get stuck, use the hint as a nudge rather
than as a replacement for a proof. Once you have a solution, compare it with

the one at the end: check each step, note any idea that you did not use, and
compare alternative solutions when they are provided.

\section*{Notations}
\addcontentsline{toc}{section}{Notations}

Only the booklet-specific conventions are listed here; standard olympiad
notation is left implicit.

\paragraph{Domains.}
$\mathbb R$, $\mathbb Q$, and $\mathbb Z$ denote the real, rational, and
integer numbers. A subscript $>0$ means positive elements, while a subscript
$\ge0$ means nonnegative elements; for example,
$\mathbb Q_{>0}$ and $\mathbb Z_{\ge0}$. A function such as
$f:\mathbb Q\to[0,\infty)$ has exactly the displayed domain and codomain.

\paragraph{Iterates and auxiliary functions.}
$f^{\circ n}$ denotes the $n$-fold iterate of $f$, not the pointwise power
$f^n$. When a candidate model $M$ is known, we may write
$g(x)=f(x)-M(x)$ for its defect, or $f(x)/M(x)$ for its normalized ratio
when $M(x)>0$.

\paragraph{Discrete approximations.}
$\lfloor t\rfloor$ and $\lceil t\rceil$ are the floor and ceiling of $t$.
The term \emph{dyadic} means having denominator a power of $2$ (for example,
$m/2^n$); dyadic numbers are used as a dense subset of the real line.

\paragraph{Prime exponents.}
For a prime $p$, $v_p(x)$ denotes the exponent of $p$ in a nonzero rational
$x$, allowing negative exponents. Thus
\[
v_p(xy)=v_p(x)+v_p(y).
\]

\section*{Fundamental Review: Facts and Tools}
\addcontentsline{toc}{section}{Fundamental Review: Facts and Tools}

The \emph{squeeze principle} says that if a quantity is bounded above and
below by the same value, it must equal that value. In these problems, one
inequality often supplies an upper bound and another a matching lower bound,
sometimes only after a substitution or iteration.

\paragraph{Order and monotonicity.}
Adding inequalities preserves their direction. Multiplying or dividing by a
positive quantity also preserves the direction, while doing so by a negative
quantity reverses it. Applying an increasing function preserves order; for
positive numbers, taking reciprocals reverses order. A function is
\emph{nondecreasing} if $x\le y$ implies $f(x)\le f(y)$. For example, a
positive superadditive function on the positive reals is strictly increasing:
if $0<x<y$, then $f(y)\ge f(x)+f(y-x)>f(x)$.

\paragraph{Means.}
For nonnegative $u,v$, the arithmetic--geometric mean inequality is
\[
\frac{u+v}{2}\ge\sqrt{uv},
\]
with equality exactly when $u=v$. The root-mean-square--arithmetic-mean
inequality is
\[
\sqrt{\frac{u^2+v^2}{2}}\ge\frac{u+v}{2}
\qquad (u,v\ge0),
\]
again with equality exactly when $u=v$. These inequalities are useful when
the problem places a function expression between two familiar means.

\paragraph{Iteration and approximation.}
Mathematical induction justifies repeated addition or multiplication, such as
deriving bounds at integer multiples or powers. The Archimedean property says
that the positive integers are unbounded. In particular, for $t>0$,
\[
0\le t-\frac{\lfloor nt\rfloor}{n}<\frac1n,
\]
so $\lfloor nt\rfloor/n\to t$. Floor approximations of this kind let us
compare a value at an arbitrary real input with values at integer multiples.

\paragraph{Density and monotonicity.}
The rational numbers are dense in the real numbers, and
$\{2^q:q\in\mathbb Q\}$ is dense in the positive reals. If a nondecreasing
function agrees with the identity on a dense subset of its domain, it agrees
with the identity throughout: nearby points from the dense subset bound its
value from below and above.

\paragraph{Defects and normalization.}
When a candidate model $M(x)$ is suggested by the problem, it is often useful
to study the defect
\[
g(x)=f(x)-M(x)
\]
or the ratio $f(x)/M(x)$ when $M(x)>0$. The hypotheses may become simpler in
terms of $g$: a polynomial expression may become a midpoint inequality, while
a multiplicative expression may become submultiplicative after taking
reciprocals or logarithms. Always check that the normalization is defined on
the stated domain.

\paragraph{Iteration and bounded growth.}
If a relation is applied repeatedly, write the first few cases and identify
the resulting sequence. A nonnegative quantity cannot have a negative partial
sum, a bounded quantity cannot grow exponentially, and a quantity bounded by
$C/2^n$ for every $n$ must be zero. These elementary observations often turn
an inequality into an identity.

\paragraph{Discrete rigidity.}
On the positive integers, injectivity permits counting arguments: $N$
distinct inputs require $N$ distinct positive outputs. On the positive
rationals, unique factorization permits comparison of prime exponents and
divisibility chains. These tools are different from real-variable arguments,
where order, density, and regularity are central.

\paragraph{Midpoints and quadratic identities.}
The parallelogram identity
\[
\frac{(x+y)^2+(x-y)^2}{2}=x^2+y^2
\]
often simplifies inequalities on the real line. A related midpoint fact is
that a monotone function on an interval satisfying
\[
f\!\left(\frac{x+y}{2}\right)=\frac{f(x)+f(y)}2
\]
for all $x,y$ in that interval is affine. This is a regularity result for the midpoint Jensen
equation; no continuity assumption needs to be added when monotonicity is
known.

Basic algebra is used throughout: factor differences of squares, use the laws
of exponents, and handle reciprocals carefully, noting that reciprocal order
reversal requires positive quantities. Also choose inputs such as $x=y$,
$y=1$, or $y=-x$ only when the domain permits them. Normalizing a function,
for example by setting $g(x)=f(x)/x^k$, can expose simpler behavior; always
check that the normalization is defined on the stated domain. When a formula
is first established on the integers, rationals, or dyadic rationals, do not
assume that it extends automatically to all real inputs: first prove the
continuity, monotonicity, boundedness, or other regularity needed for that
extension.

\section*{Fundamental Review: Core Theorems}
\addcontentsline{toc}{section}{Fundamental Review: Core Theorems}

The following facts appear repeatedly in olympiad-style functional
inequalities. Use them only when their hypotheses have been verified.

\paragraph{1. Squeezing and equality.}
If $A\le B\le A$, then $A=B$. More generally, if

\[
0\le u_n\le v_n\qquad\text{and}\qquad v_n\to0,
\]
then $u_n\to0$. One inequality often gives one side of a bound, while a
reciprocal substitution, symmetry, or iteration gives the other side. Always
record the equality condition in a standard inequality; it often identifies
the candidate function.

\paragraph{2. The basic mean inequalities.}
For $u,v\ge0$,
\[
\sqrt{uv}\le \frac{u+v}{2}
\]
with equality if and only if $u=v$, and
\[
\frac{u+v}{2}\le
\sqrt{\frac{u^2+v^2}{2}}
\]
with equality if and only if $u=v$. The Cauchy--Schwarz inequality is
\[
\left(\sum_{i=1}^n a_i b_i\right)^2
\le
\left(\sum_{i=1}^n a_i^2\right)
\left(\sum_{i=1}^n b_i^2\right),
\]
including the useful special case $2uv\le u^2+v^2$.

\paragraph{3. Triangle and norm inequalities.}
For real numbers,
\[
|x+y|\le |x|+|y|.
\]
Consequently, for $0\le\alpha\le1$,
\[
|x+y|^\alpha\le |x|^\alpha+|y|^\alpha.
\]
The Euclidean norm satisfies the triangle inequality as well:
\[
\sqrt{\sum_i(x_i+y_i)^2}
\le \sqrt{\sum_i x_i^2}+\sqrt{\sum_i y_i^2}.
\]

\paragraph{4. Induction, iteration, and growth.}
Repeated application is justified by induction. For example,
\[
g(x+y)\ge g(x)+g(y)\Longrightarrow g(nx)\ge ng(x),
\]
and
\[
g(xy)\le g(x)g(y)\Longrightarrow g(x^n)\le g(x)^n.
\]
A bounded quantity cannot grow exponentially. If a nonnegative quantity is
bounded above by $C/r^n$ for every $n$, where $r>1$, it is zero.

\paragraph{5. Archimedean approximation.}
The positive integers are unbounded, and for $t>0$,
\[
0\le t-\frac{\lfloor nt\rfloor}{n}<\frac1n,
\qquad\text{so}\qquad
\frac{\lfloor nt\rfloor}{n}\longrightarrow t.
\]
Floor and ceiling approximations transfer conclusions from integer inputs to
arbitrary real inputs, usually after monotonicity has been proved.

\paragraph{6. Monotonicity and density.}
If $x\le y$ implies $f(x)\le f(y)$, then $f$ is nondecreasing. A positive
superadditive function on $(0,\infty)$ is strictly increasing, since
\[
0<x<y\Longrightarrow f(y)\ge f(x)+f(y-x)>f(x).
\]
The rationals are dense in $\mathbb R$, and
$\{2^q:q\in\mathbb Q\}$ is dense in $(0,\infty)$. A nondecreasing function
that agrees with a known function on a dense subset agrees with it everywhere.

\paragraph{7. Midpoint and additive rigidity.}
If a monotone function satisfies
\[
f\!\left(\frac{x+y}{2}\right)=\frac{f(x)+f(y)}2,
\]
then it is affine. An additive function satisfying
$f(x+y)=f(x)+f(y)$ is linear once one proves a regularity condition such as
monotonicity, boundedness on an interval, or continuity at one point.
Additivity alone does not justify continuity.

\textbf{Continuity is proved, not assumed, in this booklet}, especially in
the additive, quadratic, and midpoint-identity problems.

\paragraph{8. Quadratic identities.}
The parallelogram identity is
\[
f(x+y)+f(x-y)=2f(x)+2f(y).
\]
For $f(x)=x^2$ this follows from
\[
(x+y)^2+(x-y)^2=2x^2+2y^2.
\]
A one-sided inequality can become this identity after symmetry supplies the
reverse inequality. A separately proved regularity condition is still needed
to rule out pathological solutions on the real line.

\paragraph{9. Integer and rational rigidity.}
Unique prime factorization gives
\[
n=\prod_p p^{v_p(n)}
\]
for positive integers, with integer exponents; for nonzero positive rationals,
the exponents may be negative. Multiplicative problems can therefore be
studied prime by prime. Bézout's identity says that if $\gcd(a,b)=1$, then
there exist integers $r,s$ such that $ra+sb=1$.

On the positive integers, injectivity gives a counting principle: $N$
distinct inputs require $N$ distinct positive outputs. This can contradict an
upper bound placing all outputs in a set with fewer than $N$ elements.

\section*{Frequently Seen Patterns}
\addcontentsline{toc}{section}{Frequently Seen Patterns}

Look for special inputs that produce bounds, identities, or fixed points, then
combine or repeat those consequences until they become decisive. The useful
substitutions depend on the domain: try $x=y$, $y=1$, $y=-x$, $x=0$, $x=1$,
reciprocal inputs, scaled inputs, and values such as $x=f(y)$ or $x=f(x)$
when these are permitted.

\paragraph{Normalize the expected answer.}
If the statement contains products, reciprocals, roots, powers, or a familiar
polynomial expression, introduce a transformed function. Useful examples are
\[
g(x)=\frac1{f(x)},\qquad g(x)=\sqrt{f(x)},\qquad
g(x)=f(x)-x^2,\qquad g(x)=\frac{f(x)}x,\qquad g(x)=\log f(x).
\]
A good normalization can turn a complicated inequality into superadditivity,
subadditivity, submultiplicativity, midpoint concavity, or a statement about
a nonnegative defect.

\paragraph{Track and eliminate a defect.}
After identifying a likely model $M(x)$, define $g(x)=f(x)-M(x)$ or
$g(x)=f(x)/M(x)$. If a positive defect grows under translation, scaling,
multiplication, or iteration while another condition keeps it bounded, the
defect must vanish. Conversely, if repeated scaling bounds it by a quantity
tending to zero, it must be zero.

\paragraph{Reverse inequalities and propagate equality.}
Interchanging variables, reversing an arithmetic progression, or applying a
condition at reciprocal inputs can produce the opposite inequality. Together
the two directions force equality. A single anchor value can then propagate
through integer multiples, powers, roots, dyadic scaling, or density. Equality
cases in standard mean and norm inequalities are especially informative.

\paragraph{Use iteration and composition.}
A pointwise estimate for $f$ is often most effective after evaluating it at
$f(x)$. Compare that result with a condition on $f(f(x))$ or a higher iterate
to eliminate constants, force fixed points, or rule out an alternative branch.

\paragraph{Use the domain's arithmetic structure.}
On the positive integers, count distinct values and use divisibility. On the
positive rationals, track prime supports and prime exponents under iteration.
On the positive reals, use order, reciprocal reversal, roots, and density. On
the full real line, exploit signs and symmetry, but prove any needed
monotonicity or continuity rather than assuming it.

\paragraph{Prove regularity before extending a formula.}
Many real-line solutions first determine a function on integers, rationals, or
dyadic rationals. That information alone is not enough. Establish continuity,
monotonicity, boundedness, or a Lipschitz estimate from the hypotheses before
passing from a dense subset to all real inputs. In particular, continuity is
proved, not assumed, in additive, quadratic, and midpoint-identity arguments.

\paragraph{Keep equality cases and exceptional parameters.}
After finding a candidate, substitute it into every original condition and
record where equality occurs. Do not discard zero slopes, zero exponents,
constant solutions, endpoint values, or exceptional values at the origin
without checking whether the domain and hypotheses genuinely exclude them.

\newpage
\setcounter{tocdepth}{3}
\tableofcontents
\newpage

\section*{Problems}
\addcontentsline{toc}{section}{Problems}

\subsection*{Part I. Additive--Multiplicative Rigidity}
\addcontentsline{toc}{subsection}{Part I. Additive--Multiplicative Rigidity}

\subsubsection*{Problem 1. The Additive--Multiplicative Rigidity}
\addcontentsline{toc}{subsubsection}{Problem 1. The Additive--Multiplicative Rigidity}

\begin{exercisebox}[hard]

Find all functions $f$ that map positive real numbers to positive real
numbers and satisfy, for all $x,y>0$,
\[
f(x+y)\geq f(x)f(y)\geq f(xy).
\]

\begin{flushright}
\textit{Composed by Vu Hung Nguyen, 2026}
\end{flushright}

\upsidehint{Evaluate at $1$, then compare repeated addition with repeated multiplication.}

\end{exercisebox}
\subsubsection*{Problem 2. IMO 2013 -- Problem 5}
\addcontentsline{toc}{subsubsection}{Problem 2. IMO 2013 -- Problem 5}

\begin{exercisebox}[medium]

Let $f$ be a function that maps positive rational
numbers to real numbers and satisfies, for all
positive rational numbers $x,y$,
\[
\begin{aligned}
f(x)f(y)&\geq f(xy),\\
f(x+y)&\geq f(x)+f(y).
\end{aligned}
\]

Suppose there exists a rational number $a>1$
such that $f(a)=a$.

Prove that
\[
f(x)=x\quad\text{for all positive rational numbers }x.
\]

\medskip
\noindent
\textit{Source: International Mathematical Olympiad
2013 (54th IMO), Day 2, Problem 5.
The wording has been slightly adapted without
changing the mathematical content.}

\upsidehint{Iterate both inequalities at sums and products of $a$; can these bounds trap values at other rationals?}

\end{exercisebox}
\subsubsection*{Problem 3. The Periodic Excess}
\addcontentsline{toc}{subsubsection}{Problem 3. The Periodic Excess}

\begin{exercisebox}[hard]

Let $a>1$ be a fixed real number. Find all functions $f$ that map
positive real numbers to positive real numbers and satisfy, for all
$x,y>0$,
\[
\begin{aligned}
f(x+y)&\geq f(x)+f(y),\\
f(xy)&\leq f(x)f(y),\\
f(x+a)&\leq f(x)+a.
\end{aligned}
\]

\begin{flushright}
\textit{Proposed by Vu Hung Nguyen, 2026}
\end{flushright}

\upsidehint{Iterate the shift by $a$ and compare it with what repeated addition tells you.}

\end{exercisebox}
\subsubsection*{Problem 4. The Dyadic Lock}
\addcontentsline{toc}{subsubsection}{Problem 4. The Dyadic Lock}

\begin{exercisebox}[medium]

Find all functions $f:(0,\infty)\to(0,\infty)$ such that
\[
f(2)f(f(2))\leq4
\]
and, for all $x,y>0$,
\[
f(x+y)\geq f(x)+f(y)
\geq f(xy)\left(\frac1{f(x)}+\frac1{f(y)}\right).
\]

\begin{flushright}
\textit{Proposed by Vu Hung Nguyen, 2026}
\end{flushright}

\upsidehint{What does the inequality at equal inputs tell you, and how might it interact with the bound at $2$?}

\end{exercisebox}
\subsubsection*{Problem 5. The Triadic Balance}
\addcontentsline{toc}{subsubsection}{Problem 5. The Triadic Balance}

\begin{exercisebox}[medium]

Find all functions $f:(0,\infty)\to(0,\infty)$ such that
\[
f(3)+f\left(\frac13\right)\le \frac{10}{3},
\]
and, for all $x,y>0$,
\[
f(x+y)\ge f(x)+f(y),
\qquad
f(xy)\le f(x)f(y).
\]

\begin{flushright}
\textit{Composed by Vu Hung Nguyen, 2026}
\end{flushright}

\upsidehint{Use repeated addition and multiplication at $3$ and $1/3$ to exploit the normalization bound.}

\end{exercisebox}
\subsubsection*{Problem 34. The Ratio--Product Squeeze}
\addcontentsline{toc}{subsubsection}{Problem 34. The Ratio--Product Squeeze}

\begin{exercisebox}[hard]

Find all functions $f:(0,\infty)\to(0,\infty)$ such that, for all
$x,y>0$,
\[
x f(x+y)\ge (x+y)f(x),
\qquad
f(xy)\le f(x)f(y),
\]
and for which there exists $a>1$ satisfying
\[
f(f(a))\le f(a).
\]

\begin{flushright}
\textit{Original problem for this collection; not an official IMO problem.}
\end{flushright}

\upsidehint{Normalize by $x$ to reveal monotonicity and submultiplicativity, then propagate equality from an interval around $1$ using integer roots.}

\end{exercisebox}
\subsubsection*{Problem 35. The Subadditive--Supermultiplicative Squeeze}
\addcontentsline{toc}{subsubsection}{Problem 35. The Subadditive--Supermultiplicative Squeeze}

\begin{exercisebox}[hard]

Find all functions $f$ that map positive real numbers to positive real
numbers, satisfy $f(f(a))=a$ for some $a>1$, and, for all $x,y>0$,
\[
\boxed{
f(x)\leq f(x+y)\leq f(x)+f(y)
\leq f(xy)\left(\frac1{f(x)}+\frac1{f(y)}\right).
}
\]

\begin{flushright}
\textit{Composed by Vu Hung Nguyen, 2026}
\end{flushright}

\upsidehint{Extract monotonicity, subadditivity, and supermultiplicativity; then use the anchor to squeeze $f(x)$ against $x$.}

\end{exercisebox}
\subsubsection*{Problem 38. The Injective Staircase}
\addcontentsline{toc}{subsubsection}{Problem 38. The Injective Staircase}

\begin{exercisebox}[medium]

Determine all injective functions $f:\mathbb Z_{>0}\to\mathbb Z_{>0}$ such that
\[
\boxed{\quad f\bigl(m+f(n)\bigr)\le n+f(m+1)-1\quad}
\]
for all positive integers $m,n$.

\begin{flushright}
\textit{Composed by Vu Hung Nguyen, 2026}
\end{flushright}

\upsidehint{Put $n=1$ and $n=2$ to get "staircase" bounds. Then compare how many distinct values must fit under a linear bound.}

\end{exercisebox}

\subsubsection*{Problem 41. The Prime-Exponent Trap}
\addcontentsline{toc}{subsubsection}{Problem 41. The Prime-Exponent Trap}

\begin{exercisebox}[hard]

Determine all functions $f:\mathbb Q_{>0}\to\mathbb Q_{>0}$ such that, for all positive rational numbers $x,y$,
\[
f(xy)=f(x)f(y)
\]
and
\[
\boxed{\qquad f\bigl(f(x)^3\bigr)\ge x^2f(x).\qquad}
\]

\begin{flushright}
\textit{Composed by Vu Hung Nguyen, 2026}
\end{flushright}

\upsidehint{Replace $x$ by $1/x$, then investigate how prime exponents behave along the iterates of $f(x)/x$.}

\end{exercisebox}

\subsubsection*{Problem 42. Bounded Errors, Exact Answer}
\addcontentsline{toc}{subsubsection}{Problem 42. Bounded Errors, Exact Answer}

\begin{exercisebox}[medium]

Determine all functions $f:\mathbb R\to\mathbb R$ satisfying $f(2)>1$ and, for all real numbers $x,y$,
\[
\boxed{
\begin{aligned}
\bigl|f(x+y)-f(x)-f(y)\bigr|&\le1,\\
\bigl|f(xy)-f(x)f(y)\bigr|&\le1.
\end{aligned}}
\]

No continuity, monotonicity, or boundedness is assumed.

\begin{flushright}
\textit{Proposed by Vu Hung Nguyen, 2026}
\end{flushright}

\upsidehint{First prove that $f$ is unbounded, then compare the approximations to $f((xy)t)$ and $f(x(yt))$.}

\end{exercisebox}

\subsubsection*{Problem 46. The Missing Middle Value}
\addcontentsline{toc}{subsubsection}{Problem 46. The Missing Middle Value}

\begin{exercisebox}[hard]

Determine all functions $f:\mathbb Z_{>0}\to\mathbb Z_{>0}$ satisfying, for all positive integers $m,n$,
\[
\boxed{
\begin{aligned}
f(mn)&=f(m)f(n),\\
f(m+n)&\ge f(m)+f(n),\\
f(f(n))&\le n^4.
\end{aligned}}
\]

\begin{flushright}
\textit{Composed by Vu Hung Nguyen, 2026}
\end{flushright}

\upsidehint{First show $f(2)\in\{2,3,4\}$; to eliminate $3$, compare $3^2$ with $2^3$, and $3^5$ with $2^8$.}

\end{exercisebox}

\subsubsection*{Problem 49. The Prime-Pairing Machine}
\addcontentsline{toc}{subsubsection}{Problem 49. The Prime-Pairing Machine}

\begin{exercisebox}[hard]

Determine all functions $f:\mathbb Q_{>0}\to\mathbb Q_{>0}$ satisfying
\[
\boxed{
f(xy)=f(x)f(y),
\qquad
f(f(x))\ge x^2
}
\]
for all positive rational numbers $x,y$, and such that
\[
\boxed{\qquad f(n)\in\mathbb Z_{>0}\quad\text{whenever }n\in\mathbb Z_{>0}.\qquad}
\]

\begin{flushright}
\textit{Proposed by Vu Hung Nguyen, 2026}
\end{flushright}

\upsidehint{Use reciprocals to force equality, then track what happens to each prime divisor of $f(p)$.}

\end{exercisebox}

\subsubsection*{Problem 53. When Even Numbers Become Smaller}
\addcontentsline{toc}{subsubsection}{Problem 53. When Even Numbers Become Smaller}

\begin{exercisebox}[hard]

Determine all functions $f:\mathbb Q\to[0,\infty)$ satisfying
\[
f(0)=0,\qquad f(2)=\frac12,
\]
and, for all rational numbers $x,y$,
\[
\boxed{
\begin{aligned}
f(xy)&=f(x)f(y),\\
f(x+y)&\le f(x)+f(y).
\end{aligned}}
\]

\begin{flushright}
\textit{Proposed by Vu Hung Nguyen, 2026}
\end{flushright}

\upsidehint{Use binary expansions to bound $f(n)$ uniformly on positive integers, then apply the bound to $n^r$.}

\end{exercisebox}

\subsubsection*{Problem 58. The Four-Step Prime Labyrinth}
\addcontentsline{toc}{subsubsection}{Problem 58. The Four-Step Prime Labyrinth}

\begin{exercisebox}[hard]

Determine all functions $f:\mathbb Q_{>0}\to\mathbb Q_{>0}$ satisfying
\[
f(xy)=f(x)f(y)
\]
for all positive rational $x,y$, such that $f(n)$ is a positive integer whenever $n$ is a positive integer, and
\[
\boxed{\qquad
f(f(f(f(x))))\ge x^4
\qquad(x\in\mathbb Q_{>0}).
\qquad}
\]

\begin{flushright}
\textit{Composed by Vu Hung Nguyen, 2026}
\end{flushright}

\upsidehint{Follow prime divisors through four applications of $f$, then restart the same chain at its second prime.}

\end{exercisebox}

\subsubsection*{Problem 59. The Hidden Product Rule}
\addcontentsline{toc}{subsubsection}{Problem 59. The Hidden Product Rule}

\begin{exercisebox}[hard]

Determine all functions $f:\mathbb R\to\mathbb R$ satisfying
\[
\boxed{
\begin{aligned}
f(x+y)&\ge f(x)+f(y),\\
f(xy)+xy&\ge x f(y)+y f(x)
\end{aligned}}
\]
for all real numbers $x,y$, and
\[
\boxed{\qquad
f(f(t))+t\ge2f(t)
\qquad(t\ge0).
\qquad}
\]

\begin{flushright}
\textit{Proposed by Vu Hung Nguyen, 2026}
\end{flushright}

\upsidehint{Set $g(x)=f(x)-x$, use negative inputs to force equalities, then investigate $g(g(x))$.}

\end{exercisebox}

\subsubsection*{Problem 62. The Hidden Square Root}
\addcontentsline{toc}{subsubsection}{Problem 62. The Hidden Square Root}

\begin{exercisebox}[hard]

Determine all functions $f:\mathbb Q\to[0,\infty)$ satisfying
\[
f(0)=0,\qquad f(1)=1,\qquad f(4)=2,
\]
and, for all rational numbers $x,y$,
\[
\boxed{
\begin{aligned}
f(x+y)&\le f(x)+f(y),\\
f(xy)&\ge f(x)f(y),
\end{aligned}}
\]
together with
\[
\boxed{\qquad
f(x)f(1/x)\le1
\qquad(x\ne0).
\qquad}
\]

\begin{flushright}
\textit{Proposed by Vu Hung Nguyen, 2026}
\end{flushright}

\upsidehint{First force multiplicativity, then compare base-$2$ and base-$b$ expansions of large powers.}

\end{exercisebox}

\subsubsection*{Problem 64. The Invisible Conjugate}
\addcontentsline{toc}{subsubsection}{Problem 64. The Invisible Conjugate}

\begin{exercisebox}[hard]

Let
\[
K=\{a+b\sqrt2:a,b\in\mathbb Q\}.
\]
Determine all functions $f:K\to[0,\infty)$ satisfying
\[
f(0)=0,\qquad f(1)=1,\qquad f(2)=2,
\]
and, for all $x,y\in K$,
\[
\boxed{
\begin{aligned}
f(x+y)&\le f(x)+f(y),\\
f(xy)&=f(x)f(y).
\end{aligned}}
\]

No continuity, monotonicity, or symmetry is assumed.

\begin{flushright}
\textit{Composed by Vu Hung Nguyen, 2026}
\end{flushright}

\upsidehint{For rational $r$ close to $\sqrt2$, examine the product $f(\sqrt2-r)f(\sqrt2+r)$.}

\end{exercisebox}

\subsubsection*{Problem 65. Three Hidden Primes}
\addcontentsline{toc}{subsubsection}{Problem 65. Three Hidden Primes}

\begin{exercisebox}[hard]

Determine all functions $f:\mathbb Q\to[0,\infty)$ satisfying
\[
f(0)=0,\qquad f(30)=\frac1{30},
\]
and, for all rational numbers $x,y$,
\[
\boxed{
\begin{aligned}
f(x+y)&\le f(x)+f(y),\\
f(xy)&=f(x)f(y).
\end{aligned}}
\]

No continuity, monotonicity, or integer-valuedness is assumed.

\begin{flushright}
\textit{Composed by Vu Hung Nguyen, 2026}
\end{flushright}

\upsidehint{Use base-$30$ expansions to bound $f$ on integers, then show that two distinct primes cannot both have values below $1$.}

\end{exercisebox}

\subsubsection*{Problem 68. The Conditioned Bilinear Sign Test}
\addcontentsline{toc}{subsubsection}{Problem 68. The Conditioned Bilinear Sign Test}

\begin{exercisebox}[hard]

Find all functions $f:\mathbb R\to\mathbb R$ satisfying $f(0)=0$, $f(1)=1$, and,
for all $a,b,c,d\in\mathbb R$:
\begin{enumerate}
    \item If $ab+cd\geq0$, then
          \[
          f(a)f(b)+f(c)f(d)\geq0.
          \]
    \item If $ab+cd\leq0$, then
          \[
          f(a)f(b)+f(c)f(d)\leq0.
          \]
\end{enumerate}

\begin{flushright}
\textit{Composed by Vu Hung Nguyen, 2026.}
\end{flushright}

\upsidehint{Use the boundary case $ab+cd=0$ to obtain exact identities. Try
$(t,1,-t,1)$ and $(x,y,-xy,1)$, then take logarithms on the positive half-line.}

\end{exercisebox}

\subsection*{Part II. Reciprocal and Harmonic Transforms}
\addcontentsline{toc}{subsection}{Part II. Reciprocal and Harmonic Transforms}

\subsubsection*{Problem 6. The Sum--Product Squeeze}
\addcontentsline{toc}{subsubsection}{Problem 6. The Sum--Product Squeeze}

\begin{exercisebox}[medium]

Find all functions $f$ of a single variable that map
positive real numbers to positive real numbers and
satisfy, for all $x,y>0$,
\[
f(x+y)\bigl(f(x)+f(y)\bigr)
\le 2f(2)f(x)f(y)
\le f(xy).
\]

\begin{flushright}
\textit{Proposed by Vu Hung Nguyen, 2026}
\end{flushright}

\upsidehint{What do the inequalities become for carefully chosen pairs such as $(x,1)$ or $(x,x)$?}

\end{exercisebox}
\subsubsection*{Problem 7. The Harmonic Squeeze}
\addcontentsline{toc}{subsubsection}{Problem 7. The Harmonic Squeeze}

\begin{exercisebox}[medium]

Find all functions $f:(0,\infty)\to(0,\infty)$ such that
\[
f(2)\ge \frac12
\]
and, for all $x,y>0$,
\[
f(x+y)\bigl(f(x)+f(y)\bigr)
\le f(x)f(y)
\le f(xy).
\]

\begin{flushright}
\textit{Proposed by Vu Hung Nguyen, 2026}
\end{flushright}

\upsidehint{Set $x=y$ first, then compare the resulting doubling relation with the product bound.}

\end{exercisebox}
\subsubsection*{Problem 8. The Multiplicative Mirror}
\addcontentsline{toc}{subsubsection}{Problem 8. The Multiplicative Mirror}

\begin{exercisebox}[medium]

Find all functions $f$ of a single variable that map
positive real numbers to positive real numbers such that
\[
f(x)f\left(\frac1x\right)=1,
\qquad f(f(2))\le 2,
\]
and, for all $x,y>0$,
\[
\frac{f(x)f(y)}
{f(x+y)\bigl(f(x)+f(y)\bigr)}
\ge 1 \ge
\frac{f(x)^2}{f(xy)f(x/y)}.
\]

\begin{flushright}
\textit{Proposed by Vu Hung Nguyen, 2026}
\end{flushright}

\upsidehint{What do the two inequalities say after replacing both variables by their reciprocals?}

\end{exercisebox}
\subsubsection*{Problem 9. The Reflected Product}
\addcontentsline{toc}{subsubsection}{Problem 9. The Reflected Product}

\begin{exercisebox}[medium]

Find all functions $f:(0,\infty)\to(0,\infty)$ such that
\[
f(x)f\left(\frac1x\right)=1,
\qquad f(2)\leq 2,
\]
and, for all $x,y>0$,
\[
f(x+y)\geq f(x)+f(y),
\qquad
f(x)^2\leq f(xy)f\left(\frac{x}{y}\right).
\]

\begin{flushright}
\textit{Composed by Vu Hung Nguyen, 2026}
\end{flushright}

\upsidehint{Combine the reciprocal identity with the product inequality after interchanging $x$ and $y$.}

\end{exercisebox}
\subsubsection*{Problem 10. The Reciprocal Pinch}
\addcontentsline{toc}{subsubsection}{Problem 10. The Reciprocal Pinch}

\begin{exercisebox}[medium]

Find all functions $f:(0,\infty)\to(0,\infty)$ such that
\[
f(x)f\left(\frac1x\right)=1,
\qquad
f(2)+f\left(\frac12\right)\le\frac52,
\]
and, for all $x,y>0$,
\[
f(x+y)\ge f(x)+f(y),
\qquad
f(xy)\ge f(x)f(y).
\]

\begin{flushright}
\textit{Proposed by Vu Hung Nguyen, 2026}
\end{flushright}

\upsidehint{Use reciprocal symmetry at $2$ and compare the resulting values with their given sum bound.}

\end{exercisebox}
\subsubsection*{Problem 11. The Harmonic--Product Squeeze}
\addcontentsline{toc}{subsubsection}{Problem 11. The Harmonic--Product Squeeze}

\begin{exercisebox}[hard]

Find all functions $f$ that map positive real numbers to positive real
numbers, provided that $f(f(c))\geq c$ for some $c\in(0,1)$, and, for
all $x,y>0$,
\[
f\!\left(\frac{xy}{x+y}\right)
\leq\frac{f(x)f(y)}{f(x)+f(y)},
\qquad
f(xy)\geq f(x)f(y).
\]

\begin{flushright}
\textit{Composed by Vu Hung Nguyen, 2026}
\end{flushright}

\upsidehint{Try reciprocal substitutions, then bring the special point $c$ into play.}

\end{exercisebox}
\subsubsection*{Problem 54. The Harmonic Involution}
\addcontentsline{toc}{subsubsection}{Problem 54. The Harmonic Involution}

\begin{exercisebox}[medium]

Determine all functions $f:(0,\infty)\to(0,\infty)$ satisfying
\[
\boxed{\qquad f(f(x))=x\qquad}
\]
for every $x>0$, and
\[
\boxed{\qquad
f(x+y)\bigl(f(x)+f(y)\bigr)\le f(x)f(y)
\qquad}
\]
for all $x,y>0$.

No continuity, monotonicity, or normalisation is assumed.

\begin{flushright}
\textit{Proposed by Vu Hung Nguyen, 2026}
\end{flushright}

\upsidehint{Prove $f(nx)\le f(x)/n$, then apply the same bound with input $f(nx)$ and use the involution.}

\end{exercisebox}

\subsubsection*{Problem 60. The Squared-Logarithm Trap}
\addcontentsline{toc}{subsubsection}{Problem 60. The Squared-Logarithm Trap}

\begin{exercisebox}[hard]

Determine all functions $f:(0,\infty)\to[1,\infty)$ satisfying
\[
f(2)=2,\qquad f(t)\le t\quad(1\le t\le2),
\]
and, for all positive real numbers $x,y$,
\[
\boxed{\qquad
f(xy)\,f(x/y)\le f(x)^2f(y)^2,
\qquad} \tag{1}
\]
together with
\[
\boxed{\qquad
f(x^2)\ge f(x)^4.
\qquad} \tag{2}
\]

No continuity or monotonicity is assumed.

\begin{flushright}
\textit{Composed by Vu Hung Nguyen, 2026}
\end{flushright}

\upsidehint{Set $g(t)=\log_2 f(2^t)$, then use the substitution $(u,v)\mapsto(u+v,u-v)$.}

\end{exercisebox}

\subsubsection*{Problem 61. The Hidden Entropy Function}
\addcontentsline{toc}{subsubsection}{Problem 61. The Hidden Entropy Function}

\begin{exercisebox}[hard]

Determine all functions $f:(0,\infty)\to[0,\infty)$ satisfying
\[
\boxed{\qquad
f(xy)\ge x f(y)+y f(x)+(x-1)(y-1)
\qquad}
\]
for all positive real numbers $x,y$.

No continuity, monotonicity, or normalisation is assumed.

\begin{flushright}
\textit{Composed by Vu Hung Nguyen, 2026}
\end{flushright}

\upsidehint{Investigate $g(x)=\bigl(f(x)+x-1\bigr)/x$, and split $x$ into arbitrarily many equal multiplicative factors.}

\end{exercisebox}

\subsubsection*{Problem 66. The Irrational Iteration Root}
\addcontentsline{toc}{subsubsection}{Problem 66. The Irrational Iteration Root}

\begin{exercisebox}[hard]

Determine all functions $f:(0,\infty)\to(0,\infty)$ satisfying
\[
\boxed{
\begin{aligned}
f(x+y)&\ge f(x)+f(y),\\
f(xy)&\le f(x)f(y)
\end{aligned}}
\]
for all positive real numbers $x,y$, together with
\[
\boxed{\qquad f(f(x))=x^2\qquad(x>0).\qquad}
\]

No continuity, monotonicity, or normalisation is assumed.

\begin{flushright}
\textit{Composed by Vu Hung Nguyen, 2026}
\end{flushright}

\upsidehint{Apply $f$ to $f(xy)\le f(x)f(y)$, then apply submultiplicativity to $f(x)$ and $f(y)$.}

\end{exercisebox}


\subsection*{Part III. Iteration, Fixed Points, and Growth}
\addcontentsline{toc}{subsection}{Part III. Iteration, Fixed Points, and Growth}

\subsubsection*{Problem 12. IMO 1977 -- Problem 6}
\addcontentsline{toc}{subsubsection}{Problem 12. IMO 1977 -- Problem 6}

\begin{exercisebox}[medium]

Let $f$ be a function that maps positive integers to positive integers
and satisfies, for all $n\in\mathbb Z_{>0}$,
\[
f(n+1)>f(f(n)).
\]
Prove that
\[
f(n)=n\qquad(n\in\mathbb Z_{>0}).
\]

\medskip
\noindent
\textit{Source: International Mathematical Olympiad 1977 (19th IMO),
Day 2, Problem 6.}

\upsidehint{Can you compare consecutive values of $f$ by combining the condition with iteration?}

\end{exercisebox}
\subsubsection*{Problem 13. The Fixed-Point Lever}
\addcontentsline{toc}{subsubsection}{Problem 13. The Fixed-Point Lever}

\begin{exercisebox}[hard]

Find all functions $f$ that map positive real numbers
to positive real numbers, if there exists $a>0$ such that
$f(f(a))=a$, and for all $x,y>0$,
\[
f(x)\bigl(x+f(y)\bigr)
\geq f(f(x))\bigl(f(x)+y\bigr).
\]

\begin{flushright}
\textit{Composed by Vu Hung Nguyen, 2026}
\end{flushright}

\upsidehint{Try using the special point $a$ in each variable, where $f(f(a))=a$.}

\end{exercisebox}
\subsubsection*{Problem 36. The Multiplicative Mirror: Positive Recurrence}
\addcontentsline{toc}{subsubsection}{Problem 36. The Multiplicative Mirror: Positive Recurrence}

\begin{exercisebox}[hard]

Determine all functions $f:(0,\infty)\to(0,\infty)$ satisfying
\[
f(1)=1,\qquad f(f(x))=x
\]
for every $x>0$, and
\[
f(xy)+f\left(\frac yx\right)
\le \left(x+\frac1x\right)f(y)
\]
for all $x,y>0$.

\begin{flushright}
\textit{Proposed by Vu Hung Nguyen, 2026}
\end{flushright}

\upsidehint{Follow the sequence $f(t^n)$ for $t>1$. Use positivity to obtain bounds, then apply the involution to those bounds.}

\end{exercisebox}
\subsubsection*{Problem 37. The Crossed Contraction Lock}
\addcontentsline{toc}{subsubsection}{Problem 37. The Crossed Contraction Lock}

\begin{exercisebox}[medium]

Determine all functions $f:(0,\infty)\to(0,\infty)$ such that
\[
f(f(x))\ge x
\]
for all $x>0$, and
\[
\boxed{f\bigl(xf(y)\bigr)+f\bigl(yf(x)\bigr)\le 2xy}
\]
for all $x,y>0$.

\begin{flushright}
\textit{Proposed by Vu Hung Nguyen, 2026}
\end{flushright}

\upsidehint{Put $y=1$ and let $c=f(1)$. First obtain a linear upper bound for $f$, then iterate it against $f(f(x))\ge x$.}

\end{exercisebox}
\subsubsection*{Problem 14. IMO 2023 Shortlist -- Algebra A4}
\addcontentsline{toc}{subsubsection}{Problem 14. IMO 2023 Shortlist -- Algebra A4}

\begin{exercisebox}[medium]

Find all functions $f$ of a single variable that map
positive real numbers to positive real numbers such that,
for all $x,y>0$,
\[
 x\bigl(f(x)+f(y)\bigr)
 \ge \bigl(f(f(x))+y\bigr)f(y).
\]

\noindent
\textit{Source: This problem is taken from the
IMO 2023 Shortlist, Algebra Problem A4.
The wording has been slightly adapted without
changing the mathematical content.}

\upsidehint{Fix $x$ and let $y$ vary; what does the inequality force about the possible size of $f(y)$?}

\end{exercisebox}
\subsubsection*{Problem 15. FEOO 2020 Shortlist -- Algebra A1}
\addcontentsline{toc}{subsubsection}{Problem 15. FEOO 2020 Shortlist -- Algebra A1}

\begin{exercisebox}[medium]

Find all functions $f$ of a single variable that
map real numbers to real numbers and satisfy,
for all $x,y\in\mathbb R$,
\[
2xyf(xy)\geq x^2y^2+xf(x)f(y^2).
\]

\medskip
\noindent
\textit{Source: Functional Equations Online Olympiad
(FEOO) 2020, Shortlist, Algebra Problem A1.
The wording has been slightly adapted without
changing the mathematical content.}

\upsidehint{Start with $y=0$ and exploit the fact that $y^2$ stays nonnegative.}

\end{exercisebox}
\subsubsection*{Problem 16. The Iterated Ratio Squeeze}
\addcontentsline{toc}{subsubsection}{Problem 16. The Iterated Ratio Squeeze}

\begin{exercisebox}[hard]

Let $a>1$ be a fixed real number. Find all functions $f$ that map
positive real numbers to positive real numbers, satisfy $f(a)=a$, and,
for all $x,y>0$,
\[
\frac{f(x)}{f(y)}+\frac{f(y)}{f(x)}
\geq\frac{x}{y}+\frac{y}{x}
\geq\frac{f^{\circ 3}(x)}{f^{\circ 3}(y)}
+\frac{f^{\circ 3}(y)}{f^{\circ 3}(x)}.
\]
Here $f^{\circ 3}(x)=f(f(f(x)))$.

\begin{flushright}
\textit{Composed by Vu Hung Nguyen, 2026}
\end{flushright}

\upsidehint{Fix one variable at $a$ and follow what the squeeze says through repeated applications of $f$.}

\end{exercisebox}
\subsubsection*{Problem 17. IMO 2022 Shortlist -- Algebra A3}
\addcontentsline{toc}{subsubsection}{Problem 17. IMO 2022 Shortlist -- Algebra A3}

\begin{exercisebox}[medium]

Find all functions $f$ of a single variable that map
positive real numbers to positive real numbers such
that, for every $x>0$, there exists exactly one
$y>0$ satisfying
\[
xf(y)+yf(x)\leq2.
\]

\medskip
\noindent
\textit{Source: International Mathematical Olympiad
2022 (63rd IMO), Shortlist, Algebra Problem A3,
proposed by the Netherlands. The wording has been
slightly adapted without changing the mathematical
content.}

\upsidehint{For a fixed $x$, what shape does the set of admissible $y$ have, and how does uniqueness restrict it?}

\end{exercisebox}
\subsubsection*{Problem 18. The Geometric-Mean Trap}
\addcontentsline{toc}{subsubsection}{Problem 18. The Geometric-Mean Trap}

\begin{exercisebox}[hard]

Find all functions $f$ of a single variable that map positive real
numbers to positive real numbers and satisfy the following property:
for every $x>0$, there exists exactly one $y>0$ such that
\[
\frac{x}{f(y)}+\frac{y}{f(x)}
\leq\frac{f(x)}x+\frac{f(y)}y.
\]

\begin{flushright}
\textit{Composed by Vu Hung Nguyen, 2026}
\end{flushright}

\upsidehint{For fixed $x$, rewrite the condition using the ratios $f(x)/x$ and $f(y)/y$; what does uniqueness imply?}

\end{exercisebox}
\subsubsection*{Problem 19. The Reciprocal-Slope Pinch}
\addcontentsline{toc}{subsubsection}{Problem 19. The Reciprocal-Slope Pinch}

\begin{exercisebox}[hard]

Find all functions $f:(0,\infty)\to(0,\infty)$ such that, for all
$x,y>0$,
\[
xy+f\bigl(xf(x)\bigr)\leq x f(x+y).
\]

\begin{flushright}
\textit{Composed by Vu Hung Nguyen, 2026}
\end{flushright}

\upsidehint{For fixed $x$, compare the lower bounds at nearby inputs and watch how the term $xf(x)$ enters.}

\end{exercisebox}
\subsubsection*{Problem 20. The Cross-Dilated Shift}
\addcontentsline{toc}{subsubsection}{Problem 20. The Cross-Dilated Shift}

\begin{exercisebox}[hard]

Find all functions $f:(0,\infty)\to(0,\infty)$ such that, for all
$x,y>0$,
\[
x^2+f\bigl(xf(y)\bigr)\leq x f(x+y).
\]

\begin{flushright}
\textit{Proposed by Vu Hung Nguyen, 2026}
\end{flushright}

\upsidehint{Compare the inequality at exchanged or rescaled inputs; what role does the sum $x+y$ play?}

\end{exercisebox}
\subsubsection*{Problem 40. The Monotonicity Ratchet}
\addcontentsline{toc}{subsubsection}{Problem 40. The Monotonicity Ratchet}

\begin{exercisebox}[medium]

Determine all functions $f:(0,\infty)\to(0,\infty)$ such that
\[
\boxed{\quad f\bigl(x+f(y)\bigr)\ge f(x)+y\qquad\text{and}\qquad f\bigl(f(y)\bigr)\le y\quad}
\]
for all $x,y>0$.

\begin{flushright}
\textit{Composed by Vu Hung Nguyen, 2026}
\end{flushright}

\upsidehint{Show that $f(t)>y$ whenever $t>f(y)$, then substitute $y=f(w)$ to prove $f$ is strictly increasing.}

\end{exercisebox}

\subsubsection*{Problem 43. The Moving Zero}
\addcontentsline{toc}{subsubsection}{Problem 43. The Moving Zero}

\begin{exercisebox}[hard]

Determine all functions $f:\mathbb R\to[0,\infty)$ satisfying
\[
\boxed{\quad
f(x+y)+f(x-y)\ge 2f(x)+2y^2
\quad} \tag{1}
\]
for all real $x,y$, and
\[
\boxed{\quad f(f(x))\le x^4\quad} \tag{2}
\]
for every real $x$.

No continuity or monotonicity is assumed.

\begin{flushright}
\textit{Proposed by Vu Hung Nguyen, 2026}
\end{flushright}

\upsidehint{Locate a zero $a$, repeatedly bisect the segment from $a$ to $t$, and use the resulting lower bound inside the composition inequality.}

\end{exercisebox}

\subsubsection*{Problem 44. The Hidden Golden Exponent}
\addcontentsline{toc}{subsubsection}{Problem 44. The Hidden Golden Exponent}

\begin{exercisebox}[medium]

Determine all functions $f:(0,\infty)\to(0,\infty)$ satisfying
\[
f(xy)=f(x)f(y)
\]
and
\[
\boxed{\qquad
f\bigl(x+f(y)\bigr)\ge f(x)+y f(y)
\qquad}
\]
for all positive real numbers $x,y$.

No continuity or monotonicity is assumed.

\begin{flushright}
\textit{Composed by Vu Hung Nguyen, 2026}
\end{flushright}

\upsidehint{Set $y=1$, use multiplicativity to obtain superadditivity, and only then investigate power functions.}

\end{exercisebox}

\subsubsection*{Problem 45. The Conjugate-Radical Lock}
\addcontentsline{toc}{subsubsection}{Problem 45. The Conjugate-Radical Lock}

\begin{exercisebox}[medium]

Determine all functions $f:[0,\infty)\to[1,\infty)$ satisfying $f(1)=2$ and, for all $x,y\ge0$,
\[
\boxed{
f(x+y)\ge
f(x)f(y)+
\sqrt{\bigl(f(x)^2-1\bigr)\bigl(f(y)^2-1\bigr)}
}
\tag{1}
\]
and
\[
\boxed{\qquad f(2x)\le 2f(x)^2-1.\qquad} \tag{2}
\]

No continuity or monotonicity is assumed.

\begin{flushright}
\textit{Composed by Vu Hung Nguyen, 2026}
\end{flushright}

\upsidehint{Investigate $g(x)=\log\!\left(f(x)+\sqrt{f(x)^2-1}\right)$, then combine superadditivity with exact doubling.}

\end{exercisebox}

\subsubsection*{Problem 48. The One-Point Cubic Lock}
\addcontentsline{toc}{subsubsection}{Problem 48. The One-Point Cubic Lock}

\begin{exercisebox}[hard]

Determine all functions $f:(0,\infty)\to(0,\infty)$ satisfying
\[
\boxed{
\begin{aligned}
f(x+y)&\ge f(x)+f(y)+3xy(x+y),\\
f(xy)&\le f(x)f(y)
\end{aligned}}
\]
for all positive real numbers $x,y$, together with
\[
\boxed{\qquad f(f(2))\le512.\qquad}
\]

No continuity or monotonicity is assumed.

\begin{flushright}
\textit{Proposed by Vu Hung Nguyen, 2026}
\end{flushright}

\upsidehint{Split $t$ into $n$ equal pieces to prove $f(t)\ge t^3$, then use the composition bound to lock $f(2)$.}

\end{exercisebox}

\subsubsection*{Problem 50. The Golden Staircase}
\addcontentsline{toc}{subsubsection}{Problem 50. The Golden Staircase}

\begin{exercisebox}[hard]

Determine all functions $f:\mathbb Z_{>0}\to\mathbb Z_{>0}$ such that, for all positive integers $m,n$,
\[
\boxed{\qquad
0\le f(m+n)-f(m)-f(n)\le1
\qquad} \tag{1}
\]
and
\[
\boxed{\qquad
\left|f(f(n))-f(n)-n\right|\le1.
\qquad} \tag{2}
\]

\begin{flushright}
\textit{Composed by Vu Hung Nguyen, 2026}
\end{flushright}

\upsidehint{Compare $f(mn)$ with both $m f(n)$ and $n f(m)$ to uncover a limiting slope.}

\end{exercisebox}

\subsubsection*{Problem 52. The Almost-Cosine Sequence}
\addcontentsline{toc}{subsubsection}{Problem 52. The Almost-Cosine Sequence}

\begin{exercisebox}[hard]

Determine all functions $f:\mathbb Z\to\mathbb Z_{>0}$ satisfying $f(1)=2$ and
\[
\boxed{\qquad
\left|f(m+n)+f(m-n)-2f(m)f(n)\right|\le1
\qquad}
\]
for all integers $m,n$.

\begin{flushright}
\textit{Proposed by Vu Hung Nguyen, 2026}
\end{flushright}

\upsidehint{First prove unboundedness, then expand the sum of $f(x+y+z)$, $f(x+y-z)$, $f(x-y+z)$, and $f(x-y-z)$ in two ways.}

\end{exercisebox}

\subsubsection*{Problem 55. The Reversible Difference}
\addcontentsline{toc}{subsubsection}{Problem 55. The Reversible Difference}

\begin{exercisebox}[medium]

Determine all functions $f:\mathbb Z\to\mathbb Z_{\ge0}$ satisfying
\[
f(0)=0,\qquad f(1)=1,\qquad f(f(2))\le16,
\]
and, for all integers $m,n$,
\[
\boxed{\qquad
f(m+3n)+3f(m+n)
\ge 3f(m+2n)+f(m).
\qquad}
\]

\begin{flushright}
\textit{Composed by Vu Hung Nguyen, 2026}
\end{flushright}

\upsidehint{Replace $(m,n)$ by $(m+3n,-n)$, then examine consecutive second differences.}

\end{exercisebox}

\subsubsection*{Problem 57. The Hidden Interval}
\addcontentsline{toc}{subsubsection}{Problem 57. The Hidden Interval}

\begin{exercisebox}[medium]

Determine all functions $f:\mathbb R\to\mathbb R$ satisfying
\[
\boxed{\qquad
\bigl(f(x)-f(y)\bigr)^2
\le (x-y)\bigl(f(x)-f(y)\bigr)
\qquad} \tag{1}
\]
for all real $x,y$, and
\[
\boxed{\qquad
f(f(f(x)))=f(x)
\qquad} \tag{2}
\]
for every real $x$.

No continuity, monotonicity, or boundedness is assumed.

\begin{flushright}
\textit{Proposed by Vu Hung Nguyen, 2026}
\end{flushright}

\upsidehint{Show that $f$ is nondecreasing and cannot increase distances, then determine what happens between two fixed points.}

\end{exercisebox}

\subsubsection*{Problem 67. The Forbidden Coefficient}
\addcontentsline{toc}{subsubsection}{Problem 67. The Forbidden Coefficient}

\begin{exercisebox}[hard]

Determine all functions $f:(0,\infty)\to(0,\infty)$ satisfying, for all $x,y>0$,
\[
\boxed{
0\le f(x+y)-f(x)-f(y)-3xy(x+y)\le(x-y)^2,
}
\tag{1}
\]
\[
\boxed{\qquad f(xy)\le f(x)f(y),\qquad} \tag{2}
\]
and
\[
\boxed{\qquad f(f(1))\le10.\qquad} \tag{3}
\]

No continuity or monotonicity is assumed.

\begin{flushright}
\textit{Composed by Vu Hung Nguyen, 2026}
\end{flushright}

\upsidehint{Subtract $x^3$, shrink the inputs repeatedly, and then test the multiplicative inequality near zero.}

\end{exercisebox}


\subsection*{Part IV. Mean and Norm Squeezes}
\addcontentsline{toc}{subsection}{Part IV. Mean and Norm Squeezes}

\subsubsection*{Problem 21. 2026 IMO Problem 5}
\addcontentsline{toc}{subsubsection}{Problem 21. 2026 IMO Problem 5}

\begin{exercisebox}[medium]

Determine all functions $f:\mathbb{R}_{>0}\to\mathbb{R}_{>0}$ such that,
for all $x,y>0$,
\[
\sqrt{\frac{x^2+f(y)^2}{2}}
\ge \frac{f(x)+y}{2}
\ge \sqrt{x f(y)}.
\]

\begin{flushright}
\textit{Source: 2026 IMO, Problem 5.}
\end{flushright}

\upsidehint{What can you learn by choosing $y$ in relation to $f(x)$ and comparing the two mean bounds?}

\end{exercisebox}
\subsubsection*{Problem 22. The Hidden Quadratic Function}
\addcontentsline{toc}{subsubsection}{Problem 22. The Hidden Quadratic Function}

\begin{exercisebox}[medium]

Find all functions $f$ of a single variable that map
positive real numbers to positive real numbers such that
$f(f(2))\le 16$ and, for all $x,y>0$,
\[
\sqrt{f(x+y)}
\ge \sqrt{f(x)}+\sqrt{f(y)}
\ge \sqrt{f(xy)}
\left(
\frac{1}{\sqrt{f(x)}}+\frac{1}{\sqrt{f(y)}}
\right).
\]

\begin{flushright}
\textit{Composed by Vu Hung Nguyen, 2026}
\end{flushright}

\upsidehint{Try setting $x=y$ and compare the resulting bounds as the input is doubled.}

\end{exercisebox}
\subsubsection*{Problem 23. The Euclidean--Multiplicative Squeeze}
\addcontentsline{toc}{subsubsection}{Problem 23. The Euclidean--Multiplicative Squeeze}

\begin{exercisebox}[hard]

Find all functions $f$ that map positive real numbers to positive real
numbers, satisfy
\[
f(f(\sqrt2))\leq\sqrt2,
\]
and, for all $x,y>0$,
\[
f\!\left(\sqrt{x^2+y^2}\right)^2
\geq f(x)^2+f(y)^2
\geq f(xy)^2\left(\frac1{f(x)^2}+\frac1{f(y)^2}\right).
\]

\begin{flushright}
\textit{Composed by Vu Hung Nguyen, 2026}
\end{flushright}

\upsidehint{Can you choose inputs so that the Euclidean combination and the product are closely related?}

\end{exercisebox}
\subsubsection*{Problem 24. The Quadratic Mean Squeeze}
\addcontentsline{toc}{subsubsection}{Problem 24. The Quadratic Mean Squeeze}

\begin{exercisebox}[hard]

Find all functions $f$ that map positive real numbers to positive real
numbers, satisfy $f(f(a))\leq a^4$ for some $a>0$, and, for all
$x,y>0$,
\[
\sqrt{\frac{x^4+f(y)^2}{2}}
\geq\frac{f(x)+y^2}{2}
\geq x\sqrt{f(y)}.
\]

\begin{flushright}
\textit{Composed by Vu Hung Nguyen, 2026}
\end{flushright}

\upsidehint{Look for choices of $x$ and $y$ that make the two mean bounds especially tight.}

\end{exercisebox}
\subsubsection*{Problem 47. The Three-Variable Quadratic Trap}
\addcontentsline{toc}{subsubsection}{Problem 47. The Three-Variable Quadratic Trap}

\begin{exercisebox}[hard]

Determine all functions $f:\mathbb R\to[0,\infty)$ satisfying
\[
f(1)=1,\qquad f(t)\le1\quad(0\le t\le1),
\]
and, for all real $x,y,z$,
\[
\boxed{
f(x+y)+f(y+z)+f(z+x)
\le f(x+y+z)+f(x)+f(y)+f(z),
}
\tag{1}
\]
together with
\[
\boxed{\qquad f(2x)\ge4f(x).\qquad} \tag{2}
\]

No continuity is assumed.

\begin{flushright}
\textit{Proposed by Vu Hung Nguyen, 2026}
\end{flushright}

\upsidehint{Use $(x,y,z)=(t,t,-t)$, then $(x,y,z)=(x,y,-y)$, and finally replace $(x,y)$ by $(x+y,x-y)$.}

\end{exercisebox}

\subsubsection*{Problem 51. The Logarithm Hidden in a Chain}
\addcontentsline{toc}{subsubsection}{Problem 51. The Logarithm Hidden in a Chain}

\begin{exercisebox}[medium]

Determine all functions $f:(0,\infty)\to\mathbb R$ satisfying
\[
\boxed{\qquad
f(xy)\le f(x)+f(y)\le x+y-2
\qquad}
\]
for all positive real numbers $x,y$.

\begin{flushright}
\textit{Proposed by Vu Hung Nguyen, 2026}
\end{flushright}

\upsidehint{First find $f(1)$, then apply the first inequality repeatedly to $n$ identical factors $x^{1/n}$.}

\end{exercisebox}

\subsubsection*{Problem 56. The Microscopic Error Trap}
\addcontentsline{toc}{subsubsection}{Problem 56. The Microscopic Error Trap}

\begin{exercisebox}[medium]

Determine all functions $f:\mathbb R\to\mathbb R$ satisfying
\[
\boxed{\qquad
0\le f(x)+f(y)-f(x+y)\le(x-y)^2
\qquad} \tag{1}
\]
for all real $x,y$, and
\[
\boxed{\qquad
f\bigl(f(x^2)\bigr)\ge f(x)^2
\qquad} \tag{2}
\]
for every real $x$.

No continuity, monotonicity, or normalisation is assumed.

\begin{flushright}
\textit{Composed by Vu Hung Nguyen, 2026}
\end{flushright}

\upsidehint{Use exact doubling to shrink the additive error, then study the sign of the composite function $f\circ f$.}

\end{exercisebox}

\subsubsection*{Problem 63. The Hidden Euclidean Length}
\addcontentsline{toc}{subsubsection}{Problem 63. The Hidden Euclidean Length}

\begin{exercisebox}[hard]

Determine all functions
\[
F:\mathbb Q^2\to[0,\infty)
\]
satisfying
\[
F(0,0)=0,\qquad F(1,0)=1,\qquad F(2,0)=2,
\]
and, for all rational numbers $a,b,c,d$,
\[
\boxed{
F(a+c,b+d)\le F(a,b)+F(c,d)
}
\tag{1}
\]
and
\[
\boxed{
F(ac-bd,ad+bc)=F(a,b)F(c,d).
}
\tag{2}
\]

No continuity, monotonicity, or symmetry is assumed.

\begin{flushright}
\textit{Composed by Vu Hung Nguyen, 2026}
\end{flushright}

\upsidehint{First determine $F$ on the coordinate axes, then take large powers under the multiplication $(a,b)(c,d)=(ac-bd,ad+bc)$.}

\end{exercisebox}


\subsection*{Part V. Real-Line and Polynomial Inequalities}
\addcontentsline{toc}{subsection}{Part V. Real-Line and Polynomial Inequalities}

\subsubsection*{Problem 25. Vietnam Mathematical Olympiad 1991}
\addcontentsline{toc}{subsubsection}{Problem 25. Vietnam Mathematical Olympiad 1991}

\begin{exercisebox}[medium]

Find all functions $f$ of a single variable that map
real numbers to real numbers and satisfy, for all
real numbers $x,y,z$,
\[
\frac{f(xy)+f(xz)}{2}
-f(x)f(yz)\geq\frac14.
\]

\noindent
\textit{Source: 29th Vietnamese Mathematical Olympiad,
1991, First Day, Problem 1.}

\upsidehint{What do the substitutions $y=z$ and $y=0$ reveal?}

\end{exercisebox}
\subsubsection*{Problem 26. IMO 2023 Shortlist -- Algebra A2}
\addcontentsline{toc}{subsubsection}{Problem 26. IMO 2023 Shortlist -- Algebra A2}

\begin{exercisebox}[medium]

Let $f$ be a function that maps real numbers
to real numbers such that, for all $x,y\in\mathbb R$,
\[
f(x+y)f(x-y)\geq f(x)^2-f(y)^2.
\]

Suppose that the inequality is strict for
at least one pair of real numbers $(x_0,y_0)$.

Prove that either
\[
f(x)\geq0\quad\text{for all }x\in\mathbb R
\]
or
\[
f(x)\leq0\quad\text{for all }x\in\mathbb R.
\]

\medskip
\noindent
\textit{Source: IMO 2023 Shortlist, Algebra Problem A2,
proposed by Malaysia. The wording has been slightly
adapted without changing the mathematical content.}

\upsidehint{Compare the inequality with the one obtained by swapping $x$ and $y$.}

\end{exercisebox}
\subsubsection*{Problem 27. IMO 2011 -- Problem 3}
\addcontentsline{toc}{subsubsection}{Problem 27. IMO 2011 -- Problem 3}

\begin{exercisebox}[hard]

Let $f$ be a function that maps real numbers to real numbers and
satisfies, for all $x,y\in\mathbb R$,
\[
f(x+y)\leq yf(x)+f(f(x)).
\]
Prove that
\[
f(x)=0\qquad\text{for all }x\leq0.
\]

\medskip
\noindent
\textit{Source: International Mathematical Olympiad 2011 (52nd IMO),
Day 1, Problem 3.}

\upsidehint{Begin with $y=0$ and $y=-x$; what do these say about the values at $0$ and $f(x)$?}

\end{exercisebox}
\subsubsection*{Problem 28. The Degree-Normalized Cubic Squeeze}
\addcontentsline{toc}{subsubsection}{Problem 28. The Degree-Normalized Cubic Squeeze}

\begin{exercisebox}[medium]

Find all functions $f$ that map positive real numbers to real numbers
and satisfy, for all $x,y>0$,
\[
\frac{f(x)}{x^2}+\frac{f(y)}{y^2}
\leq\frac{f(x+y)}{(x+y)^2}
\qquad\text{and}\qquad
\frac{f(x)}{x^3}+\frac{f(y)}{y^3}
\geq\frac{2f(x+y)}{(x+y)^3}.
\]

\upsidehint{Normalize by a power of the input; which choice makes one of the inequalities simpler?}

\end{exercisebox}
\subsubsection*{Problem 29. The Conditioned Cubic Squeeze}
\addcontentsline{toc}{subsubsection}{Problem 29. The Conditioned Cubic Squeeze}

\begin{exercisebox}[hard]

Find all functions $f$ that map real numbers to real numbers such that
$f(0)=0$, $f(1)=1$, and, for all $a,b,c\in\mathbb R$:
\begin{enumerate}
    \item If $a+b+c\geq0$, then
          $f(a^3)+f(b^3)+f(c^3)\geq3f(abc)$.
    \item If $a+b+c\leq0$, then
          $f(a^3)+f(b^3)+f(c^3)\leq3f(abc)$.
\end{enumerate}

\begin{flushright}
\textit{Composed by Vu Hung Nguyen, 2026}
\end{flushright}

\upsidehint{Try triples with a zero entry or two equal entries, and compare the cases of opposite signs.}

\end{exercisebox}
\subsubsection*{Problem 30. The Quadratic--Difference Squeeze}
\addcontentsline{toc}{subsubsection}{Problem 30. The Quadratic--Difference Squeeze}

\begin{exercisebox}[medium]

Find all functions $f$ that map real numbers to real numbers and
satisfy, for all $x,y\in\mathbb R$,
\[
f(x)+f(y)\geq2xy
\qquad\text{and}\qquad
f(x+y)+f(x-y)\leq2f(x)+2y^2.
\]

\begin{flushright}
\textit{Composed by Vu Hung Nguyen, 2026}
\end{flushright}

\upsidehint{Use the first inequality to bound selected values, then feed those bounds into the second.}

\end{exercisebox}
\subsubsection*{Problem 31. The Iterated Parallelogram--Quadratic Squeeze}
\addcontentsline{toc}{subsubsection}{Problem 31. The Iterated Parallelogram--Quadratic Squeeze}

\begin{exercisebox}[hard]

Find all functions $f$ that map real numbers to real numbers, satisfy
\[
f(f(2))\leq16,
\]
and, for all $x,y\in\mathbb R$,
\[
 x^2-y^2+2yf(x)+f(y)-f(f(x))
 \leq\frac{f(x+y)+f(x-y)}{2}
 \leq f(x)+y^2.
\]

\begin{flushright}
\textit{Proposed by Vu Hung Nguyen, 2026}
\end{flushright}

\upsidehint{Start with $x=0$ and $y=0$; what bounds do they give before you iterate $f$?}

\end{exercisebox}
\subsubsection*{Problem 32. The Skew-Symmetric Parallelogram Squeeze}
\addcontentsline{toc}{subsubsection}{Problem 32. The Skew-Symmetric Parallelogram Squeeze}

\begin{exercisebox}[hard]

Find all functions $f$ that map real numbers to real numbers and
satisfy, for all $x,y\in\mathbb R$,
\[
2xy+x^2f(y)-y^2f(x)
\leq\frac{f(x+y)+f(x-y)}{2}
\leq f(x)+y^2.
\]

\begin{flushright}
\textit{Composed by Vu Hung Nguyen, 2026}
\end{flushright}

\upsidehint{Swap $x$ and $y$ in the lower bound and compare the result with the upper bound.}

\end{exercisebox}
\subsubsection*{Problem 33. The Dyadic Parallelogram--Quadratic Squeeze}
\addcontentsline{toc}{subsubsection}{Problem 33. The Dyadic Parallelogram--Quadratic Squeeze}

\begin{exercisebox}[hard]

Find all functions $f$ that map real numbers to real numbers, satisfy
\[
f(f(1))\leq1,
\]
and, for all $x,y\in\mathbb R$,
\[
8xy+3f(x)-3f(y)
\leq\frac{5\bigl(f(x+y)+f(x-y)\bigr)}2
\leq5f(x)+5y^2.
\]

\begin{flushright}
\textit{Composed by Vu Hung Nguyen, 2026}
\end{flushright}

\upsidehint{First set one variable to zero; can you extract a simpler relation before using the iterate bound?}

\end{exercisebox}
\subsubsection*{Problem 39. The Twin Squeeze at Zero}
\addcontentsline{toc}{subsubsection}{Problem 39. The Twin Squeeze at Zero}

\begin{exercisebox}[medium]

Determine all functions $f:\mathbb R\to\mathbb R$ such that for all real $x,y$,
\[
\boxed{\quad f(x)f(y)\le f(xy)+x+y \qquad\text{and}\qquad f(x+y)\le f(x)+f(y)-1\quad}
\]

\begin{flushright}
\textit{Proposed by Vu Hung Nguyen, 2026}
\end{flushright}

\upsidehint{Plug $x=y=0$ into both inequalities to pin down $f(0)$, then use $y=-x$.}

\end{exercisebox}


\newpage
\section*{Solutions}
\addcontentsline{toc}{section}{Solutions}

\subsection*{Part I. Additive--Multiplicative Rigidity}
\addcontentsline{toc}{subsection}{Part I. Additive--Multiplicative Rigidity}

\addcontentsline{toc}{subsubsection}{Solution 1. The Additive--Multiplicative Rigidity}

\begin{dlsummary}[title={\bfseries Solution 1. The Additive--Multiplicative Rigidity},breakable]

\begin{proof}
Set $g(x)=\ln f(x)$. Since $\ln$ is strictly increasing, the given
inequalities are equivalent to
\[
g(x+y)\geq g(x)+g(y),
\qquad
g(xy)\leq g(x)+g(y). \tag{1}
\]

\medskip
\noindent
\textbf{Step 1. Show that $g$ vanishes on the positive integers.}

Taking $x=y=1$ in the second inequality of (1) gives $g(1)\geq0$.
Repeated superadditivity then gives
\[
g(n)\geq n g(1)\geq0\qquad(n\in\mathbb N).
\]
Fix an integer $n\geq2$. For every positive integer $m$, repeated
application of the two inequalities in (1) gives
\[
n^{m-1}g(n)\leq g(n^m)\leq m g(n).
\]
If $g(n)>0$, division by $g(n)$ would imply
$n^{m-1}\leq m$ for every positive integer $m$, which is impossible
for sufficiently large $m$. Since $g(n)\geq0$, it follows that
$g(n)=0$ for each integer $n\geq2$.

Superadditivity gives $g(2)\geq2g(1)$. As $g(2)=0$ and $g(1)\geq0$,
we obtain $g(1)=0$. Thus
\[
g(n)=0\qquad(n\in\mathbb N). \tag{2}
\]

\medskip
\noindent
\textbf{Step 2. Show that $g$ vanishes on positive rationals.}

Write a positive rational as $r=p/q$, where $p,q$ are positive
integers. By superadditivity and (2),
\[
0=g(p)=g(\underbrace{r+\cdots+r}_{q\text{ terms}})
\geq qg(r),
\]
so $g(r)\leq0$. By the second inequality in (1) and (2),
\[
0=g(p)=g(rq)\leq g(r)+g(q)=g(r),
\]
so $g(r)\geq0$. Therefore
\[
g(r)=0\qquad(r\in\mathbb Q_{>0}). \tag{3}
\]

\medskip
\noindent
\textbf{Step 3. Show that $g$ vanishes everywhere.}

Fix $x>0$ and an integer $n\geq2$. Using $g(n)=0$ in (1),
\[
g(x)=g\!\left(\frac xn\cdot n\right)\leq g\!\left(\frac xn\right),
\]
while superadditivity gives
\[
g(x)\geq n g\!\left(\frac xn\right).
\]
Hence
\[
g(x)\leq g\!\left(\frac xn\right)\leq\frac{g(x)}n.
\]
If $g(x)>0$, these bounds contradict $n\geq2$. Thus $g(x)\leq0$
for every $x>0$.

Finally, by (1) and (2),
\[
0=g(1)=g\!\left(x\cdot\frac1x\right)
\leq g(x)+g\!\left(\frac1x\right).
\]
Both terms on the right are nonpositive, so each must be zero. In
particular, $g(x)=0$ for all $x>0$. Therefore $f(x)=e^{g(x)}=1$.

\medskip
\noindent
\textbf{Verification.}

The constant function $f(x)=1$ makes both inequalities equalities,
so it is indeed the unique solution.
\end{proof}

\end{dlsummary}
\addcontentsline{toc}{subsubsection}{Solution 2. IMO 2013 -- Problem 5}

\begin{dlsummary}[title={\bfseries Solution 2. IMO 2013 -- Problem 5},breakable]

\begin{proof}
From $f(a)f(1)\geq f(a)$, we have $f(1)\geq1$, hence
$f(n)\geq n$ for positive integers $n$. Also,
$f(q)f(p/q)\geq f(p)>0$ shows $f(r)>0$ for every positive rational
$r$. Thus superadditivity makes $f$ strictly increasing.

Submultiplicativity gives $f(a^n)\leq a^n$, while
$f(a^n)f(a^{-n})\geq f(1)\geq1$ yields $f(a^{-n})\geq a^{-n}$.
For rational $x>0$, set $m_n=\lfloor a^n x\rfloor$. For large $n$,
monotonicity and superadditivity give
\[
f(x)\geq f(m_na^{-n})\geq m_nf(a^{-n})
\geq\frac{\lfloor a^n x\rfloor}{a^n}\longrightarrow x.
\]
Hence $f(x)\geq x$.

Applying this to $a^n$ shows $f(a^n)=a^n$. For any positive
rational $x$, choose $n$ with $a^n>x$. Then
\[
a^n=f(a^n)\geq f(x)+f(a^n-x)\geq f(x)+a^n-x,
\]
so $f(x)\leq x$. Therefore $f(x)=x$ for all positive rationals;
the identity satisfies both inequalities and the fixed-point condition.
\end{proof}

\end{dlsummary}
\addcontentsline{toc}{subsubsection}{Solution 3. The Periodic Excess}

\begin{dlsummary}[title={\bfseries Solution 3. The Periodic Excess},breakable]

\begin{proof}
\textbf{Step 1. Establish monotonicity and a growth bound.}

Since $f$ is positive and superadditive, it is strictly increasing:
if $y>x>0$, then
\[
f(y)=f(x+(y-x))\geq f(x)+f(y-x)>f(x).
\]
Taking $x=y=1$ in the multiplicative inequality gives
$f(1)\leq f(1)^2$, so $f(1)\geq1$. Repeated superadditivity then
implies
\[
f(n)\geq n\qquad(n\in\mathbb N).
\]
Fix $t>1$ and a positive integer $n$. Repeated submultiplicativity
and monotonicity give
\[
f(t)^n\geq f(t^n)\geq f(\lfloor t^n\rfloor)
\geq\lfloor t^n\rfloor.
\]
Taking $n$th roots and letting $n\to\infty$ yields
\[
f(t)\geq t\qquad(t>1). \tag{1}
\]

\medskip
\noindent
\textbf{Step 2. Force an exact translation law.}

Superadditivity and the third condition imply
\[
f(x)+f(a)\leq f(x+a)\leq f(x)+a.
\]
Thus $f(a)\leq a$. Since $a>1$, (1) also gives $f(a)\geq a$.
Therefore
\[
f(a)=a,\qquad f(x+a)=f(x)+a\quad(x>0). \tag{2}
\]

\medskip
\noindent
\textbf{Step 3. Study the periodic excess.}

Define $h(x)=f(x)-x$. Equation (2) gives
\[
h(x+a)=h(x),
\]
so $h$ is periodic with period $a$ on the positive reals.
Superadditivity of $f$ gives
\[
h(x+y)\geq h(x)+h(y). \tag{3}
\]
For any $x>0$, choose a positive integer $k$ large enough that
$x+ka>1$. By periodicity and (1),
\[
h(x)=h(x+ka)=f(x+ka)-(x+ka)\geq0.
\]
Since $f$ is increasing and $f(a)=a$, for $0<x\leq a$ we have
\[
h(x)=f(x)-x\leq a-x<a.
\]
By periodicity, this yields
\[
0\leq h(x)\leq a\qquad(x>0). \tag{4}
\]

\medskip
\noindent
\textbf{Step 4. Show that the periodic excess vanishes.}

Repeated use of (3) gives $h(nx)\geq n h(x)$ for every positive
integer $n$. By (4),
\[
0\leq n h(x)\leq h(nx)\leq a.
\]
Hence $0\leq h(x)\leq a/n$ for every positive integer $n$, forcing
$h(x)=0$. Thus $f(x)=x$ for all $x>0$.

\medskip
\noindent
\textbf{Verification.}

The identity function satisfies all three inequalities with equality,
so it is the unique solution.
\end{proof}

\end{dlsummary}
\addcontentsline{toc}{subsubsection}{Solution 4. The Dyadic Lock}

\begin{dlsummary}[title={\bfseries Solution 4. The Dyadic Lock},breakable]

\begin{proof}
The second inequality in the problem implies
\[
f(xy)\leq f(x)f(y). \tag{1}
\]
Also, $f$ is strictly increasing: if $u>v>0$, then
$f(u)=f(v+(u-v))\geq f(v)+f(u-v)>f(v)$.
Taking $x=y=1$ in (1) gives $f(1)\geq1$. Superadditivity then gives
$f(2)\geq2f(1)\geq2$. Since $f$ is increasing,
$f(f(2))\geq f(2)$, and therefore
\[
4\geq f(2)f(f(2))\geq f(2)^2\geq4.
\]
It follows that $f(2)=2$ and, from $f(2)\geq2f(1)\geq2$, also
$f(1)=1$. Superadditivity and (1), applied repeatedly, now give
\[
f(2^n)=2^n\qquad(n=0,1,2,\ldots). \tag{2}
\]

We next obtain a common upper bound. Fix $t>0$, and for all sufficiently
large $n$ let
\[
m_n=\left\lfloor\frac{2^n}{t}\right\rfloor.
\]
By monotonicity, superadditivity and (2),
\[
2^n=f(2^n)\geq f(m_nt)\geq m_nf(t).
\]
Since $2^n/m_n\to t$, this proves
\[
f(t)\leq t\qquad(t>0). \tag{3}
\]

For every positive integer $k$, superadditivity and (3) give
\[
k=kf(1)\leq f(k)\leq k,
\]
so $f(k)=k$. If $r=p/q$ is a positive rational, where $p,q$ are
positive integers, then (1) yields
\[
p=f(p)=f(qr)\leq f(q)f(r)=qf(r).
\]
Hence $f(r)\geq r$, and (3) gives $f(r)=r$. For any positive real $x$,
choose positive rationals $r<x<s$ arbitrarily close to $x$. Monotonicity
then gives $r=f(r)\leq f(x)\leq f(s)=s$. Letting $r,s$ approach $x$
proves $f(x)=x$.

The identity function satisfies the original inequalities and numerical
condition with equality, so it is the unique solution.
\end{proof}

\end{dlsummary}
\addcontentsline{toc}{subsubsection}{Solution 5. The Triadic Balance}

\begin{dlsummary}[title={\bfseries Solution 5. The Triadic Balance},breakable]

\begin{proof}
\textbf{Step 1. Unlock the value at $3$.}

Taking $x=y=1$ in submultiplicativity gives
\[
f(1)\le f(1)^2,
\]
so $f(1)\ge1$. By superadditivity,
\[
f(3)\ge3f(1)\ge3.
\]
Also, submultiplicativity gives
\[
f(1)=f\left(3\cdot\frac13\right)
\le f(3)f\left(\frac13\right),
\]
and hence
\[
f\left(\frac13\right)\ge\frac{f(1)}{f(3)}\ge\frac1{f(3)}.
\]
Writing $a=f(3)\ge3$, we obtain
\[
\frac{10}{3}
\ge f(3)+f\left(\frac13\right)
\ge a+\frac1a
\ge 3+\frac13=\frac{10}{3}.
\]
Equality holds throughout, so
\[
f(1)=1,\qquad f(3)=3,\qquad f\left(\frac13\right)=\frac13.
\]

\medskip
\noindent\textbf{Step 2. Establish an upper bound.}

By induction, superadditivity and submultiplicativity give
\[
f(3^n)=3^n\qquad(n=0,1,2,\ldots).
\]
The function is strictly increasing: if $u>v>0$, then
\[
f(u)=f\bigl(v+(u-v)\bigr)
\ge f(v)+f(u-v)>f(v).
\]

Fix $t>0$, and for all sufficiently large $n$ let
\[
m_n=\left\lfloor\frac{3^n}{t}\right\rfloor.
\]
Since $m_nt\le3^n$, monotonicity and superadditivity yield
\[
3^n=f(3^n)\ge f(m_nt)\ge m_nf(t).
\]
As $3^n/m_n\to t$, it follows that
\[
f(t)\le t\qquad(t>0). \tag{1}
\]

\medskip
\noindent\textbf{Method I. Rational squeeze.}

For every positive integer $k$, superadditivity and (1) give
\[
k=kf(1)\le f(k)\le k,
\]
so $f(k)=k$. If $r=p/q$ is a positive rational, where $p,q$ are
positive integers, submultiplicativity gives
\[
p=f(p)=f(qr)\le f(q)f(r)=qf(r).
\]
Thus $f(r)\ge r$, and (1) shows $f(r)=r$. By monotonicity and the
density of the positive rationals, $f(x)=x$ for all $x>0$.

\end{proof}

The identity function satisfies the original conditions with equality
throughout, so it is the unique solution.

\end{dlsummary}
\addcontentsline{toc}{subsubsection}{Solution 34. The Ratio--Product Squeeze}

\begin{dlsummary}[title={\bfseries Solution 34. The Ratio--Product Squeeze},breakable]

\begin{proof}
Define
\[
g(x)=\frac{f(x)}x\qquad(x>0).
\]
The first inequality implies
\[
g(x+y)\ge g(x),
\]
so $g$ is nondecreasing. The second inequality becomes
\[
g(xy)\le g(x)g(y). \tag{1}
\]
Taking $x=y=1$ in (1) gives $g(1)\le g(1)^2$. Since $g(1)>0$,
\[
g(1)\ge1. \tag{2}
\]

Let $a>1$ satisfy the composition condition. By monotonicity and (2),
\[
g(a)\ge g(1)\ge1,
\qquad
f(a)=ag(a)\ge a>1.
\]
Therefore $g(f(a))\ge g(a)\ge1$. On the other hand,
\[
f(f(a))=f(a)g(f(a))\le f(a),
\]
so $g(f(a))\le1$. It follows that
\[
g(f(a))=g(a)=g(1)=1.
\]
By monotonicity,
\[
g(x)=1\qquad(1\le x\le a). \tag{3}
\]

Now let $1/a\le x\le1$. From monotonicity and (3), $g(x)\le1$ and
$g(1/x)=1$. Applying (1) to $x$ and $1/x$ gives
\[
1=g(1)\le g(x)g(1/x)=g(x),
\]
so $g(x)=1$. Hence
\[
g(x)=1\qquad(1/a\le x\le a). \tag{4}
\]

Finally, fix any $x>0$. Choose a positive integer $n$ large enough that
$t=x^{1/n}\in[1/a,a]$. By (4), $g(t)=g(1/t)=1$. Repeatedly applying
(1) yields
\[
g(x)=g(t^n)\le g(t)^n=1,
\qquad
g(1/x)=g((1/t)^n)\le g(1/t)^n=1.
\]
But (1) also gives
\[
1=g(1)\le g(x)g(1/x).
\]
Both factors are positive and at most $1$, so they are both equal to $1$.
In particular, $g(x)=1$. Thus $g\equiv1$, and consequently
\[
f(x)=x\qquad(x>0).
\]
\end{proof}

The identity function satisfies both inequalities with equality and
$f(f(a))=f(a)$, so it is the unique solution.

\end{dlsummary}
\addcontentsline{toc}{subsubsection}{Solution 35. The Subadditive--Supermultiplicative Squeeze}

\begin{dlsummary}[title={\bfseries Solution 35. The Subadditive--Supermultiplicative Squeeze},breakable]

\begin{proof}
The first inequality shows that $f$ is nondecreasing, and the second
shows that $f$ is subadditive. Since $f(x)+f(y)>0$, the last inequality
is equivalent to
\[
f(xy)\geq f(x)f(y). \tag{1}
\]
Thus $f$ is also supermultiplicative.

\medskip
\noindent
\textbf{Step 1. The anchor is a fixed point.}

Let $b=f(a)$. If $b>a$, monotonicity gives
\[
a=f(b)\geq f(a)=b>a,
\]
a contradiction. If $b<a$, then
\[
a=f(b)\leq f(a)=b<a,
\]
again a contradiction. Hence
\[
f(a)=a. \tag{2}
\]

\medskip
\noindent
\textbf{Step 2. Determine $f(1)$.}

Set $c=f(1)$. By (1), $c\geq c^2$, so $c\leq1$. For every positive
integer $n$, repeated use of (1) and (2) gives
\[
f(a^n)\geq f(a)^n=a^n.
\]
On the other hand, monotonicity and subadditivity give
\[
f(a^n)\leq f(\lceil a^n\rceil)
\leq \lceil a^n\rceil c.
\]
Therefore
\[
c\geq\frac{a^n}{\lceil a^n\rceil}\longrightarrow1,
\]
and so
\[
f(1)=1. \tag{3}
\]

\medskip
\noindent
\textbf{Step 3. Show that $f(x)\geq x$.}

Fix $x>0$. For all sufficiently large $n$, set
\[
m_n=\left\lfloor\frac{a^n}{x}\right\rfloor,
\qquad r_n=a^n-m_nx,
\]
so $m_n\geq1$ and $0\leq r_n<x$. By (1) and (2),
\[
a^n\leq f(a^n).
\]
If $r_n>0$, subadditivity, repeated subadditivity, and monotonicity give
\[
f(a^n)=f(m_nx+r_n)
\leq f(m_nx)+f(r_n)
\leq m_nf(x)+f(x).
\]
If $r_n=0$, the same final bound follows directly from
$f(m_nx)\leq m_nf(x)$. Thus
\[
f(x)\geq\frac{a^n}{m_n+1}
=\frac{a^n}{\left\lfloor a^n/x\right\rfloor+1}.
\]
Letting $n\to\infty$ yields
\[
f(x)\geq x. \tag{4}
\]

\medskip
\noindent
\textbf{Step 4. Force equality.}

Using (1), (3), and (4) applied to $1/x$, we obtain
\[
1=f(1)=f\!\left(x\cdot\frac1x\right)
\geq f(x)f\!\left(\frac1x\right)
\geq\frac{f(x)}x.
\]
Hence $f(x)\leq x$. Together with (4), this proves
\[
f(x)=x\qquad(x>0).
\]

The identity function satisfies the chain and $f(f(a))=a$ for every
$a>1$, so it is the unique solution.
\end{proof}

\end{dlsummary}
\addcontentsline{toc}{subsubsection}{Solution 38. The Injective Staircase}

\begin{dlsummary}[title={\bfseries Solution 38. The Injective Staircase},breakable]

\begin{proof}
The only solution is \(f(n)=n\). It satisfies the inequality with equality: \(m+n\le n+(m+1)-1\).

\paragraph{1. Show \(f(1)=1\).} Put \(n=1\) and \(c=f(1)\). Then
\[
f(m+c)\le f(m+1)\quad\text{for all } m\ge1.
\]
Suppose \(c\ge2\). Fix \(j\ge2\) and apply the bound repeatedly. This gives
\[
f\bigl(j+t(c-1)\bigr)\le f(j)\quad\text{for all } t\ge0.
\]
So infinitely many distinct inputs would map into the finite set \(\{1,\dots,f(j)\}\). This contradicts injectivity, so \(c=1\).

\paragraph{2. Show \(f(2)=2\).} Let \(k=f(2)\). Since \(f(1)=1\) and \(f\) is injective, \(k\ge2\). Put \(n=2\):
\[
f(m+k)\le f(m+1)+1.
\]
Writing \(j=m+1\ge2\), this reads
\[
f\bigl(j+(k-1)\bigr)\le f(j)+1.
\]
Suppose \(k\ge3\). Every \(n\ge2\) can be written as \(n=j+t(k-1)\) with \(j\in\{2,\dots,k\}\) and \(t\le\frac{n}{k-1}\). Let \(M=\max_{2\le j\le k}f(j)\). Then
\[
f(n)\le M+\frac{n}{k-1}\le M+\frac n2 .
\]
Now consider \(f(1),\dots,f(N)\). These are \(N\) distinct positive integers, all at most \(M+\frac N2\). Hence \(N\le M+\frac N2\), which fails for large \(N\). So \(k=2\).

\paragraph{3. Conclude \(f(n)=n\).} With \(k=2\), the bound from Step 2 becomes \(f(j+1)\le f(j)+1\). Since \(f(1)=1\), induction gives \(f(n)\le n\) for all \(n\). An injective map of \(\{1,\dots,N\}\) into itself is a bijection. So \(f\) permutes \(\{1,\dots,N\}\) for every \(N\), and therefore \(f(n)=n\). \(\blacksquare\)
\end{proof}
\end{dlsummary}

\begin{takeaways}
\begin{itemize}
\item \textbf{Use \(n=1\) and \(n=2\) as probes:} Small values of \(n\) turn a two-variable inequality into one-step shift bounds.
\item \textbf{Injectivity beats descent:} A rule like \(f(x+d)\le f(x)\) cannot continue forever on an injective map into the positive integers.
\item \textbf{Counting is a quantitative injectivity argument:} Compare \(N\) distinct values with a linear upper bound.
\item \textbf{Show the hypothesis is needed:} The constant function \(f\equiv1\) satisfies the inequality, so injectivity cannot be dropped.
\end{itemize}
\end{takeaways}

\addcontentsline{toc}{subsubsection}{Solution 41. The Prime-Exponent Trap}

\begin{dlsummary}[title={\bfseries Solution 41. The Prime-Exponent Trap},breakable]

\begin{proof}
We prove that the only solution is
\[
\boxed{f(x)=x\qquad(x\in\mathbb Q_{>0}).}
\]

\paragraph{1. Turn the inequality into equality.}

Positivity and multiplicativity give
\[
f(1)=1,\qquad f(1/x)=\frac1{f(x)}.
\]
Applying the inequality at \(1/x\), we obtain
\[
\frac1{f(f(x)^3)}\ge\frac1{x^2f(x)}.
\]
Together with the original inequality, this forces
\[
f(f(x))^3=x^2f(x). \tag{1}
\]

\paragraph{2. Isolate the deviation from the identity.}

Fix \(x>0\), and put
\[
h=\frac{f(x)}x.
\]
By multiplicativity,
\[
f(f(x))=f(xh)=f(x)f(h).
\]
Substituting into (1) gives
\[
f(h)^3=h^{-2}. \tag{2}
\]

We will prove that a positive rational \(h\) satisfying (2) must equal \(1\).

\paragraph{3. Iterate and use prime exponents.}

Set \(h_0=h\) and \(h_{n+1}=f(h_n)\). Applying \(f\) repeatedly to (2), using multiplicativity, gives
\[
h_{n+1}^3=h_n^{-2}\qquad(n\ge0).
\]

For a prime \(p\), let \(e_n\) be the exponent of \(p\) in the prime factorisation of the positive rational \(h_n\). Each \(e_n\) is an integer, possibly negative, and
\[
3e_{n+1}=-2e_n.
\]
Consequently,
\[
3^n e_n=(-2)^n e_0.
\]
Since \(2\) and \(3\) are coprime, \(3^n\) divides \(e_0\) for every \(n\). Therefore \(e_0=0\).

This holds for every prime \(p\), so \(h=1\). Hence \(f(x)=x\) for every positive rational \(x\).

Conversely, \(f(x)=x\) satisfies both conditions, with equality.
\end{proof}

\paragraph{Takeaways.}
\begin{itemize}
\item **Reciprocals reverse inequalities:** Multiplicativity makes the substitution \(x\mapsto1/x\) especially powerful.
\item **Normalise the defect:** The ratio \(f(x)/x\) separates the desired solution from its possible deviations.
\item **Iteration creates divisibility:** A relation that looks harmless once can become impossible when repeated indefinitely.
\item **The rational domain matters:** Prime exponents must remain integers at every iteration.
\end{itemize}

\end{dlsummary}

\addcontentsline{toc}{subsubsection}{Solution 42. Bounded Errors, Exact Answer}

\begin{dlsummary}[title={\bfseries Solution 42. Bounded Errors, Exact Answer},breakable]

\begin{proof}
The only solution is
\[
\boxed{f(x)=x\qquad(x\in\mathbb R).}
\]

\paragraph{1. Force unboundedness.}

Repeatedly applying the additive inequality gives
\[
f(2n)\ge nf(2)-(n-1)
=1+n\bigl(f(2)-1\bigr).
\]
Since \(f(2)>1\), these values tend to infinity.

\paragraph{2. Eliminate the multiplicative error.}

Write
\[
E(x,y)=f(xy)-f(x)f(y),\qquad |E(x,y)|\le1.
\]
Evaluating the same triple product in two ways yields
\[
E(x,y)f(t)
=E(x,yt)-E(xy,t)+f(x)E(y,t).
\]
Consequently,
\[
|E(x,y)|\,|f(t)|\le2+|f(x)|.
\]
For fixed \(x,y\), the right-hand side is constant, while \(|f(t)|\) can be arbitrarily large. Therefore \(E(x,y)=0\), proving
\[
f(xy)=f(x)f(y). \tag{1}
\]

\paragraph{3. Eliminate the additive error.}

Apply the additive inequality to \(tx,ty\). Using (1),
\[
|f(t)|\,|f(x+y)-f(x)-f(y)|\le1.
\]
Again, unboundedness forces
\[
f(x+y)=f(x)+f(y). \tag{2}
\]

\paragraph{4. Recover order, then the function.}

Multiplicativity gives \(f(1)^2=f(1)\); since \(f\) is not identically zero, \(f(1)=1\). Additivity now gives
\[
f(q)=q\qquad(q\in\mathbb Q).
\]

For \(u>0\),
\[
f(u)=f(\sqrt u)^2>0,
\]
where strict positivity follows from
\[
f(\sqrt u)\,f(1/\sqrt u)=1.
\]
Thus, if \(b>a\), then
\[
f(b)-f(a)=f(b-a)>0.
\]
So \(f\) is strictly increasing. Comparing any real \(x\) with rational numbers immediately below and above it forces \(f(x)=x\).

The identity function plainly satisfies both inequalities and \(f(2)>1\).
\end{proof}
\end{dlsummary}

\begin{takeaways}
\begin{itemize}
\item \textbf{Associativity can remove errors:} Comparing two expansions of a triple product makes a fixed discrepancy multiply an unbounded quantity.
\item \textbf{Order can be a conclusion:} Squares establish positivity, and additivity turns positivity into monotonicity.
\item \textbf{A threshold can matter:} The condition \(f(2)>1\) excludes \(f\equiv1\), which satisfies both error bounds.
\item \textbf{Avoid assuming regularity:} The argument establishes every property it needs directly from the hypotheses.
\end{itemize}
\end{takeaways}

\addcontentsline{toc}{subsubsection}{Solution 46. The Missing Middle Value}

\begin{dlsummary}[title={\bfseries Solution 46. The Missing Middle Value},breakable]

\begin{proof}
The only solutions are
\[
\boxed{f(n)=n\qquad\text{and}\qquad f(n)=n^2.}
\]

\paragraph{1. Establish order and bound \(f(2)\).}

Multiplicativity and positivity give \(f(1)=1\). Superadditivity then implies
\[
f(n+1)\ge f(n)+1.
\]
Thus \(f\) is strictly increasing and \(f(n)\ge n\).

Write \(k=f(2)\), so \(k\ge2\). If \(k\ge5\), strict increase and multiplicativity give
\[
f(f(2))=f(k)>f(4)=k^2\ge25,
\]
contradicting \(f(f(2))\le16\). Therefore
\[
\boxed{k\in\{2,3,4\}.}
\]

\paragraph{2. Eliminate the middle value.}

Suppose \(f(2)=3\). Since \(9>8\),
\[
f(3)^2=f(9)>f(8)=27.
\]
Because \(f(3)\) is an integer, this forces \(f(3)\ge6\).

But \(243<256\), so
\[
f(3)^5=f(243)<f(256)=3^8.
\]
This is impossible, since
\[
f(3)^5\ge6^5=7776>6561=3^8.
\]
Hence
\[
f(2)=2\quad\text{or}\quad f(2)=4.
\]

\paragraph{3. One value determines the whole function.}

Write \(f(2)=2^d\), where \(d\in\{1,2\}\). Fix \(n\ge2\).

For each positive integer \(r\), choose an integer \(s\ge0\) such that
\[
2^s\le n^r<2^{s+1}.
\]
Monotonicity and multiplicativity give
\[
2^{ds}\le f(n)^r<2^{d(s+1)}.
\]
Also,
\[
2^{ds}\le n^{dr}<2^{d(s+1)}.
\]
Dividing the bounds appropriately,
\[
2^{-d}<\left(\frac{f(n)}{n^d}\right)^r<2^d.
\]
These bounds hold for every positive integer \(r\). A positive number other than \(1\) has powers tending either to \(0\) or to infinity, so
\[
f(n)=n^d.
\]
Together with \(f(1)=1\), this proves the classification.

\paragraph{4. Verify both solutions.}

For \(f(n)=n\), all conditions are immediate. For \(f(n)=n^2\),
\[
(m+n)^2\ge m^2+n^2,
\qquad
f(f(n))=n^4,
\]
and multiplicativity also holds.
\end{proof}

\paragraph{Takeaways.}
\begin{itemize}
\item **Iteration can bound a small value:** The condition at \(n=2\) sharply restricts \(f(2)\).
\item **Integrality creates gaps:** The inequality \(f(3)^2>27\) becomes the stronger statement \(f(3)\ge6\).
\item **Compare the same number in different ways:** Multiplicativity translates numerical power comparisons into restrictions on function values.
\item **Use powers to amplify discrepancies:** A fixed multiplicative error cannot remain bounded through arbitrarily high powers.
\end{itemize}

\end{dlsummary}

\addcontentsline{toc}{subsubsection}{Solution 49. The Prime-Pairing Machine}

\begin{dlsummary}[title={\bfseries Solution 49. The Prime-Pairing Machine},breakable]

\begin{proof}
\textbf{The solutions:} Partition the set of all primes into pairs. Within each pair, designate one prime \(p\) and the other \(q\), and set
\[
\boxed{f(p)=q,\qquad f(q)=p^2.}
\]
Extend these assignments multiplicatively to all positive rationals. Every solution arises this way.

\paragraph{1. Remove the inequality.}

Positivity and multiplicativity give
\[
f(1)=1,\qquad f(1/x)=\frac1{f(x)}.
\]
Applying the composition inequality to \(1/x\),
\[
\frac1{f(f(x))}\ge\frac1{x^2}.
\]
Together with the original inequality, this proves
\[
\boxed{f(f(x))=x^2\qquad(x>0).} \tag{1}
\]

\paragraph{2. Show that primes map to prime powers.}

Let \(p\) be prime. Since \(f(p)\) is an integer and
\[
f(f(p))=p^2>1,
\]
we have \(f(p)>1\).

Choose a prime divisor \(q\) of \(f(p)\). Factoring \(f(p)\) into primes and applying \(f\), all factors are positive integers and their product is \(p^2\). Also, \(f(q)>1\), by the same reasoning as above. Thus
\[
f(q)=p^b
\]
for some positive integer \(b\).

Apply \(f\) once more:
\[
q^2=f(f(q))=f(p)^b.
\]
Unique prime factorisation forces
\[
\boxed{f(p)=q^a}
\]
for a positive integer \(a\). In particular, \(f(p)\) has only one prime divisor.

\paragraph{3. Identify the allowed pairs.}

Now
\[
p^2=f(f(p))=f(q)^a=p^{ab},
\]
so
\[
ab=2.
\]
Therefore
\[
(a,b)=(1,2)\quad\text{or}\quad(2,1).
\]

Moreover, \(q\ne p\): otherwise \(f(p)=p^a\) would give \(a^2=2\), impossible for an integer \(a\).

Thus every prime belongs to a two-prime pair, with assignments
\[
f(p)=q,\qquad f(q)=p^2
\]
after choosing the orientation. These pairs are disjoint because the prime associated with any given prime is uniquely determined.

\paragraph{4. Construct and verify every solution.}

Choose any partition of the primes into oriented pairs, make the assignments above, and define
\[
f\!\left(\prod_r r^{e_r}\right)=\prod_r f(r)^{e_r},
\]
where the exponents \(e_r\) are integers and only finitely many are nonzero.

This is well-defined and multiplicative. It sends positive integers to positive integers, and on every prime,
\[
f(f(r))=r^2.
\]
Multiplicativity therefore gives \(f(f(x))=x^2\) for every positive rational \(x\), verifying all conditions.

\paragraph{Example.}

Pair consecutive primes and orient them as
\[
(2,3),\ (5,7),\ (11,13),\ldots.
\]
Then
\[
f(2)=3,\quad f(3)=4,\quad f(5)=7,\quad f(7)=25,
\]
so, for instance,
\[
f(6)=12,\qquad f(12)=36.
\]
This construction also passed direct composition checks on the integers from 1 through 100.[1]
\end{proof}
\end{dlsummary}

\begin{takeaways}
\begin{itemize}
\item \textbf{Reciprocals can force exactness:} A global multiplicative inequality may secretly be an equation.
\item \textbf{Integrality controls factorisation:} Integer-valued images prevent unwanted prime factors from cancelling.
\item \textbf{Follow a prime through two iterations:} The return to \(p^2\) restricts both the prime support and the exponents.
\item \textbf{Classification need not mean uniqueness:} A complete answer can be a structural recipe for infinitely many functions.
\end{itemize}
\end{takeaways}

\addcontentsline{toc}{subsubsection}{Solution 53. When Even Numbers Become Smaller}

\begin{dlsummary}[title={\bfseries Solution 53. When Even Numbers Become Smaller},breakable]

\begin{proof}
The unique solution is
\[
\boxed{
f(0)=0,\qquad
f\!\left(2^k\frac ab\right)=2^{-k},
}
\]
where \(k\in\mathbb Z\), and \(a,b\) are odd integers with \(b>0\).

In other words, only the exponent of \(2\) in a nonzero rational number matters.

\paragraph{1. Establish the basic identities.}

Since
\[
f(2)=f(2)f(1)>0,
\]
we have \(f(1)=1\). Consequently, for \(x\ne0\),
\[
f(x)f(1/x)=1,
\]
so \(f(x)>0\) and
\[
f(1/x)=\frac1{f(x)}.
\]
Also, \(f(-1)^2=f(1)=1\), so nonnegativity gives \(f(-1)=1\). Thus
\[
f(-x)=f(x).
\]

\paragraph{2. Use binary expansions to bound all integers.}

Write a positive integer \(n\) in binary:
\[
n=\varepsilon_0+\varepsilon_1 2+\cdots+\varepsilon_s2^s,
\qquad \varepsilon_j\in\{0,1\}.
\]
Repeated subadditivity gives
\[
f(n)\le\sum_{j=0}^s\varepsilon_j f(2^j)
=\sum_{j=0}^s\varepsilon_j2^{-j}<2.
\]
Therefore
\[
\boxed{f(n)<2\qquad(n\in\mathbb Z_{>0}).} \tag{1}
\]

\paragraph{3. Sharpen the bound by taking powers.}

Apply (1) to \(n^r\). Multiplicativity yields
\[
f(n)^r=f(n^r)<2
\]
for every positive integer \(r\).

If \(f(n)>1\), its powers would eventually exceed \(2\). Hence
\[
\boxed{f(n)\le1\qquad(n\in\mathbb Z_{>0}).} \tag{2}
\]

\paragraph{4. Determine the values at odd integers.}

Let \(n\) be a positive odd integer. For every positive integer \(r\),
\[
j=\frac{n^r-1}{2}
\]
is a nonnegative integer. Since \(1=n^r-2j\), subadditivity, evenness, and (2) imply
\[
1=f(1)\le f(n^r)+f(2j)
=f(n)^r+\frac12f(j)
\le f(n)^r+\frac12.
\]
Thus
\[
f(n)^r\ge\frac12
\]
for every \(r\). If \(f(n)<1\), its powers would tend to zero. Together with (2), this proves
\[
\boxed{f(n)=1\qquad(n\text{ odd}).}
\]

\paragraph{5. Determine all rational values and verify.}

Every nonzero rational has a unique representation
\[
x=2^k\frac ab
\]
with \(a,b\) odd and \(b>0\). Multiplicativity now forces
\[
f(x)=f(2)^k\frac{f(a)}{f(b)}=2^{-k}.
\]

Conversely, define \(f\) by this formula. Multiplication adds the exponents of \(2\), so multiplicativity holds.

For nonzero \(x,y\), write
\[
x=2^k u,\qquad y=2^\ell v,\qquad k\le\ell,
\]
where \(u,v\) have odd numerators and denominators. Then
\[
x+y=2^k\bigl(u+2^{\ell-k}v\bigr).
\]
Unless the sum is zero, the expression in parentheses has a nonnegative exponent of \(2\). Therefore
\[
f(x+y)\le2^{-k}=\max\{f(x),f(y)\}\le f(x)+f(y).
\]
Zero cases are immediate, completing the verification.
\end{proof}

\paragraph{Takeaways.}
\begin{itemize}
\item **A small value at \(2\) makes binary expansions powerful:** They produce a uniform bound independent of the integer.
\item **Powers sharpen uniform bounds:** From \(f(n^r)<2\), multiplicativity forces \(f(n)\le1\).
\item **Parity supplies the reverse bound:** The identity \(n^r-2j=1\) prevents odd integers from having values below \(1\).
\item **Stronger inequalities can emerge:** The solution satisfies the maximum bound \(f(x+y)\le\max\{f(x),f(y)\}\), although only a sum bound was assumed.
\end{itemize}

\end{dlsummary}

\addcontentsline{toc}{subsubsection}{Solution 58. The Four-Step Prime Labyrinth}

\begin{dlsummary}[title={\bfseries Solution 58. The Four-Step Prime Labyrinth},breakable]

\begin{proof}
Write \(f^{\circ j}\) for \(j\) successive applications of \(f\).

\paragraph{1. Force equality by taking reciprocals.}

Positivity and multiplicativity give
\[
f(1)=1,\qquad f(1/x)=\frac1{f(x)}.
\]
The same reciprocal identity holds for every iterate. Applying the given inequality at \(1/x\) therefore gives the reverse bound, so
\[
\boxed{f^{\circ4}(x)=x^4.} \tag{1}
\]

For every prime \(p\), we must have \(f(p)>1\); otherwise all subsequent iterates would equal \(1\), contradicting (1).

Also, multiplicativity and integer preservation imply
\[
a\mid b\quad\Longrightarrow\quad f(a)\mid f(b)
\]
for positive integers \(a,b\).

\paragraph{2. Prove that each prime maps to a prime power.}

Fix a prime \(p\). Choose primes successively so that
\[
q\mid f(p),\qquad r\mid f(q),\qquad s\mid f(r).
\]
These choices are possible because the image of every prime exceeds \(1\).

By repeated divisibility,
\[
f(s)\mid f^{\circ4}(p)=p^4.
\]
Since \(f(s)>1\), we have \(p\mid f(s)\). Thus we have a closed divisibility chain:
\[
p\longrightarrow q\longrightarrow r\longrightarrow s\longrightarrow p,
\]
where an arrow \(u\longrightarrow v\) means \(v\mid f(u)\).

Now let \(q'\) be \textbf{any} prime divisor of \(f(p)\). Restarting at \(q\), we obtain
\[
q\longrightarrow r\longrightarrow s\longrightarrow p\longrightarrow q'.
\]
Hence
\[
q'\mid f^{\circ4}(q)=q^4,
\]
so \(q'=q\).

Therefore \(f(p)\) has only one prime divisor:
\[
\boxed{f(p)=\sigma(p)^{a_p},}
\]
where \(\sigma(p)\) is prime and \(a_p\) is a positive integer.

\paragraph{3. Classify the prime cycles.}

Applying \(f\) four times,
\[
f^{\circ4}(p)
=
\sigma^4(p)^{
a_p a_{\sigma(p)}a_{\sigma^2(p)}a_{\sigma^3(p)}
}.
\]
Comparing with \(p^4\),
\[
\boxed{
\sigma^4(p)=p,\qquad
a_p a_{\sigma(p)}a_{\sigma^2(p)}a_{\sigma^3(p)}=4.
} \tag{2}
\]

Thus \(\sigma\) is a permutation whose cycles have lengths \(1,2,\) or \(4\).

\begin{itemize}
\item **Length \(1\) is impossible:** It would require \(a_p^4=4\).
\item **Length \(2\):** If \(p\leftrightarrow q\), then \((a_pa_q)^2=4\), so \(a_pa_q=2\).
\item **Length \(4\):** The four exponents must have product \(4\).
\end{itemize}

\paragraph{Complete classification.}

Partition the primes into disjoint directed cycles of the following types:

| Cycle | Assignments |
|---|---|
| Two primes \(p,q\) | \(f(p)=q,\quad f(q)=p^2\), after choosing an orientation |
| Four primes \(p_1,p_2,p_3,p_4\) | \(f(p_i)=p_{i+1}^{a_i}\), with indices taken cyclically and \(a_1a_2a_3a_4=4\) |

For a four-cycle, the exponent pattern consists either of \textbf{one \(4\) and three \(1\)s}, or \textbf{two \(2\)s and two \(1\)s}, in any cyclic arrangement.

Extend these prime assignments multiplicatively:
\[
f\!\left(\prod_p p^{e_p}\right)
=\prod_p f(p)^{e_p},
\qquad e_p\in\mathbb Z.
\]
These are exactly all solutions.

\textbf{Verification.} The construction preserves positive integers and is multiplicative. On a two-cycle, four iterations produce exponent \((1\cdot2)^2=4\); on a four-cycle, they produce exponent \(a_1a_2a_3a_4=4\). Therefore \(f^{\circ4}(x)=x^4\) for every positive rational \(x\).
\end{proof}

\paragraph{Takeaways.}
\emph{No takeaways were provided in the source draft.}

\end{dlsummary}

\addcontentsline{toc}{subsubsection}{Solution 59. The Hidden Product Rule}

\begin{dlsummary}[title={\bfseries Solution 59. The Hidden Product Rule},breakable]

\begin{proof}
The unique solution is
\[
\boxed{f(x)=x\qquad(x\in\mathbb R).}
\]

\paragraph{1. Remove the expected solution.}

Define
\[
g(x)=f(x)-x.
\]
The first two inequalities become
\[
g(x+y)\ge g(x)+g(y), \tag{1}
\]
\[
g(xy)\ge xg(y)+yg(x). \tag{2}
\]

We first show that both are actually equalities.

\paragraph{2. Establish the normalisation and oddness.}

Putting \(y=0\) in (2),
\[
g(0)\ge xg(0)\qquad(x\in\mathbb R),
\]
which forces \(g(0)=0\).

Putting \(x=1\) in (2),
\[
0\ge yg(1)\qquad(y\in\mathbb R),
\]
so \(g(1)=0\).

Taking \(x=y=-1\) in (2) gives
\[
0\ge-2g(-1),
\]
so \(g(-1)\ge0\). But (1) at \((1,-1)\) gives \(g(-1)\le0\). Hence \(g(-1)=0\).

Now (2) at \((-1,x)\) gives
\[
g(-x)\ge-g(x),
\]
while (1) at \((x,-x)\) gives the reverse inequality. Therefore
\[
\boxed{g(-x)=-g(x).} \tag{3}
\]

\paragraph{3. Force additivity and the product identity.}

Apply (1) to \(-x,-y\), then use oddness:
\[
g(x+y)\le g(x)+g(y).
\]
Thus
\[
\boxed{g(x+y)=g(x)+g(y).} \tag{4}
\]

Similarly, applying (2) to \(-x,y\) reverses its direction. Hence
\[
\boxed{g(xy)=xg(y)+yg(x).} \tag{5}
\]

\paragraph{4. Interpret the composition inequality.}

Since \(f(t)=t+g(t)\), additivity gives
\[
f(f(t))
=f(t)+g(f(t))
=t+2g(t)+g(g(t)).
\]
Consequently, the third original inequality becomes
\[
\boxed{g(g(t))\ge0\qquad(t\ge0).} \tag{6}
\]

Let \(h=g\circ g\). It is additive, and (6) makes it nondecreasing: if \(b>a\), then
\[
h(b)-h(a)=h(b-a)\ge0.
\]
Moreover,
\[
h(1)=g(g(1))=0.
\]
An additive, nondecreasing function with value \(0\) at \(1\) is identically zero: it vanishes on the rationals, and rational bounds squeeze its value at every real input. Thus
\[
\boxed{g(g(x))=0\qquad(x\in\mathbb R).} \tag{7}
\]

\paragraph{5. Apply the product identity twice.}

From (5),
\[
g(x^2)=2xg(x).
\]
Apply \(g\) again. By additivity and (5),
\[
\begin{aligned}
g(g(x^2))
&=2g(xg(x))\\
&=2xg(g(x))+2g(x)^2.
\end{aligned}
\]
Both composite terms vanish by (7), leaving
\[
0=2g(x)^2.
\]
Hence \(g(x)=0\) everywhere, so \(f(x)=x\).

The identity function satisfies all three inequalities with equality.
\end{proof}
\end{dlsummary}

\begin{takeaways}
\begin{itemize}
\item \textbf{Subtract the expected model:} This exposes the algebraic structure hidden in the inequalities.
\item \textbf{Negative inputs reverse information:} Oddness turns superadditivity and the product inequality into exact identities.
\item \textbf{Regularise the composite first:} Positivity on nonnegative inputs makes \(g\circ g\) monotone, even though no regularity of \(g\) is assumed.
\item \textbf{Apply a structural identity twice:} The final square \(g(x)^2\) eliminates every remaining possibility.
\end{itemize}
\end{takeaways}

\addcontentsline{toc}{subsubsection}{Solution 62. The Hidden Square Root}

\begin{dlsummary}[title={\bfseries Solution 62. The Hidden Square Root},breakable]

\begin{proof}
The unique solution is
\[
\boxed{f(x)=\sqrt{|x|}\qquad(x\in\mathbb Q).}
\]

\paragraph{1. Turn the multiplicative inequality into equality.}

For \(x\ne0\), supermultiplicativity gives
\[
1=f(1)\ge f(x)f(1/x).
\]
This repeats the stated bound, so we need a stronger argument.

Apply supermultiplicativity to \(4x\) and \(1/x\):
\[
2=f(4)\ge f(4x)f(1/x)\ge2f(x)f(1/x).
\]
Again, this only supplies the same upper bound.

In fact, these hypotheses do \textbf{not} force multiplicativity: the reciprocal condition points in the same direction as supermultiplicativity at a reciprocal pair.

\textbf{The corrected problem replaces the reciprocal condition by}
\[
\boxed{\qquad f(x)f(1/x)\ge1\qquad(x\ne0).\qquad} \tag{3}
\]

With this correction, the two bounds force
\[
f(x)f(1/x)=1.
\]
In particular, \(f(x)>0\) for \(x\ne0\).

Now, for nonzero \(x,y\),
\[
f(x)=f((xy)(1/y))
\ge f(xy)f(1/y)=\frac{f(xy)}{f(y)}.
\]
Combined with supermultiplicativity, this proves
\[
\boxed{f(xy)=f(x)f(y).} \tag{4}
\]

Consequently,
\[
f(2)=\sqrt2,\qquad f(-1)=1,\qquad f(-x)=f(x).
\]

\paragraph{2. Prove an upper bound using binary expansions.}

Write a positive integer \(N\) in binary:
\[
N=\sum_{j=0}^{s}\varepsilon_j2^j,
\qquad \varepsilon_j\in\{0,1\},\quad 2^s\le N.
\]
Subadditivity and multiplicativity give
\[
f(N)\le\sum_{j=0}^{s}(\sqrt2)^j
< C\sqrt N,
\qquad
C=\frac{\sqrt2}{\sqrt2-1}.
\]
Apply this to \(N=n^r\):
\[
f(n)^r=f(n^r)<C\,n^{r/2}.
\]
Taking \(r\)-th roots and letting \(r\to\infty\),
\[
\boxed{f(n)\le\sqrt n.} \tag{5}
\]

\paragraph{3. Prove the reverse bound using base-\(b\) expansions.}

Fix an integer \(b\ge2\), and write \(A=f(b)>0\). Let
\[
M=\max\{f(0),f(1),\ldots,f(b-1)\}.
\]

Write \(2^k\) in base \(b\):
\[
2^k=d_0+d_1b+\cdots+d_sb^s,
\qquad 0\le d_j<b,\quad b^s\le2^k.
\]
Then
\[
2^{k/2}=f(2^k)\le M\sum_{j=0}^{s}A^j. \tag{6}
\]

If \(A\le1\), the right side is at most \(M(s+1)\), which grows at most linearly in \(k\). The left side grows exponentially—a contradiction. Thus \(A>1\).

The geometric sum in (6) now gives
\[
2^{k/2}
\le \frac{MA}{A-1}A^s
\le \frac{MA}{A-1}A^{k\log_b2}.
\]
Taking \(k\)-th roots and letting \(k\to\infty\),
\[
\sqrt2\le A^{\log_b2}.
\]
Therefore \(A\ge\sqrt b\). Together with (5),
\[
\boxed{f(n)=\sqrt n\qquad(n\in\mathbb Z_{>0}).}
\]

\paragraph{4. Extend to rational numbers and verify.}

For positive integers \(m,n\), multiplicativity gives
\[
f(m/n)=\frac{f(m)}{f(n)}=\sqrt{m/n}.
\]
Evenness and \(f(0)=0\) complete the classification.

Conversely, \(f(x)=\sqrt{|x|}\) is multiplicative, has the required normalisations, and satisfies the corrected reciprocal condition with equality. Finally,
\[
\sqrt{|x+y|}
\le\sqrt{|x|+|y|}
\le\sqrt{|x|}+\sqrt{|y|},
\]
so it is subadditive.

\subsubsection{Final corrected statement}\label{final-corrected-statement}

Use the problem above with the reciprocal condition \textbf{\(f(x)f(1/x)\ge1\)}, not \(\le1\). With that correction, the classification and proof are complete.
\end{proof}

\paragraph{Takeaways.}
\begin{itemize}
\item **Check inequality directions carefully:** Reciprocal substitutions only create rigidity when the bounds oppose one another.
\item **Do not assume monotonicity:** Numeral expansions provide growth estimates without comparing \(f(a)\) and \(f(b)\) directly.
\item **Change the base strategically:** Binary expansions give one bound; base-\(b\) expansions give the reverse bound at \(b\).
\item **Amplify away constants:** Taking high powers eliminates fixed multiplicative losses.
\end{itemize}

\end{dlsummary}

\addcontentsline{toc}{subsubsection}{Solution 64. The Invisible Conjugate}

\begin{dlsummary}[title={\bfseries Solution 64. The Invisible Conjugate},breakable]

\begin{proof}
The only solutions are
\[
\boxed{f(a+b\sqrt2)=|a+b\sqrt2|}
\]
and
\[
\boxed{f(a+b\sqrt2)=|a-b\sqrt2|.}
\]

\paragraph{1. Establish symmetry and the reverse triangle inequality.}

Multiplicativity gives
\[
f(-1)^2=f(1)=1.
\]
Nonnegativity implies \(f(-1)=1\), hence
\[
f(-x)=f(x).
\]
Subadditivity now gives
\[
\boxed{|f(x)-f(y)|\le f(x-y).} \tag{1}
\]

\paragraph{2. Determine the function on the rationals.}

Repeated addition gives
\[
f(n)\le n\qquad(n\in\mathbb Z_{>0}),
\]
while multiplicativity gives \(f(2^k)=2^k\).

For any positive integer \(n\), choose \(2^k>n\). Then
\[
2^k=f(2^k)
\le f(n)+f(2^k-n)
\le f(n)+2^k-n.
\]
Thus \(f(n)\ge n\), proving equality.

Multiplicativity and symmetry consequently yield
\[
\boxed{f(q)=|q|\qquad(q\in\mathbb Q).} \tag{2}
\]
Also,
\[
f(\sqrt2)^2=f(2)=2,
\]
so \(f(\sqrt2)=\sqrt2\).

\paragraph{3. Force one conjugate direction to vanish.}

Choose rational numbers \(r_n\to\sqrt2\), and define
\[
A_n=f(\sqrt2-r_n),\qquad
B_n=f(\sqrt2+r_n).
\]

By (1) and (2),
\[
|A_n-A_m|\le|r_n-r_m|,
\qquad
|B_n-B_m|\le|r_n-r_m|.
\]
Therefore both sequences converge; write their limits as \(A,B\ge0\).

Multiplicativity gives
\[
A_nB_n
=f\bigl((\sqrt2-r_n)(\sqrt2+r_n)\bigr)
=f(2-r_n^2)
=|2-r_n^2|.
\]
Hence
\[
AB=0.
\]
But subadditivity gives
\[
A_n+B_n\ge f(2\sqrt2)=2\sqrt2,
\]
so
\[
A+B\ge2\sqrt2.
\]
Thus \textbf{exactly one} of \(A,B\) is zero.

\paragraph{4. The vanishing direction determines every value.}

If \(A=0\), then, for any rational \(a,b\), equation (1) gives
\[
\begin{aligned}
\left|f(a+b\sqrt2)-|a+br_n|\right|
&=\left|f(a+b\sqrt2)-f(a+br_n)\right|\\
&\le f\bigl(b(\sqrt2-r_n)\bigr)\\
&=|b|A_n.
\end{aligned}
\]
Taking limits,
\[
f(a+b\sqrt2)=|a+b\sqrt2|.
\]

If \(B=0\), compare instead with the rational number \(a-br_n\):
\[
\left|f(a+b\sqrt2)-|a-br_n|\right|
\le |b|B_n.
\]
Taking limits gives
\[
f(a+b\sqrt2)=|a-b\sqrt2|.
\]

This proves completeness.

\paragraph{5. Verify both solutions.}

The first is ordinary absolute value.

For the second, define
\[
\overline{a+b\sqrt2}=a-b\sqrt2.
\]
Direct expansion shows
\[
\overline{x+y}=\bar x+\bar y,
\qquad
\overline{xy}=\bar x\,\bar y.
\]
Therefore \(f(x)=|\bar x|\) satisfies both functional conditions and all normalisations.

\paragraph{A revealing example.}

At \(x=1+\sqrt2\), the two solutions give
\[
f(x)=1+\sqrt2
\quad\text{or}\quad
f(x)=\sqrt2-1.
\]
So an input larger than \(2\) can have a function value smaller than \(1\).
\end{proof}
\end{dlsummary}

\begin{takeaways}
\begin{itemize}
\item \textbf{Identify a familiar subdomain first:} Exact values on \(\mathbb Q\) provide the approximation machinery.
\item \textbf{Conjugate products can force a choice:} Two nonnegative limiting quantities have product zero but cannot both vanish.
\item \textbf{A local choice can determine everything:} Once one direction becomes small, rational comparison fixes all inputs.
\item \textbf{Do not assume ordinary continuity:} One solution is discontinuous relative to the usual real-number distance on \(K\).
\end{itemize}
\end{takeaways}

\addcontentsline{toc}{subsubsection}{Solution 65. Three Hidden Primes}

\begin{dlsummary}[title={\bfseries Solution 65. Three Hidden Primes},breakable]

\begin{proof}
For a prime \(p\), let \(v_p(x)\) denote the exponent of \(p\) in a nonzero rational \(x\), allowing negative exponents.

The three solutions are
\[
\boxed{
f(0)=0,\qquad f(x)=30^{-v_p(x)}\quad(x\ne0),
\qquad p\in\{2,3,5\}.
}
\]

\paragraph{1. Establish the basic identities.}

Since \(f(30)>0\), multiplicativity forces \(f(1)=1\). For \(x\ne0\),
\[
f(x)f(1/x)=1.
\]
Thus nonzero inputs have positive values.

Also,
\[
f(-1)^2=f(1)=1,
\]
so \(f(-1)=1\) and
\[
f(-x)=f(x). \tag{1}
\]

\paragraph{2. Obtain a uniform bound on integers.}

Let
\[
M=\max\{f(0),f(1),\ldots,f(29)\}.
\]
Write a positive integer \(N\) in base \(30\):
\[
N=d_0+d_1 30+\cdots+d_s30^s,
\qquad 0\le d_j\le29.
\]
Subadditivity and multiplicativity give
\[
f(N)
\le\sum_{j=0}^{s}f(d_j)f(30)^j
\le M\sum_{j=0}^{s}30^{-j}
<\frac{30M}{29}.
\]
Hence \(f\) is uniformly bounded on the positive integers.

Now apply this bound to \(N=n^r\):
\[
f(n)^r=f(n^r)<\frac{30M}{29}.
\]
Since this holds for every positive integer \(r\),
\[
\boxed{f(n)\le1\qquad(n\in\mathbb Z_{>0}).} \tag{2}
\]
By (1), the same bound holds for negative integers.

\paragraph{3. Locate a prime with value below \(1\).}

Multiplicativity gives
\[
f(2)f(3)f(5)=f(30)=\frac1{30}.
\]
All three factors belong to \((0,1]\), so at least one is below \(1\).

We next prove that \textbf{at most one prime anywhere} can have value below \(1\).

\paragraph{4. Rule out two exceptional primes.}

Suppose distinct primes \(p,q\) satisfy
\[
f(p)<1,\qquad f(q)<1.
\]
Choose a positive integer \(r\) sufficiently large that
\[
f(p)^r+f(q)^r<1.
\]

Since \(p^r\) and \(q^r\) are coprime, Bézout’s identity supplies integers \(a,b\) such that
\[
ap^r+bq^r=1.
\]
Using subadditivity, multiplicativity, and the integer bound,
\[
\begin{aligned}
1=f(1)
&\le f(a)f(p)^r+f(b)f(q)^r\\
&\le f(p)^r+f(q)^r\\
&<1,
\end{aligned}
\]
a contradiction.

Therefore exactly one prime \(p\) has \(f(p)<1\). It belongs to \(\{2,3,5\}\); every other prime has value \(1\). Consequently,
\[
f(p)=\frac1{30}.
\]

\paragraph{5. Determine every rational value.}

Unique prime factorisation and multiplicativity now give
\[
f(x)=30^{-v_p(x)}
\]
for every nonzero rational \(x\).

\paragraph{6. Verify the three candidates.}

Multiplicativity and the normalisations are immediate.

For nonzero \(x,y\) with \(x+y\ne0\),
\[
v_p(x+y)\ge\min\{v_p(x),v_p(y)\}.
\]
Therefore
\[
f(x+y)
\le\max\{f(x),f(y)\}
\le f(x)+f(y).
\]
Cases involving zero or \(x+y=0\) are immediate. All three functions work.
\end{proof}
\end{dlsummary}

\begin{takeaways}
\begin{itemize}
\item \textbf{A shrinking base controls all integers:} Base-\(30\) expansions turn \(f(30)<1\) into a uniform bound.
\item \textbf{Powers sharpen boundedness:} Uniformly bounded values of \(f(n)^r\) force \(f(n)\le1\).
\item \textbf{Bézout connects unrelated primes:} Two independently shrinking primes would make the value at \(1\) too small.
\item \textbf{One normalisation can leave several solutions:} The factorisation \(30=2\cdot3\cdot5\) creates exactly three possible controlling primes.
\end{itemize}
\end{takeaways}

\addcontentsline{toc}{subsubsection}{Solution 68. The Conditioned Bilinear Sign Test}

\begin{dlsummary}[title={\bfseries Solution 68. The Conditioned Bilinear Sign Test},breakable]

\begin{proof}
We claim that all solutions are
\[
\boxed{
f(x)=
\begin{cases}
x^p,&x>0,\\
0,&x=0,\\
-(-x)^p,&x<0,
\end{cases}
\qquad p\geq0.}
\]

\medskip
\noindent\textbf{1. Boundary identities.}
If $ab+cd=0$, both conditions apply, so
\[
f(a)f(b)+f(c)f(d)=0. \tag{1}
\]
Taking $(a,b,c,d)=(t,1,-t,1)$ gives
\[
f(t)+f(-t)=0,
\]
so $f$ is odd. Next, taking $(a,b,c,d)=(x,y,-xy,1)$ in (1) gives
\[
f(x)f(y)+f(-xy)=0.
\]
Using oddness, we obtain
\[
f(xy)=f(x)f(y)\qquad(x,y\in\mathbb R). \tag{2}
\]

\medskip
\noindent\textbf{2. Positivity and monotonicity.}
For $t>0$, (2) gives
\[
f(t)=f(\sqrt t)^2\geq0.
\]
Also, $f(t)f(1/t)=f(1)=1$, so in fact $f(t)>0$ for every $t>0$.
If $s\geq t$, apply the first condition to $(a,b,c,d)=(s,1,-t,1)$:
since $s-t\geq0$, it follows that
\[
f(s)-f(t)\geq0.
\]
Thus $f$ is nondecreasing on all of $\mathbb R$.

\medskip
\noindent\textbf{3. Classification on the positive half-line.}
Define
\[
g(u)=\log f(e^u)\qquad(u\in\mathbb R).
\]
By (2), $g$ is additive, and since $f$ is nondecreasing, so is $g$.
Put $p=g(1)$. Because $g(0)=0$ and $g$ is nondecreasing, $p\geq0$.

For positive integers $n$, additivity gives $g(n)=np$; also
$g(-n)=-np$, and $g(1)=g(n\cdot 1/n)=n g(1/n)$, so
$g(1/n)=p/n$. Hence, for every rational $q=m/n$,
\[
g(q)=pq. \tag{3}
\]
Now let $u\in\mathbb R$. For rationals $q\leq u\leq r$, monotonicity and (3) give
\[
pq=g(q)\leq g(u)\leq g(r)=pr.
\]
Choose rational $q$ and $r$ tending to $u$ from below and above. Since $p\geq0$,
the two bounds tend to $pu$, and therefore $g(u)=pu$. Consequently,
\[
f(t)=t^p\qquad(t>0). \tag{4}
\]
Oddness and $f(0)=0$ determine the values for $t\leq0$.

\medskip
\noindent\textbf{4. Verification.}
Each function in the displayed family is odd, multiplicative, and nondecreasing.
For $u,v\in\mathbb R$,
\[
f(a)f(b)+f(c)f(d)=f(ab)+f(cd),
\qquad u=ab,\ v=cd.
\]
If $u+v\geq0$, then $u\geq-v$, so monotonicity and oddness give
\[
f(u)\geq f(-v)=-f(v).
\]
Thus the sum is nonnegative. If $u+v\leq0$, the same comparison is reversed,
and the sum is nonpositive. Hence every displayed function satisfies both conditions.
\end{proof}
\end{dlsummary}

\begin{takeaways}
\begin{itemize}
\item The boundary case of two opposite inequalities produces exact identities: one
gives oddness and another gives multiplicativity. Positivity and a single ordered
substitution then yield global monotonicity. Taking logarithms converts the
positive multiplicative problem into a monotone additive one; the explicit
rational calculation and squeeze rule out pathological additive functions.
\end{itemize}
\end{takeaways}

\subsection*{Part II. Reciprocal and Harmonic Transforms}
\addcontentsline{toc}{subsection}{Part II. Reciprocal and Harmonic Transforms}

\addcontentsline{toc}{subsubsection}{Solution 6. The Sum--Product Squeeze}

\begin{dlsummary}[title={\bfseries Solution 6. The Sum--Product Squeeze},breakable]

\begin{proof}
Set
\[
g(x)=\frac1{f(x)},\qquad k=\frac{g(2)}2.
\]
The given inequalities are equivalent to
\begin{align}
g(x+y)&\ge k(g(x)+g(y)), \tag{1}\\
g(xy)&\le k\,g(x)g(y). \tag{2}
\end{align}

\textbf{Step 1.} Setting $x=y=1$ gives
\[
g(1)\le1,\qquad kg(1)\ge1,
\]
so $k\ge1$.

If $k>1$, repeated application of (1), with
increments $1/n$, yields
\[
g(2)\ge k^n g(1).
\]
Letting $n\to\infty$ gives a contradiction.
Thus $k=1$, and consequently
\[
g(1)=1,\qquad g(2)=2.
\]

\textbf{Step 2.} By (1) and (2),
\[
2g(x)\le g(2x)\le g(2)g(x)=2g(x).
\]
Hence
\[
\boxed{g(2x)=2g(x).}
\]
In particular, $g(2^n)=2^n$ for every
positive integer $n$.

\textbf{Step 3.} As $g$ is positive and
superadditive, it is strictly increasing.

Fix $x>0$ and put
\[
m_n=\left\lfloor\frac{2^n}{x}\right\rfloor.
\]
For all sufficiently large $n$,
\[
2^n=g(2^n)\ge g(m_nx)\ge m_ng(x).
\]
Therefore,
\[
g(x)\le
\frac{2^n}{\lfloor 2^n/x\rfloor}
\xrightarrow[n\to\infty]{}x.
\]
Thus $g(x)\le x$ for every $x>0$.

Conversely, by (2),
\[
1=g(1)\le g(x)g(1/x)\le\frac{g(x)}x,
\]
so $g(x)\ge x$.

We conclude that $g(x)=x$, whence
\[
\boxed{f(x)=\frac1x}.
\]

Finally, direct substitution gives
\[
f(x+y)(f(x)+f(y))
=2f(2)f(x)f(y)
=f(xy)=\frac1{xy}.
\]
Thus the proposed function satisfies
both inequalities.
\end{proof}

\end{dlsummary}
\addcontentsline{toc}{subsubsection}{Solution 7. The Harmonic Squeeze}

\begin{dlsummary}[title={\bfseries Solution 7. The Harmonic Squeeze},breakable]

\begin{proof}
Set
\[
g(x)=\frac1{f(x)}.
\]
Since all function values are positive, taking reciprocals in the given
inequalities gives
\[
g(x+y)\ge g(x)+g(y),
\qquad
g(xy)\le g(x)g(y),
\]
and the condition on $f(2)$ gives $g(2)\le2$.

\medskip
\noindent
\textbf{Step 1. Lock the dyadic values.}

Taking $x=y=1$ in submultiplicativity yields
\[
g(1)\le g(1)^2,
\]
so $g(1)\ge1$. Therefore
\[
2\ge g(2)\ge2g(1)\ge2,
\]
which implies $g(1)=1$ and $g(2)=2$. By induction, superadditivity and
submultiplicativity give
\[
g(2^n)=2^n
\qquad(n=0,1,2,\ldots). \tag{1}
\]
Moreover, $g$ is strictly increasing, since
\[
g(y+x)\ge g(y)+g(x)>g(y)
\qquad(x,y>0).
\]

\medskip
\noindent
\textbf{Step 2. Establish an upper bound.}

Fix $t>0$. For sufficiently large $n$, let
\[
m_n=\left\lfloor\frac{2^n}{t}\right\rfloor.
\]
Since $m_nt\le2^n$, monotonicity, superadditivity, and (1) yield
\[
2^n=g(2^n)\ge g(m_nt)\ge m_ng(t).
\]
As $2^n/m_n\to t$, it follows that
\[
g(t)\le t\qquad(t>0). \tag{2}
\]

\medskip

For every positive integer $k$, we have $g(k)=k$. If $x>1$, repeated
submultiplicativity and monotonicity give
\[
g(x)^n\ge g(x^n)
\ge g\bigl(\lfloor x^n\rfloor\bigr)
=\lfloor x^n\rfloor.
\]
Taking $n$-th roots and letting $n\to\infty$ gives $g(x)\ge x$.
Together with (2), this proves $g(x)=x$ for every $x>1$.

For $0<x<1$, submultiplicativity gives
\[
1=g(1)\le g(x)g\left(\frac1x\right)=\frac{g(x)}x,
\]
where the last equality follows from the result for $1/x>1$. Thus
$g(x)\ge x$, and (2) forces equality. Therefore $g(x)=x$ for all
$x>0$, and hence $f(x)=1/x$ for all $x>0$.
\end{proof}

\begin{proof}[Verification]
For $f(x)=1/x$,
\[
f(x+y)\bigl(f(x)+f(y)\bigr)
=\frac1{xy}
=f(x)f(y)
\le \frac1{xy}=f(xy),
\]
and $f(2)=1/2$. Thus the function satisfies all the conditions and is
the unique solution.
\end{proof}

\end{dlsummary}
\addcontentsline{toc}{subsubsection}{Solution 8. The Multiplicative Mirror}

\begin{dlsummary}[title={\bfseries Solution 8. The Multiplicative Mirror},breakable]

\begin{proof}
The reciprocal condition gives $f(1)=1$ and $f(1/x)=1/f(x)$. The
inequalities become
\[
\frac1{f(x+y)}\geq\frac1{f(x)}+\frac1{f(y)},\qquad
f(xy)f(x/y)\geq f(x)^2. \tag{1}
\]
Replacing $x$ by $1/x$ in the second inequality and using reciprocity
gives the reverse inequality. Hence
\[
f(xy)f(x/y)=f(x)^2.
\]
Applying this to $(u,v)$ and $(v,u)$ and multiplying yields
$f(uv)^2=f(u)^2f(v)^2$, so positivity gives $f(uv)=f(u)f(v)$.

Set $g=1/f$. Then $g$ is positive, increasing, superadditive,
multiplicative, and $g(1)=1$. The reciprocal condition also gives
\[
g(1/x)=\frac1{f(1/x)}=f(x)=\frac1{g(x)}.
\]
In particular $g(2)\geq2$. Moreover, using this identity with
$x=f(2)$,
\[
g(g(2))=\frac1{f(1/f(2))}=f(f(2))\leq2.
\]
Since $g(2)\geq2$ and $g$ is increasing,
$2\geq g(g(2))\geq g(2)\geq2$, so $g(2)=2$.

Multiplicativity gives $g(2^q)=2^q$ for every rational $q$. These
arguments form a dense subset of $(0,\infty)$; monotonicity therefore
forces $g(x)=x$ everywhere. Thus $f(x)=1/x$. Direct substitution
shows this function satisfies all conditions with equality.
\end{proof}

\end{dlsummary}
\addcontentsline{toc}{subsubsection}{Solution 9. The Reflected Product}

\begin{dlsummary}[title={\bfseries Solution 9. The Reflected Product},breakable]

\begin{proof}
\textbf{Step 1. Derive multiplicativity.}

The reciprocal identity at $x=1$ gives $f(1)=1$. Superadditivity shows that
$f$ is strictly increasing: if $v>u>0$, then
\[
f(v)\geq f(u)+f(v-u)>f(u),
\]
because $f(v-u)>0$. In particular,
\[
f(2)\geq f(1)+f(1)=2.
\]
Together with the given bound, this gives $f(2)=2$.

Apply the second inequality with $x$ replaced by $1/x$. Using the
reciprocal identity, we obtain
\[
\frac1{f(x)^2}\leq
\frac1{f(xy)f(x/y)},
\]
so $f(x)^2\geq f(xy)f(x/y)$. Hence
\[
f(x)^2=f(xy)f(x/y)\qquad(x,y>0). \tag{1}
\]
Taking $y=x$ in (1) gives $f(x^2)=f(x)^2$. Then, for any $a,b>0$,
use (1) with $x=\sqrt{ab}$ and $y=\sqrt{a/b}$ to get
\[
f(ab)=f(\sqrt{ab})^2=f(a)f(b).
\]
Thus $f$ is multiplicative.

\medskip
\noindent
\textbf{Method I. Dyadic squeeze.}

Multiplicativity and $f(2)=2$ give $f(2^n)=2^n$ for every integer
$n\geq0$. Fix $t>0$ and, for sufficiently large $n$, set
\[
m_n=\left\lfloor\frac{2^n}{t}\right\rfloor.
\]
Since $m_nt\leq2^n$, monotonicity and repeated superadditivity yield
\[
2^n=f(2^n)\geq f(m_nt)\geq m_nf(t).
\]
As $2^n/m_n\to t$, it follows that $f(t)\leq t$. The reciprocal
identity then gives
\[
f(t)=\frac1{f(1/t)}\geq t.
\]
Therefore $f(t)=t$ for all $t>0$.

\medskip
\noindent
\textbf{Method II. Dense fixed points.}

For every rational number $q=m/n$ with $m\in\mathbb Z$ and
$n\in\mathbb Z_{>0}$, multiplicativity implies
\[
f(2^q)^n=f(2^m)=2^m,
\]
where negative integer powers follow from $f(1)=1$ and
multiplicativity. Since $f(2^q)>0$, we get $f(2^q)=2^q$.
The set $\{2^q:q\in\mathbb Q\}$ is dense in $(0,\infty)$;
monotonicity and these fixed points imply $f(x)=x$ for every $x>0$.

Finally, $f(x)=x$ satisfies the reciprocal identity and $f(2)=2$;
both inequalities hold with equality. Hence it is the unique solution.
\end{proof}

\end{dlsummary}
\addcontentsline{toc}{subsubsection}{Solution 10. The Reciprocal Pinch}

\begin{dlsummary}[title={\bfseries Solution 10. The Reciprocal Pinch},breakable]

\begin{proof}
The reciprocal identity gives $f(1)=1$. By superadditivity,
\[
f(2)\geq 2f(1)=2.
\]
Writing $a=f(2)$, we have $a\geq2$ and $f(1/2)=1/a$. Thus
\[
a+\frac1a\leq\frac52,
\]
or $(2a-1)(a-2)\leq0$. Since $a\geq2$, this forces
\[
f(2)=2. \tag{1}
\]

Applying the multiplicative inequality to $1/x,1/y$ and using the
reciprocal identity gives
\[
\frac1{f(xy)}=f\left(\frac1{xy}\right)
\geq f\left(\frac1x\right)f\left(\frac1y\right)
=\frac1{f(x)f(y)}.
\]
Together with the original inequality, this proves
\[
f(xy)=f(x)f(y). \tag{2}
\]
Also, $f$ is strictly increasing: if $x<y$, then
\[
f(y)=f(x+(y-x))\geq f(x)+f(y-x)>f(x).
\]

By (2) and (1), $f(2^n)=2^n$ for every integer $n\geq0$. Fix $t>0$,
and for sufficiently large $n$ put
\[
m_n=\left\lfloor\frac{2^n}{t}\right\rfloor.
\]
Since $m_nt\leq2^n$, monotonicity and repeated superadditivity give
\[
2^n=f(2^n)\geq f(m_nt)\geq m_nf(t).
\]
As $2^n/m_n\to t$, we obtain $f(t)\leq t$. Applying this bound to
$1/t$ and using the reciprocal identity yields
\[
f(t)=\frac1{f(1/t)}\geq t.
\]
Hence $f(t)=t$ for every $t>0$.
\end{proof}

The identity function directly satisfies all the conditions, so it is
the unique solution.

\end{dlsummary}
\addcontentsline{toc}{subsubsection}{Solution 11. The Harmonic--Product Squeeze}

\begin{dlsummary}[title={\bfseries Solution 11. The Harmonic--Product Squeeze},breakable]

\begin{proof}
Define
\[
g(t)=\frac1{f(1/t)}\qquad(t>0).
\]
Taking reciprocals in the first inequality and substituting
$x=1/u$ and $y=1/v$ gives
\[
g(u+v)\geq g(u)+g(v). \tag{1}
\]
The second inequality, applied to $1/u$ and $1/v$ and then reciprocated,
gives
\[
g(uv)\leq g(u)g(v). \tag{2}
\]
Thus $g$ is positive, superadditive, and submultiplicative.

\medskip
\noindent
\textbf{Step 1. Obtain a lower bound for $g$.}

Taking $x=y=1$ in (2) gives $g(1)\leq g(1)^2$, so $g(1)\geq1$.
Also, (1) and positivity show that $g$ is strictly increasing: if
$y>x>0$, then
\[
g(y)=g(x+(y-x))\geq g(x)+g(y-x)>g(x).
\]
Fix $t>1$ and a positive integer $n$. Submultiplicativity and
superadditivity give
\[
g(t)^n\geq g(t^n)\geq\lfloor t^n\rfloor g(1)
\geq\lfloor t^n\rfloor.
\]
The middle inequality follows by writing $t^n$ as the sum of
$\lfloor t^n\rfloor$ copies of $1$ and, if needed, a positive
remainder. Taking $n$th roots and letting $n\to\infty$ yields
\[
g(t)\geq t\qquad(t>1). \tag{3}
\]

Let $a=1/c>1$. By the definition of $g$,
\[
g(a)=\frac1{f(c)},\qquad \frac1{g(a)}=f(c),
\]
and therefore
\[
g(g(a))=\frac1{f(1/g(a))}
       =\frac1{f(f(c))}\leq\frac1c=a.
\]
On the other hand, (3) applied first to $a$ and then to $g(a)>1$ gives
\[
g(g(a))\geq g(a)\geq a.
\]
Therefore,
\[
g(a)=a. \tag{4}
\]

\medskip
\noindent
\textbf{Step 2. Establish a global upper bound.}

By (2) and (4), $g(a^n)\leq g(a)^n=a^n$ for every positive integer
$n$. Fix $x>0$ and, for sufficiently large $n$, put
\[
m_n=\left\lfloor\frac{a^n}{x}\right\rfloor\geq1.
\]
Since $m_nx\leq a^n$, monotonicity and repeated superadditivity give
\[
a^n\geq g(a^n)\geq g(m_nx)\geq m_ng(x).
\]
Consequently,
\[
g(x)\leq\frac{a^n}{\lfloor a^n/x\rfloor}\longrightarrow x,
\]
so $g(x)\leq x$ for all $x>0$.

\medskip
\noindent
\textbf{Step 3. Force equality.}

The upper bound gives $g(1)\leq1$, and hence $g(1)=1$. By (2),
\[
1=g(1)=g(x\cdot1/x)\leq g(x)g(1/x)\leq\frac{g(x)}x.
\]
Thus $g(x)\geq x$. Together with the upper bound, this proves
$g(x)=x$ for all $x>0$. Returning to the definition of $g$,
\[
\frac1{f(1/x)}=x,
\]
so $f(x)=x$ for every $x>0$.

\medskip
\noindent
\textbf{Verification.}

For $f(x)=x$, both inequalities hold with equality, and
$f(f(c))=c$. Hence the identity function is indeed the unique solution.
\end{proof}

\end{dlsummary}
\addcontentsline{toc}{subsubsection}{Solution 54. The Harmonic Involution}

\begin{dlsummary}[title={\bfseries Solution 54. The Harmonic Involution},breakable]

\begin{proof}
The complete family of solutions is
\[
\boxed{f(x)=\frac{c}{x}\qquad(x>0),\quad c>0.}
\]

\paragraph{1. Recover strict monotonicity.}

Rewrite the inequality as
\[
f(x+y)\le\frac{f(x)f(y)}{f(x)+f(y)}<f(x).
\]
Whenever \(b>a>0\), take \(x=a\), \(y=b-a\). Thus
\[
f(b)<f(a),
\]
so \(f\) is strictly decreasing.

\paragraph{2. Obtain an integer scaling bound.}

Taking reciprocals gives
\[
\frac1{f(x+y)}\ge\frac1{f(x)}+\frac1{f(y)}.
\]
Repeated addition therefore yields
\[
\boxed{f(nx)\le\frac{f(x)}n} \tag{1}
\]
for every positive integer \(n\).

\paragraph{3. Make the scaling bound exact.}

Fix \(x>0\) and a positive integer \(n\). From (1),
\[
nf(nx)\le f(x).
\]
Since \(f\) is decreasing,
\[
f\bigl(nf(nx)\bigr)\ge f(f(x))=x. \tag{2}
\]

But applying (1) with input \(f(nx)\) gives
\[
f\bigl(nf(nx)\bigr)
\le\frac{f(f(nx))}{n}
=\frac{nx}{n}=x. \tag{3}
\]
Hence equality holds in (2) and (3). Strict monotonicity makes \(f\) injective, so
\[
nf(nx)=f(x).
\]
Therefore
\[
\boxed{f(nx)=\frac{f(x)}n.} \tag{4}
\]

\paragraph{4. Extend from rationals to all positive reals.}

Put \(c=f(1)>0\). Equation (4) gives
\[
f(1/n)=nc,
\qquad
f(m/n)=\frac{nc}{m}=\frac{c}{m/n}.
\]
Thus \(f(q)=c/q\) for every positive rational \(q\).

For any real \(x>0\), choose positive rationals \(r_k<x<s_k\) with both sequences tending to \(x\). Monotonicity gives
\[
\frac{c}{s_k}=f(s_k)\le f(x)\le f(r_k)=\frac{c}{r_k}.
\]
The bounds converge to \(c/x\), proving \(f(x)=c/x\).

\paragraph{5. Verify every member of the family.}

For any \(c>0\),
\[
f(f(x))=\frac{c}{c/x}=x,
\]
and
\[
f(x+y)\bigl(f(x)+f(y)\bigr)
=\frac{c}{x+y}\left(\frac cx+\frac cy\right)
=\frac{c^2}{xy}
=f(x)f(y).
\]
Thus every \(c>0\) works.
\end{proof}
\end{dlsummary}

\begin{takeaways}
\begin{itemize}
\item \emph{No takeaways were provided in the source draft.}
\end{itemize}
\end{takeaways}

\addcontentsline{toc}{subsubsection}{Solution 60. The Squared-Logarithm Trap}

\begin{dlsummary}[title={\bfseries Solution 60. The Squared-Logarithm Trap},breakable]

\begin{proof}
The unique solution is
\[
\boxed{
f(x)=2^{(\log_2 x)^2}
\qquad(x>0).
}
\]

\paragraph{1. Translate multiplication into addition.}

Define
\[
g(t)=\log_2 f(2^t)\qquad(t\in\mathbb R).
\]
The range and normalisations give
\[
g(t)\ge0,\qquad g(1)=1,\qquad
g(t)\le t\quad(0\le t\le1).
\tag{3}
\]

The two functional inequalities become
\[
g(u+v)+g(u-v)\le2g(u)+2g(v), \tag{4}
\]
and
\[
g(2u)\ge4g(u). \tag{5}
\]

\paragraph{2. Recover evenness and exact scaling.}

At \(u=0\), (5) and nonnegativity force \(g(0)=0\).

Taking \(u=0\) in (4),
\[
g(v)+g(-v)\le2g(v),
\]
so \(g(-v)\le g(v)\). Replacing \(v\) by \(-v\) proves
\[
g(-v)=g(v).
\]

Taking \(v=u\) in (4), then combining with (5), gives
\[
\boxed{g(2u)=4g(u).} \tag{6}
\]

\paragraph{3. Turn the inequality into the parallelogram equation.}

Apply (4) to \(u+v,u-v\):
\[
g(2u)+g(2v)\le2g(u+v)+2g(u-v).
\]
Using (6), this reverses (4). Hence
\[
\boxed{
g(u+v)+g(u-v)=2g(u)+2g(v).
}
\tag{7}
\]

\paragraph{4. Determine rational values.}

For fixed \(t\), equation (7) gives
\[
g((n+1)t)+g((n-1)t)=2g(nt)+2g(t).
\]
Induction from \(g(0)=0\) proves
\[
g(nt)=n^2g(t)
\]
for every integer \(n\).

Since \(g(1)=1\), it follows that
\[
\boxed{g(q)=q^2\qquad(q\in\mathbb Q).} \tag{8}
\]

\paragraph{5. Eliminate discontinuous possibilities.}

From evenness and (3),
\[
0\le g(h)\le |h|\qquad(|h|\le1).
\]
Thus \(g(h)\to0\) as \(h\to0\).

We now transfer this control to any fixed input \(x\). For any positive integer \(N\), use (7) with inputs \(x/N\) and \(h\). Nonnegativity and integer scaling give
\[
g(x+Nh)\le2g(x)+2N^2g(h).
\]

Also, (6) and (7) imply midpoint convexity:
\[
g\!\left(\frac{a+b}{2}\right)
\le\frac{g(a)+g(b)}2.
\]
For \(N\) a power of two, repeated midpoint convexity therefore gives
\[
\begin{aligned}
g(x+h)
&\le\left(1-\frac1N\right)g(x)
+\frac1N g(x+Nh)\\
&\le g(x)+\frac{g(x)}N+2Ng(h).
\end{aligned}
\]
The same upper bound holds for \(g(x-h)\). Since (7) gives
\[
g(x+h)+g(x-h)\ge2g(x),
\]
we conclude
\[
|g(x+h)-g(x)|
\le\frac{g(x)}N+2Ng(h).
\]
Choose \(N\) large first, then \(h\) small. This proves continuity at every \(x\).

Now (8) and rational approximation force
\[
g(t)=t^2\qquad(t\in\mathbb R).
\]
Undoing the transformation yields the claimed formula.

\paragraph{6. Verify.}

Write \(a=\log_2 x\) and \(b=\log_2 y\). Then
\[
(a+b)^2+(a-b)^2=2a^2+2b^2,
\]
so (1) holds with equality. Likewise, (2) holds with equality.

Finally, if \(1\le t\le2\), then \(s=\log_2t\in[0,1]\), so \(s^2\le s\). Therefore
\[
f(t)=2^{s^2}\le2^s=t.
\]
All conditions are satisfied.
\end{proof}

\paragraph{Takeaways.}
\begin{itemize}
\item **Transform both inputs and outputs:** Logarithms expose the additive structure behind a multiplicative inequality.
\item **Change coordinates to reverse a bound:** The substitution \((u+v,u-v)\) forces equality.
\item **Rational values are not the finish line:** A separate regularity argument is essential.
\item **A local condition can have global force:** The bound on \([1,2]\) eliminates all irregular quadratic solutions.
\end{itemize}

\end{dlsummary}

\addcontentsline{toc}{subsubsection}{Solution 61. The Hidden Entropy Function}

\begin{dlsummary}[title={\bfseries Solution 61. The Hidden Entropy Function},breakable]

\begin{proof}
The unique solution is
\[
\boxed{f(x)=x\ln x-x+1\qquad(x>0).}
\]

\paragraph{1. Find the forced normalisation.}

Putting \(x=y=1\),
\[
f(1)\ge2f(1).
\]
Nonnegativity therefore gives
\[
f(1)=0.
\]

\paragraph{2. Absorb the correction term.}

Define
\[
g(x)=\frac{f(x)+x-1}{x}.
\]
A direct rearrangement of the original inequality gives
\[
\boxed{g(xy)\ge g(x)+g(y).} \tag{1}
\]
Also, \(g(1)=0\), and nonnegativity of \(f\) implies
\[
\boxed{g(x)\ge1-\frac1x.} \tag{2}
\]

The problem is now to sharpen the elementary bound (2).

\paragraph{3. Subdivide multiplicatively.}

For every positive integer \(n\), repeated application of (1) gives
\[
g(x)\ge n\,g(x^{1/n}).
\]
Using (2),
\[
g(x)\ge n\left(1-x^{-1/n}\right).
\]
The elementary numerical limit
\[
\lim_{n\to\infty}n\left(1-x^{-1/n}\right)=\ln x
\]
therefore yields
\[
\boxed{g(x)\ge\ln x\qquad(x>0).} \tag{3}
\]

No continuity of \(g\) has been used: the limit concerns only the explicit expression on the right.

\paragraph{4. Use reciprocals to obtain the reverse bound.}

Apply (1) to \(x\) and \(1/x\):
\[
0=g(1)\ge g(x)+g(1/x).
\]
By (3),
\[
g(x)\le-g(1/x)\le-\ln(1/x)=\ln x.
\]
Thus \(g(x)=\ln x\). Undoing the normalisation,
\[
f(x)=xg(x)-x+1=x\ln x-x+1.
\]

\paragraph{5. Verify.}

First, the standard inequality
\[
\ln u\le u-1
\]
applied to \(u=1/x\) gives
\[
x\ln x-x+1\ge0.
\]
So the required range is satisfied.

Next, direct expansion gives
\[
\begin{aligned}
&f(xy)-xf(y)-yf(x)\\
&\qquad=xy-x-y+1\\
&\qquad=(x-1)(y-1).
\end{aligned}
\]
Hence the original inequality holds with equality.
\end{proof}
\end{dlsummary}

\begin{takeaways}
\begin{itemize}
\item \textbf{Choose the normalisation strategically:} Adding \(x-1\) and dividing by \(x\) removes both the weights and the correction term.
\item \textbf{Use the range as structural information:} Nonnegativity becomes the lower bound that drives the entire argument.
\item \textbf{Repeated roots sharpen a crude estimate:} Multiplicative subdivision turns \(1-1/x\) into \(\ln x\).
\item \textbf{Reciprocals force exactness:} Once the global lower bound is established, the upper bound follows immediately.
\end{itemize}
\end{takeaways}

\addcontentsline{toc}{subsubsection}{Solution 66. The Irrational Iteration Root}

\begin{dlsummary}[title={\bfseries Solution 66. The Irrational Iteration Root},breakable]

\begin{proof}
The unique solution is
\[
\boxed{f(x)=x^{\sqrt2}\qquad(x>0).}
\]

\paragraph{1. Recover strict monotonicity.}

If \(b>a>0\), superadditivity and positivity give
\[
f(b)=f(a+(b-a))
\ge f(a)+f(b-a)>f(a).
\]
Thus \(f\) is strictly increasing and, in particular, injective.

\paragraph{2. Force multiplicativity through composition.}

By submultiplicativity,
\[
f(xy)\le f(x)f(y).
\]
Applying the increasing function \(f\),
\[
f(f(xy))\le f(f(x)f(y)).
\]
Applying submultiplicativity again, now to \(f(x),f(y)\),
\[
f(f(x)f(y))
\le f(f(x))f(f(y)).
\]
The iteration identity therefore produces the squeeze
\[
\boxed{
x^2y^2
\le f(f(x)f(y))
\le x^2y^2.
}
\]
Hence
\[
f(f(x)f(y))=f(f(xy)).
\]
Injectivity gives
\[
\boxed{f(xy)=f(x)f(y).} \tag{1}
\]

This is the main breakthrough: the multiplicative inequality cannot remain strict anywhere.

\paragraph{3. Obtain the normalisation and rational powers.}

Taking \(x=y=1\) in (1), positivity yields
\[
f(1)=1.
\]
Also,
\[
f(2)\ge2f(1)=2.
\]
Define
\[
p=\log_2 f(2),
\]
so \(p\ge1\).

Multiplicativity gives, for every rational number \(r\),
\[
\boxed{f(2^r)=f(2)^r=2^{pr}.} \tag{2}
\]
For positive rational exponents, this follows by taking integer powers; reciprocal inputs handle negative exponents.

\paragraph{4. Determine every positive real input.}

Fix \(x>0\). Choose rational sequences \(r_n,s_n\) such that
\[
r_n<\log_2x<s_n,
\qquad r_n,s_n\longrightarrow\log_2x.
\]
Monotonicity and (2) give
\[
2^{pr_n}\le f(x)\le2^{ps_n}.
\]
The explicit bounds converge to \(x^p\). Therefore
\[
\boxed{f(x)=x^p\qquad(x>0).} \tag{3}
\]
No continuity of \(f\) was assumed.

Now the iteration identity becomes
\[
x^{p^2}=x^2\qquad(x>0).
\]
Thus \(p^2=2\), and \(p\ge1\) forces
\[
p=\sqrt2.
\]

\paragraph{5. Verify the solution.}

For \(f(x)=x^{\sqrt2}\), multiplicativity is immediate, and
\[
f(f(x))=(x^{\sqrt2})^{\sqrt2}=x^2.
\]

To verify superadditivity, note that for \(p\ge1\),
\[
\begin{aligned}
(x+y)^p
&=x(x+y)^{p-1}+y(x+y)^{p-1}\\
&\ge x\,x^{p-1}+y\,y^{p-1}\\
&=x^p+y^p.
\end{aligned}
\]
Taking \(p=\sqrt2\) completes the verification.
\end{proof}
\end{dlsummary}

\begin{takeaways}
\begin{itemize}
\item \textbf{Use inequalities at two levels:} First at \(x,y\), then at their function values.
\item \textbf{Iteration can close a squeeze:} The composition identity makes the two endpoints coincide exactly.
\item \textbf{Derive multiplicativity rather than guessing it:} This is the essential classification step.
\item \textbf{Separate regularity from algebra:} Positivity and superadditivity supply order; multiplicativity and rational approximation then identify the power.
\end{itemize}
\end{takeaways}


\subsection*{Part III. Iteration, Fixed Points, and Growth}
\addcontentsline{toc}{subsection}{Part III. Iteration, Fixed Points, and Growth}

\addcontentsline{toc}{subsubsection}{Solution 12. IMO 1977 -- Problem 6}

\begin{dlsummary}[title={\bfseries Solution 12. IMO 1977 -- Problem 6},breakable]

\begin{proof}
Start with $a_0=2$ and, while $a_j>1$, set $a_{j+1}=f(a_j-1)$. If
the sequence never reaches $1$, then
\[
f(a_j)>f(f(a_j-1))=f(a_{j+1})
\]
produces an infinite decreasing sequence of positive integers. Hence
$a_j=1$ for some $j$, so $f(t)=1$ for $t=a_{j-1}-1$. If $t>1$,
the condition at $t-1$ gives $1=f(t)>f(f(t-1))\geq1$, impossible.
Thus $t=1$ and $f(1)=1$.

Inductively, $f(n)\geq k$ for all $n\geq k$: if it holds for $k$,
then for $n\geq k+1$, $f(n-1)\geq k$ and
\[
f(n)>f(f(n-1))\geq k,
\]
so $f(n)\geq k+1$. In particular, $f(n)\geq n$.

It follows that $f(f(n))\geq f(n)$, hence $f(n+1)>f(n)$; thus $f$
is strictly increasing. If $f(n)>n$, then $f(n)\geq n+1$, so
$f(f(n))\geq f(n+1)$, contradicting the hypothesis. Therefore
$f(n)=n$ for every positive integer $n$, and direct substitution verifies
the solution.
\end{proof}

\end{dlsummary}
\addcontentsline{toc}{subsubsection}{Solution 13. The Fixed-Point Lever}

\begin{dlsummary}[title={\bfseries Solution 13. The Fixed-Point Lever},breakable]

\begin{proof}
Setting $y=x$ in the given inequality gives
\[
f(x)\bigl(x+f(x)\bigr)
\geq f(f(x))\bigl(f(x)+x\bigr).
\]
Since $x+f(x)>0$, it follows that
\[
f(f(x))\leq f(x). \tag{1}
\]

\medskip
\noindent
\textbf{Step 1. Use the assumed two-cycle.}

Let $b=f(a)$. Then $f(b)=a$. Applying (1) at $a$ gives
$a\leq b$, while applying it at $b$ gives
\[
b=f(f(b))\leq f(b)=a.
\]
Thus $a=b$ and $f(a)=a$.

Now set $x=a$ in the original inequality. Since
$f(a)=f(f(a))=a$, we get
\[
a\bigl(a+f(y)\bigr)\geq a(a+y),
\]
so
\[
f(y)\geq y\qquad(y>0). \tag{2}
\]

\medskip
\noindent
\textbf{Step 2. Force idempotence.}

Applying (2) to $f(x)>0$ yields $f(f(x))\geq f(x)$.
Together with (1), this gives
\[
f(f(x))=f(x). \tag{3}
\]

Substituting (3) into the original inequality and dividing by
$f(x)>0$ yields
\[
x+f(y)\geq f(x)+y.
\]
Therefore $f(x)-x\leq f(y)-y$ for all $x,y>0$.
Interchanging $x$ and $y$ shows that $f(x)-x$ is a constant
$c$. By (2), $c\geq0$.

Finally, $f(x)=x+c$ and (3) imply
\[
x+2c=f(f(x))=f(x)=x+c,
\]
so $c=0$. Hence $f(x)=x$ for every $x>0$.

\medskip
\noindent
\textbf{Verification.}

For $f(x)=x$, the inequality holds with equality, and
$f(f(a))=a$ for every $a>0$.
\end{proof}

\end{dlsummary}
\addcontentsline{toc}{subsubsection}{Solution 36. The Multiplicative Mirror: Positive Recurrence}

\begin{dlsummary}[title={\bfseries Solution 36. The Multiplicative Mirror: Positive Recurrence},breakable]

\begin{proof}
We prove that the only solutions are
\[
f(x)=x\qquad\text{and}\qquad f(x)=\frac1x.
\]

\medskip
\noindent\textbf{1. A sharp growth bound.}

Fix $t>1$ and write $a_n=f(t^n)$ for integers $n$. Substituting
$x=t$ and $y=t^n$ gives
\[
a_{n+1}+a_{n-1}\le (t+t^{-1})a_n. \tag{1}
\]
Put
\[
U_n=\frac{t^n-t^{-n}}{t-t^{-1}}.
\]
Iterating (1), and using positivity of all the terms, gives
\[
0<a_n\le U_n a_1-U_{n-1}a_0.
\]
Put
\[
b_n=a_{n+1}-\frac{a_n}{t}.
\]
The recurrence (1) rearranges to
\[
b_n\le t b_{n-1},
\]
so $b_n\le t^n b_0$. If $b_0<0$, then $b_n\to-\infty$, and
\[
a_{n+1}=\frac{a_n}{t}+b_n<0
\]
for all sufficiently large $n$, contradicting positivity. Hence
\[
a_1\ge\frac{a_0}{t}.
\]
Applying the same argument to the sequence $a_{-n}$ yields
\[
a_1\ge \frac{a_0}{t},\qquad a_{-1}\ge \frac{a_0}{t}.
\]
Since $a_0=f(1)=1$, (1) with $n=0$ also gives
$a_1+a_{-1}\le t+t^{-1}$. Consequently $a_1,a_{-1}\le t$. As $t>1$ was arbitrary,
\[
\min\{u,u^{-1}\}\le f(u)\le\max\{u,u^{-1}\}
\qquad(u>0). \tag{2}
\]

\medskip
\noindent\textbf{2. The pointwise alternatives.}

Taking logarithms in (2) gives
\[
|\log f(u)|\le |\log u|.
\]
Applying this inequality to $f(u)$ and using $f(f(u))=u$ gives the reverse
inequality. Thus
\[
|\log f(u)|=|\log u|,
\]
so
\[
f(u)\in\{u,u^{-1}\}\qquad(u>0). \tag{3}
\]

\medskip
\noindent\textbf{3. The choice cannot switch.}

For $u>1$, let
\[
L=\{u>1:f(u)=u^{-1}\}.
\]
If $u\in L$, substitute $y=u$ and $x=z/u$ in the original inequality. Using
(2) gives
\[
f(z)\le
\frac{z/u+u/z}{u}
-\min\left\{\frac{u^2}{z},\frac{z}{u^2}\right\}. \tag{4}
\]
As $z\to u$, the right-hand side tends to $u^{-1}<u$. Hence, for $z$
sufficiently close to $u$, the alternative $f(z)=z$ in (3) is impossible;
therefore $L$ is open in $(1,\infty)$.

Conversely, if $u_k\in L$ and $u_k\to u>1$, use (4) with centre $u_k$
and $z=u$, and pass to the limit. This gives $f(u)\le u^{-1}$, so (3)
implies $u\in L$. Thus $L$ is closed in $(1,\infty)$.

Since $(1,\infty)$ is connected, $L$ is either empty or all of $(1,\infty)$.
Finally, the involution makes $f$ injective, which forces the corresponding
choice on $(0,1)$ as well. The two possibilities are therefore precisely
\[
f(x)=x\qquad\text{or}\qquad f(x)=\frac1x.
\]
Direct substitution shows that both functions satisfy the inequality with
equality.
\end{proof}
\end{dlsummary}

\begin{takeaways}
\begin{itemize}
\item Multiplicative orbits turn the two-variable inequality into a recurrence.
Positivity makes the recurrence rigid, and applying the involution twice
turns a one-sided estimate into an equality. The final connectedness argument
is needed to exclude mixed pointwise choices.
\end{itemize}
\end{takeaways}
\addcontentsline{toc}{subsubsection}{Solution 37. The Crossed Contraction Lock}

\begin{dlsummary}[title={\bfseries Solution 37. The Crossed Contraction Lock},breakable]

\begin{proof}
The only solution is
\[
\boxed{f(x)=x\qquad(x>0).}
\]

Let $c=f(1)$. Substituting $y=1$ in the main inequality gives
\[
f(cx)+f(f(x))\le2x.
\]
Since $f(f(x))\ge x$, we obtain
\[
f(cx)\le x,
\]
or, equivalently,
\[
f(t)\le\frac tc\qquad(t>0). \tag{1}
\]

Applying (1) to $t=f(x)$ yields
\[
x\le f(f(x))\le\frac{f(x)}c,
\]
and therefore
\[
f(x)\ge cx. \tag{2}
\]

Taking $x=1$ in (1), we have $f(c)\le1$. On the other hand,
$f(c)=f(f(1))\ge1$, so $f(c)=1$. Now set $x=c$ in the main
inequality. Using $f(c)=1$ gives
\[
f\bigl(cf(y)\bigr)+f(y)\le2cy.
\]
By (2), $f(cf(y))\ge c^2f(y)$; hence
\[
(1+c^2)f(y)\le2cy,
\]
so
\[
f(y)\le ky,
\qquad k=\frac{2c}{1+c^2}\le1. \tag{3}
\]

Applying (3) twice gives
\[
f(f(x))\le k^2x.
\]
Together with $f(f(x))\ge x$, this forces $k=1$. Equality in
$2c\le1+c^2$ then gives $c=1$. Consequently, (2) and (3) become
\[
x\le f(x)\le x,
\]
so $f(x)=x$ for every $x>0$.

Finally, $f(x)=x$ satisfies the iteration condition and makes the
main inequality an equality.
\end{proof}
\end{dlsummary}

\begin{takeaways}
\begin{itemize}
\item The substitution $y=1$ converts the two-variable inequality into a
linear upper bound involving the constant $c=f(1)$. The iteration
condition turns that upper bound into a lower bound. A second special
substitution, $x=c$, produces the sharp contraction factor
$2c/(1+c^2)$, and iterating the contraction closes the argument.
\end{itemize}
\end{takeaways}
\addcontentsline{toc}{subsubsection}{Solution 14. IMO 2023 Shortlist -- Algebra A4}

\begin{dlsummary}[title={\bfseries Solution 14. IMO 2023 Shortlist -- Algebra A4},breakable]

\begin{proof}
Write $f^n(x)$ for the $n$-fold composition
of $f$, with $f^0(x)=x$.

\medskip
\noindent\textbf{Step 1.}
Setting $y=x$ gives
\[
2xf(x)\ge(f^2(x)+x)f(x).
\]
Since $f(x)>0$,
\[
\boxed{f^2(x)\le x.} \tag{1}
\]

\medskip
\noindent\textbf{Step 2.}
Setting $x=f(y)$ yields
\[
f(y)-f^3(y)\ge y-f^2(y).
\]
Replacing $y$ by $f^n(y)$ gives
\[
f^{n+1}(y)-f^{n+3}(y)
\ge f^n(y)-f^{n+2}(y).
\]

Define
\[
D_n=f^n(y)-f^{n+2}(y).
\]
Then
\[
D_{n+1}\ge D_n\ge D_0
=y-f^2(y)\ge0.
\]

For every positive integer $m$,
\begin{align*}
y-f^{2m}(y)
&=\sum_{k=0}^{m-1}D_{2k}\\
&\ge m\bigl(y-f^2(y)\bigr).
\end{align*}

Since $f^{2m}(y)>0$,
\[
y>m\bigl(y-f^2(y)\bigr).
\]
As $m$ is arbitrarily large,
\[
\boxed{f^2(y)=y.} \tag{2}
\]

\medskip
\noindent\textbf{Step 3.}
Using (2), the original inequality becomes
\[
x(f(x)+f(y))\ge(x+y)f(y),
\]
hence
\[
xf(x)\ge yf(y).
\]
Interchanging $x$ and $y$ gives equality:
\[
xf(x)=yf(y)=c
\]
for some $c>0$. Therefore,
\[
\boxed{f(x)=\frac{c}{x},\qquad c>0.}
\]

Conversely, for $f(x)=c/x$,
we have $f(f(x))=x$ and
\[
x(f(x)+f(y))
=c+\frac{cx}{y}
=(f(f(x))+y)f(y).
\]
Thus every such function satisfies the
given inequality.
\end{proof}

\end{dlsummary}
\addcontentsline{toc}{subsubsection}{Solution 15. FEOO 2020 Shortlist -- Algebra A1}

\begin{dlsummary}[title={\bfseries Solution 15. FEOO 2020 Shortlist -- Algebra A1},breakable]

\begin{proof}
Let $P(x,y)$ denote the given inequality.

\medskip
\noindent
\textbf{Step 1. Determine $f(1)$.}

From $P(1,1)$,
\[
2f(1)\geq1+f(1)^2,
\]
which gives
\[
(f(1)-1)^2\leq0.
\]
Thus
\[
\boxed{f(1)=1.} \tag{1}
\]

\medskip
\noindent
\textbf{Step 2. Establish a lower bound.}

By $P(x,1)$ and (1),
\[
2xf(x)\geq x^2+xf(x).
\]
Hence
\[
\boxed{xf(x)\geq x^2
\qquad(x\in\mathbb R).} \tag{2}
\]

\medskip
\noindent
\textbf{Step 3. Force equality.}

For $x\ne0$, applying $P(x,1/x)$ gives
\[
2f(1)\geq
1+xf(x)f\left(\frac1{x^2}\right).
\]
By (2) at $1/x^2$, we have
$f(1/x^2)\geq1/x^2>0$. Thus multiplying $xf(x)\geq x^2$ by this
positive value preserves the inequality. Using (1) and (2),
\[
\begin{aligned}
1
&\geq xf(x)f\left(\frac1{x^2}\right)\\
&\geq x^2f\left(\frac1{x^2}\right)\\
&\geq x^2\cdot\frac1{x^2}=1.
\end{aligned}
\]

All inequalities are equalities.
In particular,
\[
xf(x)=x^2.
\]
Therefore,
\[
\boxed{f(x)=x\qquad(x\ne0).} \tag{3}
\]

\medskip
\noindent
\textbf{Step 4. Determine $f(0)$.}

Setting $y=0$ gives
\[
0\geq xf(x)f(0).
\]
For $x\ne0$, by (3),
\[
0\geq x^2f(0).
\]
Consequently,
\[
\boxed{f(0)\leq0.}
\]

\medskip
\noindent
\textbf{Verification.}

Let
\[
f(x)=
\begin{cases}
x,&x\ne0,\\
c,&x=0,
\end{cases}
\qquad c\leq0.
\]

If $xy\ne0$, both sides of the original
inequality equal $2x^2y^2$.

If $x=0$, both sides vanish.

If $y=0$, the inequality becomes
\[
0\geq x^2c,
\]
which holds because $c\leq0$.

Thus all such functions satisfy the
given inequality.
\end{proof}

\end{dlsummary}
\addcontentsline{toc}{subsubsection}{Solution 16. The Iterated Ratio Squeeze}

\begin{dlsummary}[title={\bfseries Solution 16. The Iterated Ratio Squeeze},breakable]

\begin{proof}
For positive $u,v$, define
\[
R(u,v)=\frac uv+\frac vu.
\]
The first inequality in the statement says
\[
R(f(x),f(y))\geq R(x,y).
\]
Applying this inequality successively to
$(f(x),f(y))$ and $(f^{\circ2}(x),f^{\circ2}(y))$ gives
\[
R(f^{\circ3}(x),f^{\circ3}(y))
\geq R(f^{\circ2}(x),f^{\circ2}(y))
\geq R(f(x),f(y))
\geq R(x,y).
\]
The second inequality in the statement gives the reverse bound
$R(x,y)\geq R(f^{\circ3}(x),f^{\circ3}(y))$. Hence equality holds
throughout, and in particular
\[
R(f(x),f(y))=R(x,y). \tag{1}
\]

\medskip
\noindent
\textbf{Step 1. Determine the form of $f$.}

Let $c=f(1)>0$. Taking $y=1$ in (1),
\[
\frac{f(x)}c+\frac c{f(x)}=x+\frac1x.
\]
For positive $t$, the equation $t+t^{-1}=x+x^{-1}$ has roots
$t=x$ and $t=x^{-1}$. Therefore, for every $x>0$,
\[
f(x)=cx\quad\text{or}\quad f(x)=\frac cx. \tag{2}
\]

The choice in (2) cannot differ for two inputs other than $1$. Indeed,
if $u,v\ne1$ satisfied $f(u)=cu$ and $f(v)=c/v$, then (1) would give
\[
uv+\frac1{uv}=\frac uv+\frac vu,
\]
which is equivalent to $(u^2-1)(v^2-1)=0$, a contradiction. Thus the
same choice holds for every $x\ne1$. At $x=1$ the two expressions
coincide, so globally either $f(x)=cx$ or $f(x)=c/x$.

\medskip
\noindent
\textbf{Step 2. Apply the fixed-point condition.}

If $f(x)=cx$, then $f(a)=a$ gives $c=1$. If $f(x)=c/x$, then
$f(a)=a$ gives $c=a^2$. Hence the only possibilities are
\[
f(x)=x\qquad\text{or}\qquad f(x)=\frac{a^2}{x}.
\]

\medskip
\noindent
\textbf{Verification.}

Both functions satisfy $f(a)=a$ and preserve $R(x,y)$, so both
inequalities hold with equality. Therefore both are solutions.
\end{proof}

\end{dlsummary}
\addcontentsline{toc}{subsubsection}{Solution 17. IMO 2022 Shortlist -- Algebra A3}

\begin{dlsummary}[title={\bfseries Solution 17. IMO 2022 Shortlist -- Algebra A3},breakable]

\begin{proof}
Call a pair $(x,y)$ \emph{good} if
\[
xf(y)+yf(x)\leq2.
\]

\medskip
\noindent
\textbf{Step 1. Every good pair has $x=y$.}

Suppose $(x,y)$ is good with $x\ne y$.
By symmetry and uniqueness, neither
$(x,x)$ nor $(y,y)$ can be good. Thus
\[
xf(x)>1,\qquad yf(y)>1.
\]

By AM--GM,
\[
\begin{aligned}
xf(y)+yf(x)
&\geq2\sqrt{xf(y)\cdot yf(x)}\\
&=2\sqrt{xf(x)\cdot yf(y)}\\
&>2,
\end{aligned}
\]
a contradiction.

Therefore every good pair satisfies $x=y$.

\medskip
\noindent
\textbf{Step 2. Establish an upper bound.}

For every $x>0$, a good pair exists.
By Step 1, it must be $(x,x)$. Hence
\[
2xf(x)\leq2,
\]
so
\[
\boxed{f(x)\leq\frac1x.} \tag{1}
\]

\medskip
\noindent
\textbf{Step 3. Force equality.}

Fix $x>0$ and let
\[
t=\frac1{f(x)}>0.
\]

By (1),
\[
f(t)\leq\frac1t=f(x).
\]
Therefore,
\[
\begin{aligned}
xf(t)+tf(x)
&\leq xf(x)+1\\
&\leq2.
\end{aligned}
\]

Thus $(x,t)$ is a good pair. By Step 1,
$t=x$, which gives
\[
\boxed{f(x)=\frac1x\qquad(x>0).}
\]

\medskip
\noindent
\textbf{Verification.}

For $f(x)=1/x$, the given inequality becomes
\[
\frac{x}{y}+\frac{y}{x}\leq2.
\]

However, AM--GM gives
\[
\frac{x}{y}+\frac{y}{x}\geq2,
\]
with equality if and only if $x=y$.

Hence, for every $x>0$, there is exactly
one admissible $y>0$, namely $y=x$.

Therefore $f(x)=1/x$ is the unique solution.
\end{proof}

\medskip
\noindent
\textbf{Alternative 1.1. A direct contradiction.}

We use the bound $f(x)\leq1/x$ established above. Suppose that
$f(x)<1/x$ for some $x>0$. Then the interval $(x,1/f(x))$ is nonempty.
For every $a$ in this interval,
\[
af(x)<1,\qquad xf(a)\leq\frac{x}{a}<1,
\]
so
\[
af(x)+xf(a)<2.
\]
Thus every such $a$ is a good partner of $x$. Since the interval contains
more than one number, this contradicts uniqueness. Hence $f(x)=1/x$.

\medskip
\noindent
\textbf{Alternative 1.2. Strict decrease.}

We again use the fact that every good pair is diagonal and
$f(x)\leq1/x$. If $y>x$, then $(x,y)$ is not good, whereas $(y,y)$ is
good. Therefore
\[
xf(y)+yf(x)>2\geq2yf(y)>xf(y)+yf(y),
\]
which gives $f(x)>f(y)$. Thus $f$ is strictly decreasing.

For any $y>x$, strict decrease and the fact that $(x,y)$ is not good give
\[
2<xf(y)+yf(x)<(x+y)f(x).
\]
Letting $y\to x^+$ yields $f(x)\geq1/x$. Together with the upper bound,
this proves $f(x)=1/x$.

\medskip
\noindent
\textbf{Alternative 2. A partner map.}

For each $x>0$, let $g(x)$ be the unique good partner of $x$. Symmetry of
the condition and uniqueness imply $g(g(x))=x$.

We first show that $f$ is strictly decreasing. If $x<y$, then
$(g(y),y)$ is good, but $(g(y),x)$ is not, since $y$ is the unique good
partner of $g(y)$. Hence
\[
xf(g(y))+g(y)f(x)>2\geq yf(g(y))+g(y)f(y).
\]
Rearranging shows $f(x)>f(y)$.

Now suppose that $(x,y)$ is a good pair with $x\ne y$. Uniqueness
implies that neither $(x,x)$ nor $(y,y)$ is good, so
$xf(x)+yf(y)>2$. But $(x,y)$ being good gives
\[
(x-y)(f(x)-f(y))
=xf(x)+yf(y)-xf(y)-yf(x)>0,
\]
contradicting strict decrease. Thus every good pair is diagonal, giving
$f(x)\leq1/x$. Alternative 1.1 now forces equality.

\medskip
\noindent
\textbf{Alternative 3. Equality at the unique partner.}

Let $g(x)$ again denote the unique good partner of $x$. Symmetry gives
$g(g(x))=x$. We claim that every good pair satisfies equality. Suppose
instead that $(x,y)$ is good and
\[
\delta=2-xf(y)-yf(x)>0.
\]
For $0<\varepsilon<\delta/(2f(x))$, uniqueness makes $(x,y+\varepsilon)$
not good, and therefore
\[
f(y+\varepsilon)
>\frac{2-(y+\varepsilon)f(x)}{x}
>f(y)+\frac{\delta}{2x}.
\]
Put $z=y+\varepsilon$. The good pair $(z,g(z))$ and the fact that
$(g(z),y)$ is not good give
\[
f(g(z))\leq\frac{2-g(z)f(z)}{z},
\qquad
f(g(z))>\frac{2-g(z)f(y)}{y}.
\]
Combining these inequalities yields
\[
2\varepsilon<g(z)\bigl[zf(y)-yf(z)\bigr]
<g(z)\left[\varepsilon f(y)-\frac{y\delta}{2x}\right].
\]
For sufficiently small positive $\varepsilon$, the final bracket is
negative, a contradiction. Thus every good pair has value exactly $2$.

Every pair other than $(x,g(x))$ is not good, so
\[
xf(y)+yf(x)\geq2\qquad(x,y>0).
\]
Taking $y=x$ gives $f(x)\geq1/x$. Applying this bound to the good pair
$(x,g(x))$ gives
\[
2\geq xf(g(x))+g(x)f(x)
\geq\frac{x}{g(x)}+\frac{g(x)}{x}\geq2.
\]
Equality throughout implies $f(x)=1/x$.

\medskip
\noindent
\textbf{Alternative 4. Strict convexity and continuity.}

Let $g(w)$ be the unique good partner of $w$, and define
\[
q_s(t)=f(t)+st.
\]
For $s=f(g(w))/g(w)$, the good-pair condition gives
$q_s(w)\leq2/g(w)$. If $t\ne w$, then $(g(w),t)$ is not good, so
$q_s(t)>2/g(w)$. Thus $q_s$ has a unique global minimum at $w$.

We claim that $f$ is strictly convex. Otherwise, for some $u<v$ and
$\lambda\in(0,1)$, setting $w=\lambda u+(1-\lambda)v$ would give
\[
f(w)\geq\lambda f(u)+(1-\lambda)f(v).
\]
Adding the linear term $st$ yields
$q_s(w)\geq\lambda q_s(u)+(1-\lambda)q_s(v)$, contradicting the
strict minimality of $q_s$ at $w$. Hence $f$ is strictly convex and, as
a finite convex function on an open interval, is continuous.

As in the first solution, every good pair is diagonal, so $f(x)\leq1/x$.
If this inequality were strict at some $x$, the continuous function
$h(y)=xf(y)+yf(x)$ would satisfy $h(x)=2xf(x)<2$. It would then be less
than $2$ for all $y$ in some interval around $x$, contradicting the
uniqueness of the good partner. Therefore $f(x)=1/x$ for all $x>0$.

\end{dlsummary}
\addcontentsline{toc}{subsubsection}{Solution 18. The Geometric-Mean Trap}

\begin{dlsummary}[title={\bfseries Solution 18. The Geometric-Mean Trap},breakable]

\begin{proof}
Write $P(x,y)$ for the inequality in the statement. It is symmetric
in $x$ and $y$. On the diagonal, $P(x,x)$ is equivalent to
\[
\frac{2x}{f(x)}\leq\frac{2f(x)}x,
\]
which, since $x,f(x)>0$, is equivalent to
\[
f(x)\geq x. \tag{1}
\]

\medskip
\noindent
\textbf{Step 1. Every admissible pair is diagonal.}

Suppose $P(x,y)$ holds for some $x\ne y$. If $f(x)\geq x$, then
$P(x,x)$ also holds, contradicting uniqueness for this fixed $x$.
Thus $f(x)<x$. By symmetry, $f(y)<y$ as well. Therefore,
\[
\frac{x}{f(y)}+\frac{y}{f(x)}
>\frac xy+\frac yx\geq2,
\]
whereas
\[
\frac{f(x)}x+\frac{f(y)}y<2.
\]
This contradicts $P(x,y)$. Hence every admissible pair has $x=y$.

For each $x>0$, an admissible $y$ exists, so it must equal $x$.
Consequently, $P(x,x)$ holds for every $x>0$, and (1) gives
\[
f(x)\geq x\qquad(x>0). \tag{2}
\]

\medskip
\noindent
\textbf{Step 2. Use a geometric-mean choice.}

Suppose $f(a)>a$ for some $a>0$, and put
\[
u=\frac{f(a)}a>1,
\qquad y=a\sqrt u=\sqrt{a f(a)}.
\]
By (2), $f(y)\geq y$. Hence
\[
\frac{a}{f(y)}+\frac{y}{f(a)}
\leq\frac ay+\frac{y}{f(a)}
=\frac1{\sqrt u}+\frac1{\sqrt u}
=\frac2{\sqrt u}<2.
\]
On the other hand,
\[
\frac{f(a)}a+\frac{f(y)}y\geq u+1>2.
\]
Thus $P(a,y)$ holds. Since $u>1$, $y\ne a$, contradicting Step 1.
Therefore $f(a)\leq a$ for every $a>0$. Together with (2), this
proves $f(x)=x$ for all $x>0$.

\medskip
\noindent
\textbf{Verification.}

For $f(x)=x$, the inequality becomes
\[
\frac xy+\frac yx\leq2.
\]
By AM--GM, the left-hand side is at least $2$, with equality if and
only if $x=y$. Thus exactly one $y>0$ works for each $x>0$, namely
$y=x$.
\end{proof}

\end{dlsummary}
\addcontentsline{toc}{subsubsection}{Solution 19. The Reciprocal-Slope Pinch}

\begin{dlsummary}[title={\bfseries Solution 19. The Reciprocal-Slope Pinch},breakable]

\begin{proof}
\textbf{Step 1. Establish a lower bound.}

Since $f(xf(x))>0$, the given inequality implies
\[
f(x+y)>y\qquad(x,y>0).
\]
Fix $z>0$ and take any $x\in(0,z)$ with $y=z-x$. Then
$f(z)>z-x$. As this holds for every such $x$, it follows that
\[
f(z)\geq z\qquad(z>0). \tag{1}
\]

\medskip
\noindent
\textbf{Step 2. Control the excess over the identity.}

Define $g(x)=f(x)-x$, so $g(x)\geq0$ by (1). Since
$f(xf(x))=xf(x)+g(xf(x))$, the original inequality is equivalent to
\[
g(x+y)-g(x)\geq\frac{g(xf(x))}{x}\qquad(x,y>0). \tag{2}
\]
In particular, $g$ is nondecreasing.

If $f(1)>1$, taking $x=1$ and $y=f(1)-1$ in the original inequality
would give
\[
f(1)-1+f(f(1))\leq f(f(1)),
\]
which is impossible. Hence $f(1)\leq1$, and (1) gives $f(1)=1$;
therefore $g(1)=0$. Since $g$ is nonnegative and nondecreasing,
\[
g(x)=0\qquad(0<x\leq1). \tag{3}
\]

\medskip
\noindent
\textbf{Step 3. Rule out positive excess above $1$.}

Suppose $g(a)=c>0$ for some $a>1$. For every $u>1$, we have
$uf(u)\geq u^2>u$. Monotonicity of $g$ and (2) therefore imply, for
all $v>u>1$,
\[
g(v)\geq g(u)+\frac{g(uf(u))}{u}
\geq g(u)+\frac{g(u)}{u}. \tag{4}
\]
Fix $X>a$. For any positive integer $N$, choose
\[
u_j=a+\frac{j(X-a)}{N+1}\qquad(0\leq j\leq N).
\]
All these points lie in $[a,X)$, so $g(u_j)\geq c$ and
$u_j<X$. Applying (4) to each successive pair gives
\[
g(u_{j+1})\geq g(u_j)+\frac{c}{X}
\qquad(0\leq j<N).
\]
Thus $g(u_N)\geq c+Nc/X$. But monotonicity gives $g(u_N)\leq g(X)$,
which is a fixed finite value. Choosing $N$ sufficiently large is a
contradiction. Hence $g(x)=0$ for all $x>1$ as well. Together with (3),
this proves $f(x)=x$ for every $x>0$.

Finally, the identity function satisfies the original inequality with
equality throughout. It is therefore the unique solution.
\end{proof}

\end{dlsummary}
\addcontentsline{toc}{subsubsection}{Solution 20. The Cross-Dilated Shift}

\begin{dlsummary}[title={\bfseries Solution 20. The Cross-Dilated Shift},breakable]

\begin{proof}
The identity function satisfies the inequality with equality. We prove it
is the only solution.

\medskip
\noindent
\textbf{Step 1. Establish a lower bound.}

Since $f(xf(y))>0$, the given inequality implies $f(x+y)>x$ for all
$x,y>0$. Fix $z>0$. For each $\varepsilon\in(0,z)$, take
$x=z-\varepsilon$ and $y=\varepsilon$. Then $f(z)>z-\varepsilon$;
hence
\[
f(z)\geq z\qquad(z>0). \tag{1}
\]

\medskip
\noindent
\textbf{Step 2. Bound the excess over the identity.}

Suppose $f(a)>a+1$ for some $a>0$, and set
\[
x=\frac{a}{f(a)-1}.
\]
Then $0<x<1$ and $xf(a)=x+a$. Applying the given inequality at
$(x,a)$ gives
\[
x^2+f(x+a)\leq x f(x+a),
\]
which is impossible because $f(x+a)>0$ and $0<x<1$. Therefore
$f(a)\leq a+1$ for all $a>0$.

Define $g(t)=f(t)-t$. By (1),
\[
0\leq g(t)\leq1\qquad(t>0). \tag{2}
\]
Substituting $f(t)=t+g(t)$ into the original inequality and cancelling
the common terms yields
\[
xg(y)+g\bigl(xf(y)\bigr)\leq xg(x+y)
\qquad(x,y>0). \tag{3}
\]

\medskip
\noindent
\textbf{Step 3. Force the excess to vanish.}

Since $g$ is nonnegative, (3) implies $g(y)\leq g(x+y)$ for all
$x,y>0$, so $g$ is nondecreasing. Taking $x=1$ in (3), and using
$f(y)\geq y$ and the monotonicity of $g$, gives
\[
2g(y)\leq g(y)+g\bigl(f(y)\bigr)\leq g(y+1).
\]
Inductively, $g(y+n)\geq2^n g(y)$ for every positive integer $n$.
By (2), this forces $g(y)=0$ for every $y>0$. Thus $f(y)=y$ on
$(0,\infty)$, as claimed.
\end{proof}

\end{dlsummary}
\addcontentsline{toc}{subsubsection}{Solution 40. The Monotonicity Ratchet}

\begin{dlsummary}[title={\bfseries Solution 40. The Monotonicity Ratchet},breakable]

\begin{proof}
The only solution is \(f(x)=x\). It satisfies both inequalities with equality.

\paragraph{1. A first consequence.} In the first inequality, \(f(x)>0\), so \(f(x+f(y))>y\). As \(x\) ranges over all positive reals, \(t=x+f(y)\) ranges over all \(t>f(y)\). Therefore
\[
t>f(y)\ \Longrightarrow\ f(t)>y. \tag{1}
\]

\paragraph{2. \(f\) is strictly increasing.} Take \(0<w<t\). By the second condition, \(f(f(w))\le w<t\). Apply (1) with \(y=f(w)\). Since \(t>f(f(w))\), we get
\[
f(t)>f(w).
\]

\paragraph{3. \(f(y)\le y\).} Suppose \(f(y)>y\) for some \(y\). Since \(f\) is strictly increasing, \(f(f(y))>f(y)>y\). This contradicts \(f(f(y))\le y\).

\paragraph{4. Collapse to a constant defect.} Using step 3 and the first condition,
\[
x+f(y)\ge f\bigl(x+f(y)\bigr)\ge f(x)+y.
\]
Rearranging gives
\[
f(y)-y\ge f(x)-x\qquad\text{for all }x,y>0.
\]
Swapping the roles of \(x\) and \(y\) gives the reverse inequality. Hence \(f(x)-x=-k\) is constant, and step 3 gives \(k\ge0\).

\paragraph{5. Fix the constant.} If \(k>0\), then \(f(k)=0\) with \(k>0\), which contradicts positivity of \(f\). So \(k=0\) and \(f(x)=x\). \(\blacksquare\)
\end{proof}
\end{dlsummary}

\begin{takeaways}
\begin{itemize}
\item \textbf{Derive monotonicity:} Substituting \(y=f(w)\) turned a composition bound into a statement about order.
\item \textbf{Use order against iteration:} For an increasing function, \(f(y)>y\) would force \(f(f(y))>y\), which contradicts an upper bound on \(f(f(y))\).
\item \textbf{Use the symmetry trick again:} Once \(f(y)-y\ge f(x)-x\) holds for all \(x,y\), swapping them gives equality.
\item \textbf{Let the domain finish the job:} Positivity of \(f\) removes the family \(f(x)=x-k\).
\end{itemize}
\end{takeaways}

\addcontentsline{toc}{subsubsection}{Solution 43. The Moving Zero}

\begin{dlsummary}[title={\bfseries Solution 43. The Moving Zero},breakable]

\begin{proof}
The only solution is
\[
\boxed{f(x)=x^2\qquad(x\in\mathbb R).}
\]

\paragraph{1. Find a zero—but do not assume it is \(0\).}

Set
\[
a=f(0)\ge0.
\]
By (2) at \(x=0\),
\[
0\le f(a)=f(f(0))\le0,
\]
so \(f(a)=0\).

\paragraph{2. Build a quadratic lower bound around that zero.}

In (1), substitute
\[
x=\frac{t+a}{2},\qquad y=\frac{t-a}{2}.
\]
Since \(f(a)=0\),
\[
f(t)\ge 2f\!\left(\frac{t+a}{2}\right)+\frac{(t-a)^2}{2}.
\]
Repeating this midpoint substitution \(n\) times gives
\[
f(t)\ge
2^n f\!\left(a+\frac{t-a}{2^n}\right)
+(1-2^{-n})(t-a)^2.
\]
Nonnegativity lets us discard the first term. Letting \(n\to\infty\), we obtain
\[
\boxed{f(t)\ge(t-a)^2\qquad(t\in\mathbb R).} \tag{3}
\]
No limit of function values is being taken, so continuity is unnecessary.

\paragraph{3. Turn the lower bound into an upper bound.}

Apply (3) at \(t=f(x)\), then use (2):
\[
(f(x)-a)^2\le f(f(x))\le x^4.
\]
Therefore
\[
|f(x)-a|\le x^2,
\qquad\text{in particular}\qquad
f(x)\le a+x^2. \tag{4}
\]

On the other hand, (1) at \(x=0\) gives
\[
f(y)+f(-y)\ge2a+2y^2.
\]
Each summand is at most \(a+y^2\) by (4). Both must therefore attain their upper bounds:
\[
f(y)=a+y^2\qquad(y\in\mathbb R).
\]
Since \(f(a)=0\),
\[
a+a^2=0.
\]
Together with \(a\ge0\), this forces \(a=0\). Hence \(f(x)=x^2\).

Finally, \(f(x)=x^2\) satisfies both original conditions with equality.
\end{proof}

\paragraph{Takeaways.}
\begin{itemize}
\item **Locate a zero before normalising:** The composition condition supplies a zero at \(f(0)\), not immediately at \(0\).
\item **Iterate a midpoint inequality:** Repeated bisection produces a sharp bound without establishing continuity.
\item **Feed bounds into compositions:** A lower bound for \(f(t)\) becomes an upper bound for \(f(x)\) through \(f(f(x))\).
\item **Exploit equality of sums:** If two terms each have the same upper bound and their sum reaches the total bound, both are determined.
\end{itemize}

\end{dlsummary}

\addcontentsline{toc}{subsubsection}{Solution 44. The Hidden Golden Exponent}

\begin{dlsummary}[title={\bfseries Solution 44. The Hidden Golden Exponent},breakable]

\begin{proof}
The unique solution is
\[
\boxed{
f(x)=x^\varphi,\qquad
\varphi=\frac{1+\sqrt5}{2}.
}
\]

\paragraph{1. Extract superadditivity.}

Multiplicativity and positivity give \(f(1)=1\). The inequality with \(y=1\) becomes
\[
f(x+1)\ge f(x)+1.
\]
For \(u,v>0\), substitute \(x=u/v\), then multiply by \(f(v)\). Multiplicativity yields
\[
\boxed{f(u+v)\ge f(u)+f(v).} \tag{1}
\]
In particular, \(f\) is strictly increasing: if \(b>a>0\), then
\[
f(b)\ge f(a)+f(b-a)>f(a).
\]

\paragraph{2. Show that \(f\) is a power function.}

Let
\[
p=\frac{\log f(2)}{\log2}.
\]
Since \(f(2)\ge2\), we have \(p\ge1\).

Multiplicativity gives, for every rational \(r\),
\[
f(2^r)=f(2)^r=2^{pr}.
\]
For any \(t>0\), choose rational numbers \(r_n,s_n\) such that
\[
r_n<\log_2t<s_n,
\qquad r_n,s_n\longrightarrow\log_2t.
\]
Monotonicity implies
\[
2^{pr_n}\le f(t)\le2^{ps_n}.
\]
Thus
\[
\boxed{f(t)=t^p\qquad(t>0).} \tag{2}
\]
This establishes continuity as a consequence, rather than assuming it.

\paragraph{3. Pin down the exponent.}

Substituting (2) into the original inequality gives
\[
(x+y^p)^p\ge x^p+y^{p+1}.
\]
For fixed \(y>0\), let \(x\to0^+\). We obtain
\[
y^{p^2}\ge y^{p+1}\qquad(y>0).
\]
Choosing \(y>1\) gives \(p^2\ge p+1\), while choosing \(0<y<1\) gives the reverse inequality. Hence
\[
p^2=p+1.
\]
Since \(p\ge1\),
\[
p=\varphi=\frac{1+\sqrt5}{2}.
\]

\paragraph{4. Verify the solution.}

Multiplicativity is immediate. For any \(a,b>0\) and \(p\ge1\),
\[
(a+b)^p
=a(a+b)^{p-1}+b(a+b)^{p-1}
\ge a^p+b^p.
\]
Taking \(a=x\), \(b=y^\varphi\), and using \(\varphi^2=\varphi+1\), we get
\[
(x+y^\varphi)^\varphi
\ge x^\varphi+y^{\varphi^2}
=x^\varphi+y^{\varphi+1},
\]
exactly the required inequality.
\end{proof}
\end{dlsummary}

\begin{takeaways}
\begin{itemize}
\item \textbf{A unit shift can control arbitrary sums:} Multiplicativity converts \(f(x+1)\ge f(x)+1\) into superadditivity.
\item \textbf{Positivity supplies regularity:} Superadditivity forces monotonicity, eliminating nonstandard multiplicative functions.
\item \textbf{Use both sides of \(1\):} A power comparison valid for every positive input often forces equality of exponents.
\item \textbf{Composition multiplies exponents:} The competition between \(f(f(y))\) and \(y f(y)\) produces \(p^2=p+1\).
\end{itemize}
\end{takeaways}

\addcontentsline{toc}{subsubsection}{Solution 45. The Conjugate-Radical Lock}

\begin{dlsummary}[title={\bfseries Solution 45. The Conjugate-Radical Lock},breakable]

\begin{proof}
The unique solution is
\[
\boxed{
f(x)=\frac{(2+\sqrt3)^x+(2-\sqrt3)^x}{2}
\qquad(x\ge0).
}
\]

\paragraph{1. Remove the radicals.}

For \(t\ge1\), define
\[
H(t)=t+\sqrt{t^2-1}.
\]
This function is strictly increasing, and
\[
H(t)^{-1}=t-\sqrt{t^2-1}.
\]

For \(a,b\ge1\), direct expansion gives
\[
H\!\left(ab+\sqrt{a^2-1}\sqrt{b^2-1}\right)=H(a)H(b).
\tag{3}
\]
Also,
\[
H(2a^2-1)=H(a)^2. \tag{4}
\]

Set
\[
g(x)=\log H(f(x)).
\]
Then \(g(x)\ge0\). Applying \(H\) and taking logarithms in (1) and (2), using (3) and (4), gives
\[
g(x+y)\ge g(x)+g(y),
\qquad
g(2x)\le2g(x).
\]
The first inequality with \(y=x\) supplies the reverse doubling bound. Therefore
\[
\boxed{g(2x)=2g(x).} \tag{5}
\]

\paragraph{2. Recover monotonicity and rational values.}

Since \(g\) is nonnegative and superadditive, it is nondecreasing:
\[
b>a\ge0
\implies
g(b)\ge g(a)+g(b-a)\ge g(a).
\]

Write
\[
L=g(1)=\log(2+\sqrt3).
\]
Superadditivity gives \(g(n)\ge nL\) for every positive integer \(n\).

Choose a power of two \(N>n\). By (5), \(g(N)=NL\), so
\[
NL=g(N)\ge g(n)+g(N-n)\ge g(n)+(N-n)L.
\]
Thus \(g(n)\le nL\), proving
\[
g(n)=nL.
\]
Repeated doubling now yields
\[
g\!\left(\frac{m}{2^k}\right)=\frac{m}{2^k}L
\]
for all nonnegative integers \(m,k\).

\paragraph{3. Squeeze arbitrary real inputs.}

Every nonnegative real number can be approached from below and above by numbers of the form \(m/2^k\). Monotonicity therefore forces
\[
g(x)=Lx\qquad(x\ge0).
\]

Hence \(H(f(x))=(2+\sqrt3)^x\). Using the reciprocal identity for \(H\),
\[
2f(x)=H(f(x))+H(f(x))^{-1},
\]
which gives the claimed formula.

\paragraph{4. Verify.}

Put \(A=2+\sqrt3\), so \(A^{-1}=2-\sqrt3\). For the proposed function,
\[
\sqrt{f(x)^2-1}=\frac{A^x-A^{-x}}2
\qquad(x\ge0).
\]
Expansion shows that (1) holds with equality, as does (2). The domain, range, and normalisation are also satisfied.
\end{proof}
\end{dlsummary}

\begin{takeaways}
\begin{itemize}
\item \textbf{Use conjugate radicals:} The expressions \(t+\sqrt{t^2-1}\) and \(t-\sqrt{t^2-1}\) are reciprocals.
\item \textbf{Linearise before estimating:} A transformation can turn a radical addition rule into superadditivity.
\item \textbf{Exact doubling is powerful:} Together with nonnegativity and superadditivity, it determines values first on integers, then on dyadic rationals.
\item \textbf{Derive regularity:} Monotonicity justifies the final squeeze; continuity need not be assumed.
\end{itemize}
\end{takeaways}

\addcontentsline{toc}{subsubsection}{Solution 48. The One-Point Cubic Lock}

\begin{dlsummary}[title={\bfseries Solution 48. The One-Point Cubic Lock},breakable]

\begin{proof}
The unique solution is
\[
\boxed{f(x)=x^3\qquad(x>0).}
\]

\paragraph{1. Extract a global cubic lower bound.}

Set
\[
g(x)=f(x)-x^3.
\]
Because
\[
(x+y)^3=x^3+y^3+3xy(x+y),
\]
the first inequality becomes
\[
g(x+y)\ge g(x)+g(y). \tag{1}
\]

For any positive integer \(n\), repeated application gives
\[
g(t)\ge n\,g(t/n).
\]
Returning to \(f\), and using its positivity,
\[
f(t)\ge n f(t/n)+t^3\left(1-\frac1{n^2}\right)
>t^3\left(1-\frac1{n^2}\right).
\]
Letting \(n\to\infty\), we obtain
\[
\boxed{f(t)\ge t^3,\quad\text{so }g(t)\ge0.} \tag{2}
\]

\paragraph{2. Lock the value at \(2\).}

Applying (2) twice,
\[
512\ge f(f(2))\ge f(2)^3\ge8^3=512.
\]
Therefore every inequality is an equality, and
\[
\boxed{f(2)=8,\qquad g(2)=0.} \tag{3}
\]

\paragraph{3. Determine the function on an interval.}

For \(0<t<2\), superadditivity gives
\[
0=g(2)\ge g(t)+g(2-t).
\]
Both summands are nonnegative, so \(g(t)=0\). Hence
\[
\boxed{f(t)=t^3\qquad(0<t\le2).} \tag{4}
\]

\paragraph{4. Propagate to every positive input.}

Fix \(x>0\). Choose a positive integer \(k\) such that \(x/2^k<2\). Repeated submultiplicativity, followed by (3) and (4), gives
\[
f(x)
\le f(2)^k f(x/2^k)
=8^k\left(\frac{x}{2^k}\right)^3
=x^3.
\]
Combined with (2), this proves \(f(x)=x^3\).

\textbf{Verification.} The cubic identity makes the first inequality an equality; multiplicativity also holds, and
\[
f(f(2))=f(8)=512.
\]
\end{proof}

\paragraph{Takeaways.}
\begin{itemize}
\item **Subtract the expected model:** Removing \(x^3\) exposes an ordinary superadditive defect.
\item **Subdivide before taking limits:** Positivity of \(f\) controls the potentially negative defect at small inputs.
\item **One exact value can determine an interval:** A nonnegative superadditive function that vanishes at \(2\) must vanish throughout \((0,2)\).
\item **Use different inequalities locally and globally:** Addition determines a bounded interval; multiplication propagates the result everywhere.
\end{itemize}

\end{dlsummary}

\addcontentsline{toc}{subsubsection}{Solution 50. The Golden Staircase}

\begin{dlsummary}[title={\bfseries Solution 50. The Golden Staircase},breakable]

\begin{proof}
The unique solution is
\[
\boxed{
f(n)=\left\lfloor\varphi n\right\rfloor,
\qquad
\varphi=\frac{1+\sqrt5}{2}.
}
\]

\paragraph{1. Prove that a limiting slope exists.}

Repeated application of (1) gives
\[
m f(n)\le f(mn)\le m f(n)+m-1.
\]
Interchanging \(m,n\) and combining the bounds yields
\[
\frac{f(n)}n<\frac{f(m)}m+\frac1m,
\qquad
\frac{f(m)}m<\frac{f(n)}n+\frac1n.
\]
Thus the sequence \(f(n)/n\) is Cauchy. Write
\[
\alpha=\lim_{n\to\infty}\frac{f(n)}n.
\]
Also, superadditivity gives \(f(n)\ge n f(1)\ge n\), so \(\alpha\ge1\).

For fixed \(n\), divide
\[
k f(n)\le f(kn)\le k f(n)+k-1
\]
by \(k\), and let \(k\to\infty\). We obtain
\[
\boxed{\alpha n-1\le f(n)\le\alpha n.} \tag{3}
\]

\paragraph{2. Use composition to determine the slope.}

Divide (2) by \(n\):
\[
\left|
\frac{f(f(n))}{n}-\frac{f(n)}n-1
\right|\le\frac1n.
\]
Since \(f(n)\ge n\), its values tend to infinity, and hence
\[
\frac{f(f(n))}{n}
=
\frac{f(f(n))}{f(n)}\frac{f(n)}n
\longrightarrow\alpha^2.
\]
Therefore
\[
\alpha^2=\alpha+1.
\]
As \(\alpha\ge1\), necessarily \(\alpha=\varphi\).

\paragraph{3. Recover the exact integer values.}

By (3),
\[
\varphi n-1\le f(n)\le\varphi n.
\]
Because \(\varphi n\) is irrational, this interval contains exactly one integer:
\[
f(n)=\lfloor\varphi n\rfloor.
\]

\paragraph{4. Verify the solution.}

The elementary floor inequality
\[
\lfloor a\rfloor+\lfloor b\rfloor
\le\lfloor a+b\rfloor
\le\lfloor a\rfloor+\lfloor b\rfloor+1
\]
proves (1).

For (2), write
\[
k=\lfloor\varphi n\rfloor,
\qquad
t=\varphi n-k\in(0,1).
\]
Using \(\varphi^2=\varphi+1\),
\[
\varphi k
=n+k-(\varphi-1)t.
\]
Since \(0<(\varphi-1)t<1\),
\[
f(f(n))=\lfloor\varphi k\rfloor=n+k-1=n+f(n)-1.
\]
Thus (2) holds, with absolute difference exactly \(1\).
\end{proof}
\end{dlsummary}

\begin{takeaways}
\begin{itemize}
\item \textbf{Compare a common input:} Two estimates for \(f(mn)\) control the ratios \(f(n)/n\).
\item \textbf{Bounded errors can reveal exact slopes:} The additive error disappears after normalisation.
\item \textbf{Composition determines an algebraic number:} The limiting slope must satisfy \(\alpha^2=\alpha+1\).
\item \textbf{Integrality finishes the classification:} A narrow interval identifies each function value exactly.
\end{itemize}
\end{takeaways}

\addcontentsline{toc}{subsubsection}{Solution 52. The Almost-Cosine Sequence}

\begin{dlsummary}[title={\bfseries Solution 52. The Almost-Cosine Sequence},breakable]

\begin{proof}
The unique solution is
\[
\boxed{
f(n)=\frac{(2+\sqrt3)^n+(2-\sqrt3)^n}{2}
\qquad(n\in\mathbb Z).
}
\]

\paragraph{1. Force growth.}

Taking \(m=1,n=0\),
\[
|4-4f(0)|\le1.
\]
Since \(f(0)\) is an integer, \(f(0)=1\).

Now take \(m=n\):
\[
f(2n)\ge2f(n)^2-2.
\]
Starting from \(f(1)=2\), this gives
\[
f(2)\ge6,\qquad f(4)\ge70,
\]
and repeated doubling proves that \(f\) is unbounded.

\paragraph{2. Introduce the error.}

Write
\[
E(a,b)=f(a+b)+f(a-b)-2f(a)f(b),
\qquad |E(a,b)|\le1.
\]

Consider
\[
S=f(x+y+z)+f(x+y-z)+f(x-y+z)+f(x-y-z).
\]

Pairing terms with centres \(x+y\) and \(x-y\), we obtain
\[
S=
4f(x)f(y)f(z)
+2f(z)E(x,y)
+E(x+y,z)+E(x-y,z).
\]

Pairing instead with centre \(x\) and offsets \(y+z\) and \(y-z\), we obtain
\[
S=
4f(x)f(y)f(z)
+2f(x)E(y,z)
+E(x,y+z)+E(x,y-z).
\]

Subtracting these expressions gives
\[
2f(x)|E(y,z)|\le2f(z)+4.
\]
For fixed \(y,z\), the right side is constant, while \(f(x)\) is unbounded. Therefore
\[
E(y,z)=0.
\]
Thus the original inequality secretly forces
\[
\boxed{f(m+n)+f(m-n)=2f(m)f(n).} \tag{1}
\]

\paragraph{3. Recover the sequence.}

Taking \(m=0\) in (1), using \(f(0)=1\), gives
\[
f(-n)=f(n).
\]

Taking \(n=1\), using \(f(1)=2\), gives
\[
f(m+1)=4f(m)-f(m-1).
\]
Consequently, \(f(0)=1\) and \(f(1)=2\) determine every value:
\[
1,\ 2,\ 7,\ 26,\ 97,\ 362,\ldots
\]

The roots of \(r^2-4r+1=0\) are \(2+\sqrt3\) and \(2-\sqrt3\). Solving the recurrence gives the stated formula.

\paragraph{4. Verify.}

Put \(A=2+\sqrt3\), so \(A^{-1}=2-\sqrt3\). Direct expansion shows that
\[
f(n)=\frac{A^n+A^{-n}}2
\]
satisfies (1), hence the original inequality with zero error.

These values are positive, and the recurrence with integer initial values proves they are integers for nonnegative \(n\). Evenness handles negative \(n\).
\end{proof}
\end{dlsummary}

\begin{takeaways}
\begin{itemize}
\item \textbf{Find growth before controlling errors:} Unboundedness supplies the amplification needed to eliminate a fixed discrepancy.
\item \textbf{Regroup a common expression:} The four-term sum links errors at different input pairs.
\item \textbf{Separate rigidity from classification:} First establish the exact equation; only then solve its recurrence.
\item \textbf{Irrational formulas can give integers:} A recurrence often explains integrality more clearly than a closed form.
\end{itemize}
\end{takeaways}

\addcontentsline{toc}{subsubsection}{Solution 55. The Reversible Difference}

\begin{dlsummary}[title={\bfseries Solution 55. The Reversible Difference},breakable]

\begin{proof}
The only solutions are
\[
\boxed{
f(n)=n^2
\quad\text{and}\quad
f(n)=\frac{n(n+1)}2
\qquad(n\in\mathbb Z).
}
\]

\paragraph{1. Reverse the step to force equality.}

Replace \(m\) by \(m+3n\) and \(n\) by \(-n\). The given inequality becomes
\[
f(m)+3f(m+2n)
\ge3f(m+n)+f(m+3n),
\]
which is precisely the reverse inequality.

Therefore
\[
\boxed{
f(m+3n)-3f(m+2n)+3f(m+n)-f(m)=0.
} \tag{1}
\]

\paragraph{2. Recover a quadratic formula.}

Take \(n=1\) in (1), and define
\[
d_m=f(m+2)-2f(m+1)+f(m).
\]
Equation (1) says
\[
d_{m+1}=d_m.
\]
Thus all second differences equal a fixed integer \(k\).

Using \(f(0)=0\) and \(f(1)=1\), the resulting recurrence gives
\[
\boxed{
f(n)=n+k\frac{n(n-1)}2
\qquad(n\in\mathbb Z).
} \tag{2}
\]
For completeness, the right-hand side has the same second differences and initial values as \(f\); the recurrence determines values both forwards and backwards.

\paragraph{3. Use nonnegativity to restrict \(k\).}

At \(n=-1\), formula (2) gives
\[
f(-1)=k-1.
\]
Since \(f(-1)\ge0\),
\[
\boxed{k\ge1.}
\]

\paragraph{4. Use the composition bound to finish.}

We have
\[
f(2)=k+2,
\]
and hence
\[
f(f(2))
=k+2+\frac{k(k+1)(k+2)}2.
\]
If \(k\ge3\), then
\[
f(f(2))\ge5+\frac{3\cdot4\cdot5}{2}=35>16,
\]
a contradiction. Consequently,
\[
k\in\{1,2\}.
\]

Substituting these values into (2) gives
\[
k=1:\quad f(n)=\frac{n(n+1)}2,
\qquad
k=2:\quad f(n)=n^2.
\]

\paragraph{5. Verify both solutions.}

Both formulas give nonnegative integers for every integer \(n\), with the required values at \(0\) and \(1\). Every quadratic satisfies (1), so the original inequality holds with equality.

Finally,
\[
f(f(2))=
\begin{cases}
6,&f(n)=n(n+1)/2,\\
16,&f(n)=n^2.
\end{cases}
\]
Both satisfy the composition bound.
\end{proof}

\paragraph{Takeaways.}
\emph{No takeaways were provided in the source draft.}

\end{dlsummary}

\addcontentsline{toc}{subsubsection}{Solution 57. The Hidden Interval}

\begin{dlsummary}[title={\bfseries Solution 57. The Hidden Interval},breakable]

\begin{proof}
The solutions are precisely the functions that \textbf{move each real number to the nearest point of a nonempty closed interval}, allowing unbounded intervals and single points.

Explicitly:

| Fixed-point interval | Function |
|---|---|
| \([a,b]\), where \(a\le b\) | \(\displaystyle f(x)=\begin{cases}a,&x<a,\\x,&a\le x\le b,\\b,&x>b;\end{cases}\) |
| \([a,\infty)\) | \(f(x)=\max\{x,a\}\) |
| \((-\infty,b]\) | \(f(x)=\min\{x,b\}\) |
| \(\mathbb R\) | \(f(x)=x\) |

The case \(a=b\) includes every constant function.

\paragraph{1. Extract order and continuity.}

For \(x>y\), put \(d=f(x)-f(y)\). Condition (1) says
\[
d^2\le(x-y)d.
\]
Equivalently,
\[
d\bigl(d-(x-y)\bigr)\le0,
\]
so
\[
\boxed{0\le f(x)-f(y)\le x-y.} \tag{3}
\]
Thus \(f\) is nondecreasing and
\[
|f(x)-f(y)|\le|x-y|.
\]
In particular, \(f\) is continuous.

\paragraph{2. Collapse the threefold iteration.}

Fix \(x\), and write
\[
u=f(x),\qquad v=f(u).
\]
By (2),
\[
f(v)=u.
\]

If \(u<v\), monotonicity gives
\[
v=f(u)\le f(v)=u,
\]
a contradiction. If \(u>v\), the same argument gives the opposite contradiction. Hence \(u=v\).

Therefore
\[
\boxed{f(f(x))=f(x).} \tag{4}
\]

\paragraph{3. Identify the fixed-point set.}

Let
\[
C=\{t\in\mathbb R:f(t)=t\}.
\]
Equation (4) shows that every function value belongs to \(C\), so \(C\) is nonempty and equals the image of \(f\). Continuity makes \(C\) closed.

Suppose \(a,b\in C\) and \(a<t<b\). Applying (3) between \(a,t\),
\[
f(t)-a\le t-a,
\]
so \(f(t)\le t\). Applying it between \(t,b\),
\[
b-f(t)\le b-t,
\]
so \(f(t)\ge t\).

Thus \(f(t)=t\). Consequently, \(C\) is an interval.

\paragraph{4. Determine the function outside the interval.}

If \(C\) has a finite lower endpoint \(a\), then \(a\in C\). For \(x<a\), monotonicity gives \(f(x)\le f(a)=a\), while \(f(x)\in C\) gives \(f(x)\ge a\). Hence
\[
f(x)=a.
\]

Similarly, above a finite upper endpoint \(b\), we have \(f(x)=b\). On \(C\), the function is the identity. This gives exactly the listed formulas.

\paragraph{5. Verify.}

Each listed function is nondecreasing and satisfies
\[
0\le f(x)-f(y)\le x-y\qquad(x>y),
\]
which implies (1). Its values lie in an interval that it fixes pointwise, so \(f(f(x))=f(x)\), proving (2).
\end{proof}

\paragraph{Takeaways.}
\begin{itemize}
\item **Factor a quadratic inequality:** It can simultaneously reveal monotonicity and a distance bound.
\item **Order rules out nontrivial two-cycles:** This turns threefold iteration into idempotence.
\item **Study the fixed-point set:** Its geometry can determine the entire function.
\item **Expect piecewise answers:** A functional inequality need not force a polynomial or power function.
\end{itemize}

\end{dlsummary}

\addcontentsline{toc}{subsubsection}{Solution 67. The Forbidden Coefficient}

\begin{dlsummary}[title={\bfseries Solution 67. The Forbidden Coefficient},breakable]

\begin{proof}
The only solutions are
\[
\boxed{f(x)=x^3\quad\text{and}\quad f(x)=x^3+x.}
\]

\paragraph{1. Remove the cubic contribution.}

Set
\[
g(x)=f(x)-x^3.
\]
The identity
\[
(x+y)^3-x^3-y^3=3xy(x+y)
\]
turns (1) into
\[
0\le g(x+y)-g(x)-g(y)\le(x-y)^2. \tag{4}
\]

Taking \(y=x\) gives
\[
g(2x)=2g(x).
\]
Now apply (4) at \(x/2^n,y/2^n\) and multiply by \(2^n\):
\[
0\le g(x+y)-g(x)-g(y)
\le\frac{(x-y)^2}{2^n}.
\]
Letting \(n\to\infty\), we obtain
\[
\boxed{g(x+y)=g(x)+g(y).} \tag{5}
\]

\paragraph{2. Use positivity to eliminate irregular additive functions.}

For \(0<t<1\), positivity of \(f\) gives
\[
g(t)=f(t)-t^3>-1.
\]
By additivity,
\[
g(t)=g(1)-g(1-t)<g(1)+1.
\]
Thus \(g\) is bounded in absolute value on \((0,1)\); write \(|g(t)|\le M\) there.

For any positive integer \(N\), if \(0<h<1/N\), then
\[
N|g(h)|=|g(Nh)|\le M.
\]
Consequently, \(g(h)\to0\) as \(h\to0^+\). Additivity then gives continuity at every positive input.

Writing \(c=g(1)\), additivity first gives \(g(q)=cq\) for positive rationals, and continuity extends this to every positive real number. Hence
\[
\boxed{f(x)=x^3+cx.} \tag{6}
\]
Positivity for arbitrarily small \(x>0\) forces \(c\ge0\).

\paragraph{3. Discover the forbidden coefficient interval.}

Substitution into (2) gives
\[
f(x)f(y)-f(xy)
=cxy\bigl(c+x^2+y^2-1\bigr)\ge0.
\]
This factorisation was also checked symbolically.[1]

If \(c>0\), division by \(cxy\) yields
\[
c+x^2+y^2\ge1.
\]
Taking arbitrarily small positive \(x,y\), we conclude \(c\ge1\). Therefore
\[
\boxed{c=0\quad\text{or}\quad c\ge1.} \tag{7}
\]

\paragraph{4. Apply the nested-value lock.}

From (6),
\[
f(1)=1+c,
\]
so (3) becomes
\[
(1+c)^3+c(1+c)\le10.
\]
Equivalently,
\[
(c-1)(c^2+5c+9)\le0.
\]
Since \(c\ge0\), the second factor is positive, and thus \(c\le1\).

Together with (7), this leaves exactly
\[
c=0\quad\text{or}\quad c=1.
\]

\paragraph{5. Verify.}

For either candidate, the middle expression in (1) is zero.

For \(f(x)=x^3\), condition (2) is equality and \(f(f(1))=1\).

For \(f(x)=x^3+x\),
\[
f(x)f(y)-f(xy)=xy(x^2+y^2)>0,
\]
and
\[
f(f(1))=f(2)=10.
\]
Both functions satisfy every condition.
\end{proof}

\paragraph{Takeaways.}
\begin{itemize}
\item **Subtract the polynomial suggested by the defect:** This exposes a nearly additive remainder.
\item **Exploit competing scaling rates:** A quadratic error bound cannot survive exact linear scaling.
\item **Use positivity locally:** It supplies the boundedness needed to rule out discontinuous additive behaviour.
\item **Look for coefficient gaps:** Submultiplicativity excludes every \(c\in(0,1)\), rather than merely bounding \(c\).
\end{itemize}

\end{dlsummary}


\subsection*{Part IV. Mean and Norm Squeezes}
\addcontentsline{toc}{subsection}{Part IV. Mean and Norm Squeezes}

\addcontentsline{toc}{subsubsection}{Solution 21. 2026 IMO Problem 5}

\begin{dlsummary}[title={\bfseries Solution 21. 2026 IMO Problem 5},breakable]

\begin{proof}
We use the arithmetic-orbit approach from the AoPS Wiki.
Write $P(x,y)$ for the given chain. Substituting $(f(x),x)$ squeezes
its middle term between two copies of $f(x)$, so
\[
f(f(x))=2f(x)-x.
\]
Set $d(x)=f(x)-x$. Then $d(f(x))=d(x)$ and
\[
f^n(x)=x+nd(x)\qquad(n\geq0),
\]
so positivity of all iterates gives $d(x)\geq0$.

Suppose $A=d(x)>0$ and $B=d(y)>0$. For integers $m\geq0$, $n\geq1$,
set $X=x+mA$ and $Y=y+nB$. Since $d(f(t))=d(t)$, we have
$f(X)=X+A$ and $f(Y-B)=Y$. Applying the right inequality to
$(X,Y-B)$ gives
\[
\frac{X+Y+A-B}{2}\geq\sqrt{XY}.
\]
Squaring and rearranging yields
\[
(X-Y)^2\geq2(B-A)(X+Y)-(B-A)^2. \tag{1}
\]
If $A<B$, then for each sufficiently large $n$ choose $m\geq0$ so
that $|X-Y|\leq A/2$; this is possible because the values $x+mA$ have
spacing $A$ and $Y$ tends to infinity. The left side of (1) is then at
most $A^2/4$, while its right side tends to $+\infty$, a contradiction.
Interchanging $x,y$ rules out $B<A$. Hence all positive values of $d$
equal one constant $c>0$, if any occur.

Thus $d$ takes only the values $0,c$. If $d(x)=c$ and $d(y)=0$, the
left inequality in $P(x,y)$ gives
\[
(x-y)^2\geq2c(x+y)+c^2>c^2.
\]
So the two level sets are relatively open; connectedness of
$(0,\infty)$ forces one to be empty. Therefore $d$ is constant and
$f(x)=x+c$ for some $c\geq0$.

For these functions, the chain reduces to the RMS--AM--GM inequalities
for $x$ and $y+c$, so every $c\geq0$ works.
\end{proof}

\end{dlsummary}
\addcontentsline{toc}{subsubsection}{Solution 22. The Hidden Quadratic Function}

\begin{dlsummary}[title={\bfseries Solution 22. The Hidden Quadratic Function},breakable]

\begin{proof}
Put $g(x)=\sqrt{f(x)}$. The given inequalities
are equivalent to
\begin{align}
g(x+y)&\ge g(x)+g(y), \tag{1}\\
g(xy)&\le g(x)g(y). \tag{2}
\end{align}
Since $g>0$, condition (1) implies that $g$
is strictly increasing. From (2), $g(1)\ge1$.

\medskip
\noindent\textbf{Step 1. Show that $g(2)=2$.}

By (1),
\[
g(2)\ge2g(1)\ge2,\qquad
g(4)\ge4g(1)\ge4.
\]
The condition $f(f(2))\le16$ gives
\[
g(g(2)^2)\le4.
\]
Since $g(2)^2\ge4$ and $g$ is increasing,
\[
4\ge g(g(2)^2)\ge g(4)
\ge4g(1)\ge4.
\]
Thus $g(1)=1$ and, by strict monotonicity,
\[
\boxed{g(2)=2.}
\]

\medskip
\noindent\textbf{Step 2. Show that $g(x)\le x$.}

For every positive integer $n$, (1) and (2)
give
\[
2^n=2^n g(1)
\le g(2^n)\le g(2)^n=2^n.
\]
Hence
\[
\boxed{g(2^n)=2^n.}
\]

Fix $x>0$ and put
\[
m_n=\left\lfloor\frac{2^n}{x}\right\rfloor.
\]
For sufficiently large $n$, we have $m_n\ge1$.
By monotonicity and (1),
\[
2^n=g(2^n)\ge g(m_nx)\ge m_ng(x).
\]
Consequently,
\[
g(x)\le
\frac{2^n}{\lfloor2^n/x\rfloor}
\xrightarrow[n\to\infty]{}x.
\]
Therefore
\[
\boxed{g(x)\le x\qquad(x>0).}
\]

\medskip
\noindent\textbf{Step 3. Establish equality.}

By (2) and the preceding bound,
\[
1=g(1)
\le g(x)g(1/x)
\le \frac{g(x)}{x}.
\]
Thus $g(x)\ge x$, and hence
\[
g(x)=x\qquad(x>0).
\]
It follows that
\[
\boxed{f(x)=x^2\qquad(x>0).}
\]

Finally, this function satisfies $f(f(2))=16$
and
\[
\sqrt{f(x+y)}
=\sqrt{f(x)}+\sqrt{f(y)}
=\sqrt{f(xy)}
\left(
\frac1{\sqrt{f(x)}}+
\frac1{\sqrt{f(y)}}
\right)
=x+y.
\]
Therefore it is indeed the unique solution.
\end{proof}

\end{dlsummary}
\addcontentsline{toc}{subsubsection}{Solution 23. The Euclidean--Multiplicative Squeeze}

\begin{dlsummary}[title={\bfseries Solution 23. The Euclidean--Multiplicative Squeeze},breakable]

\begin{proof}
Define
\[
g(t)=f(\sqrt t)^2\qquad(t>0).
\]
The first inequality in the statement gives
\[
g(x+y)\geq g(x)+g(y). \tag{1}
\]
For the second inequality, its right-hand side is
\[
f(xy)^2\frac{f(x)^2+f(y)^2}{f(x)^2f(y)^2}.
\]
Since $f(x)^2+f(y)^2>0$, cancellation yields
\[
f(xy)^2\leq f(x)^2f(y)^2.
\]
Applying this with $\sqrt x$ and $\sqrt y$ in place of $x$ and $y$ gives
\[
g(xy)\leq g(x)g(y). \tag{2}
\]
The composition condition is equivalently
\[
g(g(2))=f(f(\sqrt2))^2\leq2. \tag{3}
\]

\medskip
\noindent
\textbf{Step 1. Determine $g(1)$ and $g(2)$.}

Taking $x=y=1$ in (2) gives $g(1)\leq g(1)^2$. Since $g(1)>0$,
we have $g(1)\geq1$. By repeated use of (1),
\[
g(n)\geq n g(1)\geq n\qquad(n\in\mathbb N).
\]
Fix $t>1$ and a positive integer $n$. Submultiplicativity and
superadditivity imply
\[
g(t)^n\geq g(t^n)\geq\lfloor t^n\rfloor.
\]
Taking $n$th roots and letting $n\to\infty$ gives $g(t)\geq t$.
In particular, $g(2)\geq2$. Applying this bound again at $g(2)>1$,
\[
2\geq g(g(2))\geq g(2)\geq2,
\]
where (3) supplies the first inequality. Thus $g(2)=2$. Now (2) at
$x=y=1$ gives $g(1)\leq g(1)^2$, while (1) and $g(2)=2$ give
$2\geq2g(1)$. Therefore $g(1)=1$.

\medskip
\noindent
\textbf{Step 2. Obtain a global upper bound.}

For every positive integer $n$, (1) and (2) give
\[
2^n\leq g(2^n)\leq g(2)^n=2^n,
\]
so $g(2^n)=2^n$. Fix $x>0$ and, for sufficiently large $n$, let
\[
m_n=\left\lfloor\frac{2^n}{x}\right\rfloor\geq1.
\]
Write $2^n=m_nx+r_n$ with $r_n\geq0$. If $r_n>0$, repeated
superadditivity and positivity of $g(r_n)$ show that
\[
2^n=g(2^n)\geq m_ng(x)+g(r_n)\geq m_ng(x).
\]
If $r_n=0$, repeated superadditivity directly gives the same final
bound. Hence
\[
g(x)\leq\frac{2^n}{\lfloor 2^n/x\rfloor}\longrightarrow x,
\]
and so $g(x)\leq x$ for every $x>0$.

\medskip
\noindent
\textbf{Step 3. Force equality.}

By (2), $g(1)\leq g(x)g(1/x)$. Using $g(1)=1$ and the upper bound
just proved,
\[
1\leq g(x)g(1/x)\leq\frac{g(x)}x.
\]
Thus $g(x)\geq x$, and consequently $g(x)=x$ for every $x>0$.
Since $g(x)=f(\sqrt x)^2$ and $f$ is positive, it follows that
\[
f(x)=x\qquad(x>0).
\]

\medskip
\noindent
\textbf{Verification.}

For $f(x)=x$, both inequalities hold with equality, and
$f(f(\sqrt2))=\sqrt2$. Therefore this function is indeed the unique
solution.
\end{proof}

\end{dlsummary}
\addcontentsline{toc}{subsubsection}{Solution 24. The Quadratic Mean Squeeze}

\begin{dlsummary}[title={\bfseries Solution 24. The Quadratic Mean Squeeze},breakable]

\begin{proof}
Set $g(t)=f(\sqrt t)$. The given chain becomes
\[
\sqrt{\frac{x^2+g(y)^2}{2}}
\geq\frac{g(x)+y}{2}\geq\sqrt{xg(y)}. \tag{1}
\]
Taking $x=g(y)$ forces equality throughout, so
\[
g(g(y))=2g(y)-y.
\]
For $d(x)=g(x)-x$, this gives $d(g(x))=d(x)$ and
$g^{\circ n}(x)=x+nd(x)$. Positivity of every iterate implies $d(x)\geq0$.

Squaring (1) and factoring the difference of the outer squares yields
\[
|d(x)-d(y)|\bigl(g(x)+y+x+g(y)\bigr)
\leq (x-g(y))^2. \tag{2}
\]
Indeed, each outer square differs from $(x+g(y))^2$ by
$(x-g(y))^2$.

Any two positive values of $d$ must coincide. Otherwise, let
$d(p)=r>0$, $d(q)=s>0$, $r\ne s$. The iterates
$X_n=p+nr$ and $Y_m=q+ms$ retain these respective $d$-values.
For each large $n$, choose $m\geq0$ with $|X_n-Y_m|\leq s/2$.
In (2) at $(X_n,Y_m)$, the left side tends to infinity, while
$(X_n-g(Y_m))^2=(X_n-Y_m-s)^2\leq9s^2/4$, a contradiction.

Thus $d$ takes at most two values, $0$ and some $c>0$. If both occur,
(2) shows that any points with different values are more than $c$
apart. Bisecting an interval whose endpoints have different values
produces such pairs arbitrarily close together, impossible. Hence
$d$ is constant, say $c\geq0$, and $f(x)=g(x^2)=x^2+c$.

The composition condition gives
\[
a^4\geq f(f(a))=(a^2+c)^2+c,
\]
forcing $c=0$. Thus $f(x)=x^2$. Substitution reduces the chain to
RMS--AM--GM for $x^2,y^2$, and $f(f(a))=a^4$.
\end{proof}

\end{dlsummary}
\addcontentsline{toc}{subsubsection}{Solution 47. The Three-Variable Quadratic Trap}

\begin{dlsummary}[title={\bfseries Solution 47. The Three-Variable Quadratic Trap},breakable]

\begin{proof}
The unique solution is
\[
\boxed{f(x)=x^2\qquad(x\in\mathbb R).}
\]

\paragraph{1. Force evenness and exact doubling.}

Condition (2) at \(x=0\), together with nonnegativity, gives \(f(0)=0\).

In (1), substitute \((x,y,z)=(t,t,-t)\):
\[
f(2t)\le3f(t)+f(-t).
\]
Combining this with (2) gives \(f(t)\le f(-t)\). Replacing \(t\) by \(-t\) proves
\[
f(-t)=f(t).
\]
The same inequalities now yield
\[
\boxed{f(2t)=4f(t).} \tag{3}
\]

\paragraph{2. Make the inequality reverse itself.}

Substitute \(z=-y\) in (1). Evenness gives
\[
f(x+y)+f(x-y)\le2f(x)+2f(y). \tag{4}
\]

Now apply (4) to \(x+y\) and \(x-y\):
\[
f(2x)+f(2y)\le2f(x+y)+2f(x-y).
\]
Using (3), this is the reverse of (4). Consequently,
\[
\boxed{f(x+y)+f(x-y)=2f(x)+2f(y).} \tag{5}
\]

\paragraph{3. Determine a dense set of values.}

Taking \(y=1\) in (5),
\[
f(n+1)+f(n-1)=2f(n)+2.
\]
Starting from \(f(0)=0\) and \(f(1)=1\), induction gives
\[
f(n)=n^2\qquad(n\in\mathbb Z).
\]
Exact doubling then gives
\[
f\!\left(\frac{m}{2^k}\right)=\left(\frac{m}{2^k}\right)^2
\]
for every integer \(m\) and nonnegative integer \(k\).

\paragraph{4. Derive continuity instead of assuming it.}

Equations (3) and (5) imply midpoint convexity:
\[
f\!\left(\frac{u+v}{2}\right)
\le\frac{f(u)+f(v)}2.
\]
Iterating this gives the corresponding inequality for dyadic weights.

Fix \(x\), and let \(N\) be a power of two. If \(|h|\le1/N\), then evenness and the assumed unit-interval bound give \(f(Nh)\le1\). By (5) and nonnegativity,
\[
f(x+Nh)\le2f(x)+2f(Nh)\le2f(x)+2.
\]
Using the dyadic weight \(1/N\),
\[
f(x+h)\le
\left(1-\frac1N\right)f(x)+\frac1N f(x+Nh)
\le f(x)+\frac{f(x)+2}{N}.
\]
The same upper bound holds for \(f(x-h)\). Since midpoint convexity gives
\[
2f(x)\le f(x+h)+f(x-h),
\]
we obtain
\[
|f(x+h)-f(x)|\le\frac{f(x)+2}{N}.
\]
Thus \(f\) is continuous everywhere.

Approximating any real number by dyadic rationals now proves \(f(x)=x^2\).

\paragraph{5. Verify.}

For \(f(x)=x^2\), expansion shows that (1) holds with equality. All other conditions are immediate.
\end{proof}

\paragraph{Takeaways.}
\begin{itemize}
\item **Exploit sign changes:** The substitution \((t,t,-t)\) creates evenness from a one-sided inequality.
\item **Change coordinates:** Replacing \((x,y)\) by \((x+y,x-y)\) reverses the parallelogram inequality.
\item **Local bounds can replace continuity:** A bound on one interval, combined with the functional structure, controls the function everywhere.
\item **Do not stop at rational values:** Extending a classification to all real inputs requires a justified regularity argument.
\end{itemize}

\end{dlsummary}

\addcontentsline{toc}{subsubsection}{Solution 51. The Logarithm Hidden in a Chain}

\begin{dlsummary}[title={\bfseries Solution 51. The Logarithm Hidden in a Chain},breakable]

\begin{proof}
The unique solution is
\[
\boxed{f(x)=\log x\qquad(x>0),}
\]
where \(\log\) denotes the natural logarithm.

\paragraph{1. Find the normalisation and a pointwise bound.}

Setting \(x=y=1\), we get
\[
f(1)\le2f(1)\le0.
\]
Thus \(f(1)=0\).

Now set \(y=1\) in the upper inequality:
\[
\boxed{f(x)\le x-1.} \tag{1}
\]

\paragraph{2. Improve the upper bound by taking roots.}

Repeated application of the first inequality gives
\[
f(t^n)\le n f(t)
\]
for every positive integer \(n\).

Taking \(t=x^{1/n}\) and using (1),
\[
f(x)\le n f(x^{1/n})
\le n\bigl(x^{1/n}-1\bigr).
\]
The standard elementary limit
\[
\lim_{n\to\infty}n\bigl(x^{1/n}-1\bigr)=\log x
\]
therefore yields
\[
\boxed{f(x)\le\log x.} \tag{2}
\]

\paragraph{3. Use reciprocals to obtain the reverse bound.}

The first inequality at \((x,1/x)\) gives
\[
0=f(1)\le f(x)+f(1/x).
\]
Applying (2) at \(1/x\),
\[
f(x)\ge-f(1/x)\ge-\log(1/x)=\log x.
\]
Together with (2), this proves \(f(x)=\log x\).

\paragraph{4. Verify.}

The logarithm satisfies
\[
\log(xy)=\log x+\log y,
\]
and the standard inequality \(\log t\le t-1\) gives
\[
\log x+\log y\le x+y-2.
\]
Thus it satisfies the entire chain.
\end{proof}

\paragraph{Takeaways.}
\begin{itemize}
\item **Multiplicative subdivision sharpens bounds:** Replacing \(x\) by \(n\) identical factors improves \(f(x)\le x-1\) to the exact logarithmic bound.
\item **Reciprocals reverse one-sided information:** An upper bound at \(1/x\) becomes a lower bound at \(x\).
\item **Normalisation can be forced:** Neither \(f(1)=0\) nor a value at any other input needs to be supplied.
\item **No continuity of \(f\) is needed:** The limit concerns an explicit numerical expression, not function values approaching an input.
\end{itemize}

\end{dlsummary}

\addcontentsline{toc}{subsubsection}{Solution 56. The Microscopic Error Trap}

\begin{dlsummary}[title={\bfseries Solution 56. The Microscopic Error Trap},breakable]

\begin{proof}
The complete family is
\[
\boxed{f(x)=ax\qquad(x\in\mathbb R),\quad a\in\mathbb R.}
\]

\paragraph{1. Extract exact doubling.}

Taking \(x=y=0\) in (1) gives \(f(0)=0\).

Taking \(y=x\), the upper bound vanishes, so
\[
\boxed{f(2x)=2f(x).} \tag{3}
\]
Consequently,
\[
f(x/2^n)=\frac{f(x)}{2^n}
\]
for every nonnegative integer \(n\).

\paragraph{2. Shrink the error until it disappears.}

Apply (1) at \(x/2^n,y/2^n\). Using (3),
\[
0\le
\frac{f(x)+f(y)-f(x+y)}{2^n}
\le\frac{(x-y)^2}{4^n}.
\]
Multiplying by \(2^n\),
\[
0\le f(x)+f(y)-f(x+y)
\le\frac{(x-y)^2}{2^n}.
\]
Letting \(n\to\infty\), we obtain
\[
\boxed{f(x+y)=f(x)+f(y).} \tag{4}
\]

Additivity alone is \textbf{not} enough to conclude that \(f\) is linear.

\paragraph{3. Find regularity in the composite.}

Define
\[
h(x)=f(f(x)).
\]
By (4), \(h\) is additive. Condition (2) says
\[
h(x^2)\ge f(x)^2\ge0.
\]
Every nonnegative real number is a square, so
\[
h(t)\ge0\qquad(t\ge0).
\]
Thus \(h\) is nondecreasing: for \(u<v\),
\[
h(v)-h(u)=h(v-u)\ge0.
\]

An additive, nondecreasing function is linear. Indeed, writing \(c=h(1)\ge0\), additivity gives \(h(q)=cq\) for rational \(q\), and rational approximation plus monotonicity gives
\[
\boxed{h(x)=cx\qquad(x\in\mathbb R).} \tag{5}
\]

\paragraph{4. Transfer the regularity back to \(f\).}

Combining (2) and (5),
\[
f(x)^2\le h(x^2)=cx^2.
\]
Therefore
\[
\boxed{|f(x)|\le\sqrt c\,|x|.} \tag{6}
\]
By additivity,
\[
|f(u)-f(v)|=|f(u-v)|\le\sqrt c\,|u-v|.
\]
Hence \(f\) is continuous.

Put \(a=f(1)\). Additivity gives \(f(q)=aq\) for every rational \(q\); continuity extends this to every real input:
\[
f(x)=ax.
\]

\paragraph{5. Verify every slope.}

For any real \(a\), the function \(f(x)=ax\) has zero additive error, and
\[
f(f(x^2))=a^2x^2=f(x)^2.
\]
Thus every real slope works, including \(a=0\) and negative slopes.
\end{proof}

\paragraph{Takeaways.}
\emph{No takeaways were provided in the source draft.}

\end{dlsummary}

\addcontentsline{toc}{subsubsection}{Solution 63. The Hidden Euclidean Length}

\begin{dlsummary}[title={\bfseries Solution 63. The Hidden Euclidean Length},breakable]

\begin{proof}
The unique solution is
\[
\boxed{F(a,b)=\sqrt{a^2+b^2}\qquad(a,b\in\mathbb Q).}
\]

\paragraph{1. Recover the basic symmetries.}

Using (2),
\[
F(-1,0)^2=F(1,0)=1,
\]
so \(F(-1,0)=1\). Consequently,
\[
F(-a,-b)=F(a,b).
\]

Similarly,
\[
F(0,1)^2=F(-1,0)=1,
\]
so \(F(0,1)=1\). Thus
\[
F(0,t)=F(t,0)\qquad(t\in\mathbb Q). \tag{3}
\]

\paragraph{2. Determine the values on the axes.}

Repeated use of (1) gives
\[
F(n,0)\le n\qquad(n\in\mathbb Z_{>0}).
\]
Meanwhile, multiplicativity and the normalisation give
\[
F(2^k,0)=2^k.
\]

Fix a positive integer \(n\), and choose \(2^k>n\). By (1),
\[
2^k
\le F(n,0)+F(2^k-n,0)
\le F(n,0)+2^k-n.
\]
Therefore \(F(n,0)\ge n\), proving equality.

For positive integers \(m,n\), multiplicativity now gives
\[
m=F(m,0)=F(m/n,0)F(n,0)=nF(m/n,0).
\]
Using symmetry and (3), we obtain
\[
\boxed{F(t,0)=F(0,t)=|t|\qquad(t\in\mathbb Q).} \tag{4}
\]

In particular, (1) yields the preliminary bound
\[
\boxed{F(a,b)\le |a|+|b|.} \tag{5}
\]

\paragraph{3. Improve the bound by taking powers.}

Define multiplication of pairs by
\[
(a,b)\ast(c,d)=(ac-bd,ad+bc).
\]
For a fixed pair \((a,b)\), write its \(n\)-th power as \((A_n,B_n)\). All coordinates remain rational.

Condition (2) gives
\[
F(A_n,B_n)=F(a,b)^n.
\]
The elementary identity
\[
(ac-bd)^2+(ad+bc)^2=(a^2+b^2)(c^2+d^2)
\]
also gives
\[
A_n^2+B_n^2=(a^2+b^2)^n.
\]

Apply (5) and \((|u|+|v|)^2\le2(u^2+v^2)\):
\[
\begin{aligned}
F(a,b)^n
&\le |A_n|+|B_n|\\
&\le\sqrt2\,(a^2+b^2)^{n/2}.
\end{aligned}
\]
For a nonzero pair, take \(n\)-th roots and let \(n\to\infty\):
\[
\boxed{F(a,b)\le\sqrt{a^2+b^2}.} \tag{6}
\]

\paragraph{4. Use the conjugate pair to force equality.}

Since
\[
(a,b)\ast(a,-b)=(a^2+b^2,0),
\]
equations (2) and (4) give
\[
F(a,b)F(a,-b)=a^2+b^2.
\]
But (6) bounds both factors above by \(\sqrt{a^2+b^2}\). Their product reaches the product of these upper bounds, so both must attain them:
\[
F(a,b)=\sqrt{a^2+b^2}.
\]
The zero pair is already determined.

\paragraph{5. Verify.}

The displayed algebraic identity proves (2). For (1), use
\[
ac+bd\le\sqrt{a^2+b^2}\sqrt{c^2+d^2},
\]
which follows from \((ad-bc)^2\ge0\). Expanding the squares then proves the required triangle inequality.

All normalisations hold.
\end{proof}

\paragraph{Takeaways.}
\begin{itemize}
\item **An exact value can propagate backwards:** Knowing the values at powers of \(2\) determines every integer on an axis.
\item **High powers sharpen estimates:** The crude factor \(\sqrt2\) disappears after taking \(n\)-th roots.
\item **Conjugates produce a known product:** Pairing \((a,b)\) with \((a,-b)\) forces equality in both upper bounds.
\item **Geometry can emerge from algebra:** Neither distance nor length appears in the hypotheses.
\end{itemize}

\end{dlsummary}


\subsection*{Part V. Real-Line and Polynomial Inequalities}
\addcontentsline{toc}{subsection}{Part V. Real-Line and Polynomial Inequalities}

\addcontentsline{toc}{subsubsection}{Solution 25. Vietnam Mathematical Olympiad 1991}

\begin{dlsummary}[title={\bfseries Solution 25. Vietnam Mathematical Olympiad 1991},breakable]

\begin{proof}
\textbf{Step 1. Determine $f(0)$.}

Setting $x=y=z=0$ gives
\[
f(0)-f(0)^2\geq\frac14,
\]
or equivalently,
\[
\left(f(0)-\frac12\right)^2\leq0.
\]
Hence
\[
\boxed{f(0)=\frac12.} \tag{1}
\]

\medskip
\noindent
\textbf{Step 2. Establish a key inequality.}

Setting $z=0$ and using (1), we obtain
\[
\frac{f(xy)}2+\frac14-\frac{f(x)}2
\geq\frac14.
\]
Therefore,
\[
\boxed{f(xy)\geq f(x)}
\qquad(x,y\in\mathbb R). \tag{2}
\]

\medskip
\noindent
\textbf{Step 3. Prove that $f$ is constant.}

For any $t\neq0$, taking $x=1,y=t$ in (2)
gives
\[
f(t)\geq f(1).
\]
Taking $x=t,y=1/t$ gives
\[
f(1)\geq f(t).
\]
Thus
\[
f(t)=f(1)\qquad(t\neq0). \tag{3}
\]

Finally, setting $x=y=z=1$ gives
\[
f(1)-f(1)^2\geq\frac14,
\]
so
\[
\left(f(1)-\frac12\right)^2\leq0.
\]
Consequently, $f(1)=\frac12$.

Combining this with (1) and (3), we get
\[
\boxed{f(x)=\frac12\qquad(x\in\mathbb R).}
\]

Conversely, for $f(x)=\frac12$, we have
\[
\frac{f(xy)+f(xz)}2-f(x)f(yz)
=\frac12-\frac14=\frac14.
\]
Thus the inequality holds with equality
for all real $x,y,z$.
\end{proof}

\end{dlsummary}
\addcontentsline{toc}{subsubsection}{Solution 26. IMO 2023 Shortlist -- Algebra A2}

\begin{dlsummary}[title={\bfseries Solution 26. IMO 2023 Shortlist -- Algebra A2},breakable]

\begin{proof}
\textbf{Step 1. Introduce a symmetric substitution.}

Put
\[
s=x+y,\qquad t=x-y,
\]
so that
\[
x=\frac{s+t}{2},\qquad
y=\frac{s-t}{2}.
\]
The given inequality becomes
\[
f(s)f(t)\geq
f\left(\frac{s+t}{2}\right)^2
-f\left(\frac{s-t}{2}\right)^2.
\tag{1}
\]

\medskip
\noindent
\textbf{Step 2. Force a sign constraint.}

Replacing $t$ by $-t$ in (1), we obtain
\[
f(s)f(-t)\geq
f\left(\frac{s-t}{2}\right)^2
-f\left(\frac{s+t}{2}\right)^2.
\tag{2}
\]

Adding (1) and (2), the squared terms cancel:
\[
\boxed{
f(s)\bigl(f(t)+f(-t)\bigr)\geq0
}
\tag{3}
\]
for every $s,t\in\mathbb R$.

\medskip
\noindent
\textbf{Step 3. Apply the strictness assumption.}

By assumption, the original inequality is strict
for some pair $(x_0,y_0)$.

Set
\[
s_0=x_0+y_0,\qquad t_0=x_0-y_0.
\]
Then (1) is strict at $(s_0,t_0)$.
Adding the corresponding inequality (2) yields
\[
f(s_0)\bigl(f(t_0)+f(-t_0)\bigr)>0.
\]
Therefore,
\[
\boxed{f(t_0)+f(-t_0)\neq0.}
\]

Now fix $t=t_0$ in (3). For every
$s\in\mathbb R$,
\[
f(s)\bigl(f(t_0)+f(-t_0)\bigr)\geq0.
\]

Since the second factor is a fixed nonzero
real number, either
\[
f(s)\geq0\qquad\text{for all }s\in\mathbb R,
\]
or
\[
f(s)\leq0\qquad\text{for all }s\in\mathbb R.
\]

This proves the required conclusion.
\end{proof}

\end{dlsummary}
\addcontentsline{toc}{subsubsection}{Solution 27. IMO 2011 -- Problem 3}

\begin{dlsummary}[title={\bfseries Solution 27. IMO 2011 -- Problem 3},breakable]

\begin{proof}
Replacing $y$ by $t-x$ in the given inequality gives
\[
f(t)\leq(t-x)f(x)+f(f(x))\qquad(x,t\in\mathbb R). \tag{1}
\]

\medskip
\noindent
\textbf{Step 1. Derive a two-variable inequality.}

Set $t=f(z)$ in (1), then add the resulting inequality to the one
obtained by interchanging $x$ and $z$. The terms involving $f(f(x))$
and $f(f(z))$ cancel, yielding
\[
xf(x)+zf(z)\leq2f(x)f(z)\qquad(x,z\in\mathbb R). \tag{2}
\]
Taking $z=2f(x)$ in (2) cancels the terms involving
$f(2f(x))$ and gives
\[
xf(x)\leq0\qquad(x\in\mathbb R). \tag{3}
\]
In particular, $f(x)\leq0$ for $x>0$, while $f(x)\geq0$ for $x<0$.

\medskip
\noindent
\textbf{Step 2. Show that $f$ is nonpositive everywhere.}

Suppose $f(w)>0$ for some $w\in\mathbb R$. Taking $x=w$ in (1),
\[
f(t)\leq(t-w)f(w)+f(f(w)).
\]
Choose
\[
t<\min\!\left(0,\,w-\frac{f(f(w))}{f(w)}\right).
\]
Then $t<0$ and the right-hand side is negative, so $f(t)<0$.
This contradicts $f(t)\geq0$ for $t<0$, established in Step 1.
Therefore $f(x)\leq0$ for every real $x$.

\medskip
\noindent
\textbf{Step 3. Prove the required vanishing.}

For $x<0$, Step 1 gives $f(x)\geq0$, while Step 2 gives
$f(x)\leq0$. Hence $f(x)=0$ for all $x<0$. Using (1) with any
$x<0$ and $t<0$ now gives
\[
0=f(t)\leq(t-x)f(x)+f(f(x))=f(0),
\]
so $f(0)\geq0$. Step 2 also gives $f(0)\leq0$, and therefore
$f(0)=0$. Thus $f(x)=0$ for every $x\leq0$.
\end{proof}

\end{dlsummary}
\addcontentsline{toc}{subsubsection}{Solution 28. The Degree-Normalized Cubic Squeeze}

\begin{dlsummary}[title={\bfseries Solution 28. The Degree-Normalized Cubic Squeeze},breakable]

\begin{proof}
Define $g(x)=f(x)/x^3$, so $f(x)=x^3g(x)$. Substitution into the
given inequalities gives
\[
xg(x)+yg(y)\leq(x+y)g(x+y),
\qquad
g(x)+g(y)\geq2g(x+y).
\]
Equivalently,
\[
\frac{xg(x)+yg(y)}{x+y}
\leq g(x+y)\leq\frac{g(x)+g(y)}2
\qquad(x,y>0). \tag{1}
\]

\medskip
\noindent
\textbf{Step 1. Establish dyadic invariance and monotonicity.}

Taking $y=x$ in (1) yields
\[
g(x)\leq g(2x)\leq g(x),
\]
so
\[
g(2x)=g(x)\qquad(x>0). \tag{2}
\]
The lower and upper bounds in (1) also imply
\[
2\bigl(xg(x)+yg(y)\bigr)\leq(x+y)\bigl(g(x)+g(y)\bigr),
\]
which simplifies to
\[
(x-y)\bigl(g(x)-g(y)\bigr)\leq0.
\]
Thus $g$ is non-increasing.

\medskip
\noindent
\textbf{Step 2. Show that $g$ is constant.}

Fix $a>0$. For any $x\in[a,2a]$, monotonicity and (2) give
\[
g(a)=g(2a)\leq g(x)\leq g(a),
\]
so $g(x)=g(a)$ throughout $[a,2a]$. These intervals cover
$(0,\infty)$ as $a$ varies, and therefore $g$ is constant: there is
some $c\in\mathbb R$ such that $g(x)=c$ for all $x>0$.
Consequently,
\[
f(x)=cx^3\qquad(x>0).
\]

\medskip
\noindent
\textbf{Verification.}

For any real constant $c$, substituting $f(x)=cx^3$ into the first
inequality reduces it to
$c(x+y)\leq c(x+y)$, and the second to
$2c\geq2c$. Hence every such function satisfies both conditions.
\end{proof}

\end{dlsummary}
\addcontentsline{toc}{subsubsection}{Solution 29. The Conditioned Cubic Squeeze}

\begin{dlsummary}[title={\bfseries Solution 29. The Conditioned Cubic Squeeze},breakable]

\begin{proof}
\textbf{Step 1. Establish oddness and monotonicity.}

On the plane $a+b+c=0$, both conditions apply, so
\[
f(a^3)+f(b^3)+f(c^3)=3f(abc). \tag{1}
\]
Taking $(a,b,c)=(t,-t,0)$ and using $f(0)=0$ gives
$f(-t^3)=-f(t^3)$. Since every real number is a cube, $f$ is odd.

If $x\geq y$, take
$a=\sqrt[3]{x}$, $b=-\sqrt[3]{y}$, and $c=0$. Their sum is
nonnegative, and the first condition gives
\[
f(x)+f(-y)+f(0)\geq3f(0).
\]
Oddness and $f(0)=0$ imply $f(x)\geq f(y)$. Thus $f$ is
non-decreasing.

\medskip
\noindent
\textbf{Step 2. Derive a boundary identity.}

In (1), take $(a,b,c)=(u,v,-u-v)$. Oddness yields
\[
f((u+v)^3)=f(u^3)+f(v^3)+3f(uv(u+v)). \tag{2}
\]
Expand $f((a+b+c)^3)$ by applying (2) first to $(a+b,c)$ and then
to $(a,b)$. Expanding instead first with $(a+c,b)$ gives, after
cancelling the common terms,
\[
f(ab(a+b))+f((a+b)c(a+b+c))
=f(ac(a+c))+f((a+c)b(a+b+c)). \tag{3}
\]

\medskip
\noindent
\textbf{Step 3. Obtain Jensen's equation and affinity.}

Let $x,y>0$ with $x\ne y$, and set
\[
k=\sqrt{\frac yx}\ne1,
\qquad 2k(k^2-1)t^3=x,
\]
where the nonzero real $t$ exists uniquely. Choose
\[
a=(k+1)t,\qquad b=(k-1)t,\qquad c=k(k-1)t.
\]
Direct calculation gives
\[
ab(a+b)=x,\qquad (a+b)c(a+b+c)=y,
\]
and
\[
ac(a+c)=(a+c)b(a+b+c)=\frac{x+y}{2}.
\]
Substitution in (3) proves
\[
f(x)+f(y)=2f\!\left(\frac{x+y}{2}\right). \tag{4}
\]
This also holds when $x=y$, so (4) holds for all positive $x,y$.

Because $f$ is non-decreasing, the midpoint identity implies that $f$
is affine on $(0,\infty)$. Indeed, repeated midpoint interpolation
first gives linear interpolation at dyadic fractions of any interval;
monotonicity then squeezes the value at every intermediate point between
such dyadic approximations. Thus there are constants $m,b$ such that
$f(x)=mx+b$ for all $x>0$.

Substitute this form into (2), for positive $u,v$. The left side is
$m(u+v)^3+b$, while the right side is
$m(u^3+v^3+3uv(u+v))+5b=m(u+v)^3+5b$. Hence $b=0$.
Since $f(1)=1$, we get $m=1$, so $f(x)=x$ for $x>0$. Oddness
and $f(0)=0$ extend this to every real $x$.

\medskip
\noindent
\textbf{Verification.}

For the identity function,
\[
a^3+b^3+c^3-3abc
=\frac12(a+b+c)\bigl((a-b)^2+(b-c)^2+(c-a)^2\bigr).
\]
Its sign agrees with the sign of $a+b+c$, so both required
inequalities hold.
\end{proof}

\end{dlsummary}
\addcontentsline{toc}{subsubsection}{Solution 30. The Quadratic--Difference Squeeze}

\begin{dlsummary}[title={\bfseries Solution 30. The Quadratic--Difference Squeeze},breakable]

\begin{proof}
Taking $y=x$ in the first inequality gives $f(x)\geq x^2$ for every
real $x$. Define
\[
g(x)=f(x)-x^2.
\]
Then $g(x)\geq0$. Substituting $f(t)=t^2+g(t)$ into the second
inequality and cancelling the quadratic terms yields
\[
g(x+y)+g(x-y)\leq2g(x)\qquad(x,y\in\mathbb R). \tag{1}
\]

\medskip
\noindent
\textbf{Step 1. Show that $g$ is constant.}

Fix arbitrary $a,b\in\mathbb R$, and set
\[
x_k=a+k(b-a)\qquad(k\in\mathbb Z),
\qquad
\Delta_k=g(x_{k+1})-g(x_k).
\]
Applying (1) with $x=x_k$ and $y=b-a$ gives
\[
g(x_{k+1})+g(x_{k-1})\leq2g(x_k),
\]
so $\Delta_k\leq\Delta_{k-1}$. In particular,
$\Delta_j\leq\Delta_0=g(b)-g(a)$ for every integer $j\geq0$.
For any positive integer $n$, summing these bounds gives
\[
-g(a)\leq g(x_n)-g(a)
=\sum_{j=0}^{n-1}\Delta_j
\leq n\bigl(g(b)-g(a)\bigr),
\]
where the first inequality uses $g(x_n)\geq0$.
If $g(b)<g(a)$, choosing $n$ sufficiently large makes the rightmost
quantity less than $-g(a)$, a contradiction. Thus $g(b)\geq g(a)$.
Since $a,b$ were arbitrary, interchanging them shows $g(a)=g(b)$;
therefore $g$ is constant. Write $g(x)=d$. Its nonnegativity gives
$d\geq0$, and hence
\[
f(x)=x^2+d\qquad(x\in\mathbb R).
\]

\medskip
\noindent
\textbf{Verification.}

For $d\geq0$,
\[
f(x)+f(y)-2xy=(x-y)^2+2d\geq0,
\]
and
\[
f(x+y)+f(x-y)
=(x+y)^2+(x-y)^2+2d
=2f(x)+2y^2.
\]
Thus every function in the stated family satisfies both inequalities.
\end{proof}

\end{dlsummary}
\addcontentsline{toc}{subsubsection}{Solution 31. The Iterated Parallelogram--Quadratic Squeeze}

\begin{dlsummary}[title={\bfseries Solution 31. The Iterated Parallelogram--Quadratic Squeeze},breakable]

\begin{proof}
\textbf{Step 1. Derive a quadratic lower bound.}

Comparing the outer expressions in the given chain gives
\[
f(x)-x^2+f(f(x))-f(y)+2y^2-2yf(x)\geq0.
\]
Setting $y=f(x)$ cancels the terms involving $f(f(x))$ and the quadratic
terms in $y$, leaving
\[
f(x)\geq x^2\qquad(x\in\mathbb R). \tag{1}
\]
Define $g(x)=f(x)-x^2$. Then $g(x)\geq0$ for every real $x$.

\medskip
\noindent
\textbf{Step 2. Obtain midpoint concavity of the residual.}

The right-hand inequality and the identity
\[
\frac{(x+y)^2+(x-y)^2}{2}=x^2+y^2
\]
show that
\[
g(x+y)+g(x-y)\leq2g(x)\qquad(x,y\in\mathbb R). \tag{2}
\]

\medskip
\noindent
\textbf{Step 3. Prove that $g$ is constant.}

Fix $a,b\in\mathbb R$, and write
\[
x_k=a+k(b-a)\quad(k\geq0),
\qquad
\Delta_k=g(x_{k+1})-g(x_k).
\]
For each $k\geq1$, apply (2) with $x=x_k$ and $y=b-a$. This gives
\[
g(x_{k+1})+g(x_{k-1})\leq2g(x_k),
\]
so $\Delta_k\leq\Delta_{k-1}$. Consequently,
$\Delta_k\leq\Delta_0=g(b)-g(a)$ for all $k\geq0$. For every positive
integer $n$, nonnegativity of $g$ now yields
\[
-g(a)\leq g(x_n)-g(a)
=\sum_{k=0}^{n-1}\Delta_k
\leq n\bigl(g(b)-g(a)\bigr).
\]
Thus $g(b)-g(a)\geq-g(a)/n$ for every $n$, and letting $n\to\infty$
shows $g(b)\geq g(a)$. Since $a,b$ were arbitrary, interchanging them
proves $g(a)=g(b)$. Therefore $g(x)=d$ for a constant $d\geq0$, and
\[
f(x)=x^2+d. \tag{3}
\]

\medskip
\noindent
\textbf{Step 4. Use the composition condition and verify.}

By (3),
\[
f(f(2))=(4+d)^2+d=16+9d+d^2\leq16.
\]
Since $d\geq0$, this forces $d=0$. Hence $f(x)=x^2$ for every real $x$.

For this function, $f(f(2))=16$, the middle and right expressions in
the given chain are both $x^2+y^2$, and the left expression equals
\[
x^2+y^2-(y-x^2)^2\leq x^2+y^2.
\]
Thus $f(x)=x^2$ satisfies all the conditions, completing the proof.
\end{proof}

\end{dlsummary}
\addcontentsline{toc}{subsubsection}{Solution 32. The Skew-Symmetric Parallelogram Squeeze}

\begin{dlsummary}[title={\bfseries Solution 32. The Skew-Symmetric Parallelogram Squeeze},breakable]

\begin{proof}
\textbf{Step 1. Obtain a quadratic lower bound.}

Comparing the outer expressions in the given chain gives
\[
2xy+x^2f(y)-y^2f(x)\leq f(x)+y^2.
\]
Setting $y=x$ cancels the skew-symmetric terms, so
\[
2x^2\leq f(x)+x^2,
\qquad\text{hence}\qquad
f(x)\geq x^2\quad(x\in\mathbb R). \tag{1}
\]
Define $g(x)=f(x)-x^2$. Then $g(x)\geq0$ for all real $x$.

\medskip
\noindent
\textbf{Step 2. Establish midpoint concavity of the residual.}

The right-hand inequality and the parallelogram identity give
\[
\frac{(x+y)^2+(x-y)^2}{2}
+\frac{g(x+y)+g(x-y)}2
\leq x^2+y^2+g(x).
\]
After cancellation,
\[
g(x+y)+g(x-y)\leq2g(x)
\qquad(x,y\in\mathbb R). \tag{2}
\]

\medskip
\noindent
\textbf{Step 3. Prove that $g$ is constant.}

Fix $a,b\in\mathbb R$, and put
\[
x_k=a+k(b-a)\quad(k\geq0),
\qquad
\Delta_k=g(x_{k+1})-g(x_k).
\]
For each $k\geq1$, applying (2) with $x=x_k$ and $y=b-a$ yields
\[
g(x_{k+1})+g(x_{k-1})\leq2g(x_k),
\]
so $\Delta_k\leq\Delta_{k-1}$. Thus $\Delta_k\leq\Delta_0=g(b)-g(a)$
for all $k\geq0$. For every positive integer $n$, nonnegativity of $g$
implies
\[
-g(a)\leq g(x_n)-g(a)
=\sum_{k=0}^{n-1}\Delta_k
\leq n\bigl(g(b)-g(a)\bigr).
\]
Letting $n\to\infty$ gives $g(b)\geq g(a)$. Interchanging $a$ and $b$
shows $g(a)=g(b)$, so $g(x)=d$ for some constant $d\geq0$. Therefore
\[
f(x)=x^2+d. \tag{3}
\]

\medskip
\noindent
\textbf{Step 4. Eliminate the constant shift.}

Substituting (3) into the outer inequality gives
\[
(x-y)^2+d(1-x^2+y^2)\geq0
\qquad(x,y\in\mathbb R).
\]
Taking $y=x-1$ reduces this to
\[
1+d(2-2x)\geq0\qquad(x\in\mathbb R).
\]
If $d>0$, this fails for sufficiently large $x$. Hence $d=0$ and
$f(x)=x^2$ for all real $x$.

For $f(x)=x^2$, the middle and right-hand expressions are both
$x^2+y^2$, while the left-hand expression is $2xy\leq x^2+y^2$.
Thus this function satisfies the stated inequalities and is the unique
solution.
\end{proof}

\end{dlsummary}
\addcontentsline{toc}{subsubsection}{Solution 33. The Dyadic Parallelogram--Quadratic Squeeze}

\begin{dlsummary}[title={\bfseries Solution 33. The Dyadic Parallelogram--Quadratic Squeeze},breakable]

\begin{proof}
\textbf{Step 1. Obtain an inequality for the residual.}

Put $g(x)=f(x)-x^2$. Comparing the outer expressions in the given chain
and substituting $f(x)=x^2+g(x)$ gives
\[
2g(x)+3g(y)+2(x-2y)^2\geq0
\qquad(x,y\in\mathbb R). \tag{1}
\]
Taking $x=y=0$ shows $g(0)\geq0$, while setting $x=2t$ and $y=t$ gives
\[
2g(2t)+3g(t)\geq0\qquad(t\in\mathbb R). \tag{2}
\]

\medskip
\noindent
\textbf{Step 2. Prove that $g$ is nonnegative.}

The right-hand inequality in the problem and the parallelogram identity
imply
\[
g(x+y)+g(x-y)\leq2g(x)
\qquad(x,y\in\mathbb R). \tag{3}
\]
Taking $x=y=t$ gives $g(2t)\leq2g(t)-g(0)$. Combining this with (2),
\[
0\leq2g(2t)+3g(t)\leq7g(t)-2g(0).
\]
Since $g(0)\geq0$, it follows that $g(t)\geq\frac27g(0)\geq0$ for
all $t\in\mathbb R$.

\medskip
\noindent
\textbf{Step 3. Show that $g$ is constant.}

Fix $a,b\in\mathbb R$, and define
\[
x_k=a+k(b-a)\quad(k\geq0),
\qquad
\Delta_k=g(x_{k+1})-g(x_k).
\]
For each $k\geq1$, apply (3) with $x=x_k$ and $y=b-a$. Then
\[
g(x_{k+1})+g(x_{k-1})\leq2g(x_k),
\]
so $\Delta_k\leq\Delta_{k-1}$. Therefore
$\Delta_k\leq\Delta_0=g(b)-g(a)$ for all $k\geq0$. For any positive
integer $n$, nonnegativity of $g$ yields
\[
-g(a)\leq g(x_n)-g(a)
=\sum_{k=0}^{n-1}\Delta_k
\leq n\bigl(g(b)-g(a)\bigr).
\]
Letting $n\to\infty$ gives $g(b)\geq g(a)$. Interchanging $a$ and $b$
shows that $g$ is constant, say $g(x)=d$ with $d\geq0$. Hence
\[
f(x)=x^2+d. \tag{4}
\]

\medskip
\noindent
\textbf{Step 4. Determine $d$ and verify the solution.}

Using (4) in the composition condition gives
\[
1\geq f(f(1))=(1+d)^2+d=1+3d+d^2.
\]
As $d\geq0$, this forces $d=0$, so $f(x)=x^2$ for all real $x$.

For $f(x)=x^2$, the middle and right-hand expressions in the problem
both equal $5x^2+5y^2$, while the left-hand expression is
\[
3x^2+8xy-3y^2=5x^2+5y^2-2(x-2y)^2
\leq5x^2+5y^2.
\]
Also $f(f(1))=1$. Thus $f(x)=x^2$ satisfies all conditions and is the
unique solution.
\end{proof}

\end{dlsummary}
\addcontentsline{toc}{subsubsection}{Solution 39. The Twin Squeeze at Zero}

\begin{dlsummary}[title={\bfseries Solution 39. The Twin Squeeze at Zero},breakable]

\begin{proof}
The only solution is \(f(x)=x+1\).

\paragraph{1. Find \(f(0)\).} Write \(a=f(0)\). The second inequality at \(x=y=0\) gives \(a\le 2a-1\), so \(a\ge1\). The first inequality at \(x=y=0\) gives \(a^2\le a\), so \(a\le1\). Hence \(f(0)=1\).

\paragraph{2. Upper bound.} Put \(y=0\) in the first inequality:
\[
f(x)\cdot1\le f(0)+x=x+1.
\]
So \(f(x)\le x+1\) for every real \(x\). \((1)\)

\paragraph{3. Lower bound by symmetry.} Put \(y=-x\) in the second inequality:
\[
1=f(0)\le f(x)+f(-x)-1,
\]
so \(f(x)\ge 2-f(-x)\). Apply (1) at \(-x\) to get \(f(-x)\le 1-x\). Therefore
\[
f(x)\ge 2-(1-x)=x+1.
\]

Together with (1), \(f(x)=x+1\). It satisfies both conditions with equality. \(\blacksquare\)
\end{proof}
\end{dlsummary}

\begin{takeaways}
\begin{itemize}
\item \textbf{Opposing inequalities fix constants:} The two conditions pushed \(f(0)\) in opposite directions, which forced \(f(0)=1\).
\item \textbf{Use \(y=-x\) to reverse a bound:} An upper bound on \(f(-x)\) became a lower bound on \(f(x)\). The same move solves many inequality problems.
\item \textbf{Shift to a cleaner function:} Setting \(g=f-1\) makes the second condition ordinary subadditivity.
\item \textbf{Predict the answer by testing a family:} Trying \(f(x)=cx+d\) shows that \(x+1\) is the only candidate.
\end{itemize}
\end{takeaways}


\newpage
\section*{Conclusion}
\addcontentsline{toc}{section}{Conclusion}

The central lesson of these problems is to keep experimenting with the
inequalities: choose useful inputs, normalize the function when appropriate,
look for quantities to compare, and follow a promising bound until it becomes
decisive. Symmetry, reciprocal substitutions, iteration, density, and
discrete arithmetic can all turn a weak-looking inequality into a rigid
identity.

Not every problem ends with the same kind of conclusion. Some force a unique
function, some leave a parameter family, and some require a separate analysis
of exceptional points or equality cases. In every case, respect the domain,
prove continuity or other regularity rather than assuming it, verify the final
candidates in the original statement, and use failed approaches as information
about the hidden structure of the problem. We hope this booklet encourages you
to approach unfamiliar problems with patience, curiosity, and confidence.

\end{document}
